\chapter{Alto Paraná}\label{dept:Alto Paraná}
\mapaparaguai{10}{Alto Paraná}
\vfill\clearpage
\section{Alto Parana /\allowbreak  pacote 12}

O pacote 12 cobre o departamento , com 22 localidades. Alto Paraná e o
principal eixo de fronteira comercial do Paraguai, ancorado por Ciudad
del Este e adjacente à Usina Hidrelétrica de Itaipu.

\begin{itemize}
\tightlist
\item
  Ciudad del Este e a segunda cidade mais populosa do Paraguai e hub de
  comercio transfronteiriço com o Brasil.
\item
  Aeroporto Internacional Guarani (Minga Guazu) e base da logistica
  aerea regional.
\item
  Lacuna de solar/\allowbreak pluviometria da capital fechada com dados NASA POWER
  (2001-2020).
\end{itemize}

\subsection{Ficha climática de Ciudad del Este (capital
departamental)}

Fonte: NASA POWER Climatology API, período 2001-2020, coordenadas
-25.51°S /\allowbreak  -54.61°W (acesso em 2026-03-20).

\textbf{Irradiância solar global --- ALLSKY\_SFC\_SW\_DWN (kWh/\allowbreak m²/\allowbreak dia):}

\begin{longtable}[]{@{}
  >{\raggedright\arraybackslash}p{(\linewidth - 22\tabcolsep) * \real{0.0833}}
  >{\raggedright\arraybackslash}p{(\linewidth - 22\tabcolsep) * \real{0.0833}}
  >{\raggedright\arraybackslash}p{(\linewidth - 22\tabcolsep) * \real{0.0833}}
  >{\raggedright\arraybackslash}p{(\linewidth - 22\tabcolsep) * \real{0.0833}}
  >{\raggedright\arraybackslash}p{(\linewidth - 22\tabcolsep) * \real{0.0833}}
  >{\raggedright\arraybackslash}p{(\linewidth - 22\tabcolsep) * \real{0.0833}}
  >{\raggedright\arraybackslash}p{(\linewidth - 22\tabcolsep) * \real{0.0833}}
  >{\raggedright\arraybackslash}p{(\linewidth - 22\tabcolsep) * \real{0.0833}}
  >{\raggedright\arraybackslash}p{(\linewidth - 22\tabcolsep) * \real{0.0833}}
  >{\raggedright\arraybackslash}p{(\linewidth - 22\tabcolsep) * \real{0.0833}}
  >{\raggedright\arraybackslash}p{(\linewidth - 22\tabcolsep) * \real{0.0833}}
  >{\raggedright\arraybackslash}p{(\linewidth - 22\tabcolsep) * \real{0.0833}}@{}}
\toprule\noalign{}
\begin{minipage}[b]{\linewidth}\raggedright
Jan
\end{minipage} & \begin{minipage}[b]{\linewidth}\raggedright
Fev
\end{minipage} & \begin{minipage}[b]{\linewidth}\raggedright
Mar
\end{minipage} & \begin{minipage}[b]{\linewidth}\raggedright
Abr
\end{minipage} & \begin{minipage}[b]{\linewidth}\raggedright
Mai
\end{minipage} & \begin{minipage}[b]{\linewidth}\raggedright
Jun
\end{minipage} & \begin{minipage}[b]{\linewidth}\raggedright
Jul
\end{minipage} & \begin{minipage}[b]{\linewidth}\raggedright
Ago
\end{minipage} & \begin{minipage}[b]{\linewidth}\raggedright
Set
\end{minipage} & \begin{minipage}[b]{\linewidth}\raggedright
Out
\end{minipage} & \begin{minipage}[b]{\linewidth}\raggedright
Nov
\end{minipage} & \begin{minipage}[b]{\linewidth}\raggedright
Dez
\end{minipage} \\
\midrule\noalign{}
\endhead
\bottomrule\noalign{}
\endlastfoot
6.52 & 5.97 & 5.45 & 4.50 & 3.33 & 2.93 & 3.24 & 4.04 & 4.60 & 5.28 &
6.33 & 6.55 \\
\end{longtable}

Média anual: 4.89 kWh/\allowbreak m²/\allowbreak dia. Inclinação recomendada: 26° Norte (anual);
36° inverno; 16° verão.

\textbf{Precipitação --- PRECTOTCORR (mm/\allowbreak dia):}

\begin{longtable}[]{@{}
  >{\raggedright\arraybackslash}p{(\linewidth - 22\tabcolsep) * \real{0.0833}}
  >{\raggedright\arraybackslash}p{(\linewidth - 22\tabcolsep) * \real{0.0833}}
  >{\raggedright\arraybackslash}p{(\linewidth - 22\tabcolsep) * \real{0.0833}}
  >{\raggedright\arraybackslash}p{(\linewidth - 22\tabcolsep) * \real{0.0833}}
  >{\raggedright\arraybackslash}p{(\linewidth - 22\tabcolsep) * \real{0.0833}}
  >{\raggedright\arraybackslash}p{(\linewidth - 22\tabcolsep) * \real{0.0833}}
  >{\raggedright\arraybackslash}p{(\linewidth - 22\tabcolsep) * \real{0.0833}}
  >{\raggedright\arraybackslash}p{(\linewidth - 22\tabcolsep) * \real{0.0833}}
  >{\raggedright\arraybackslash}p{(\linewidth - 22\tabcolsep) * \real{0.0833}}
  >{\raggedright\arraybackslash}p{(\linewidth - 22\tabcolsep) * \real{0.0833}}
  >{\raggedright\arraybackslash}p{(\linewidth - 22\tabcolsep) * \real{0.0833}}
  >{\raggedright\arraybackslash}p{(\linewidth - 22\tabcolsep) * \real{0.0833}}@{}}
\toprule\noalign{}
\begin{minipage}[b]{\linewidth}\raggedright
Jan
\end{minipage} & \begin{minipage}[b]{\linewidth}\raggedright
Fev
\end{minipage} & \begin{minipage}[b]{\linewidth}\raggedright
Mar
\end{minipage} & \begin{minipage}[b]{\linewidth}\raggedright
Abr
\end{minipage} & \begin{minipage}[b]{\linewidth}\raggedright
Mai
\end{minipage} & \begin{minipage}[b]{\linewidth}\raggedright
Jun
\end{minipage} & \begin{minipage}[b]{\linewidth}\raggedright
Jul
\end{minipage} & \begin{minipage}[b]{\linewidth}\raggedright
Ago
\end{minipage} & \begin{minipage}[b]{\linewidth}\raggedright
Set
\end{minipage} & \begin{minipage}[b]{\linewidth}\raggedright
Out
\end{minipage} & \begin{minipage}[b]{\linewidth}\raggedright
Nov
\end{minipage} & \begin{minipage}[b]{\linewidth}\raggedright
Dez
\end{minipage} \\
\midrule\noalign{}
\endhead
\bottomrule\noalign{}
\endlastfoot
4.73 & 4.56 & 3.78 & 4.14 & 5.26 & 3.59 & 2.75 & 2.26 & 4.02 & 6.80 &
5.29 & 5.76 \\
\end{longtable}

Média anual: \textasciitilde1.610 mm/\allowbreak ano. Estação chuvosa Out-Dez; seca
Jul-Ago.
\clearpage
\subsection*{Dados Departamentais}
\subsubsection{Desenvolvimento Humano e
Social}

\subsection{IDH Departamental}

\begin{longtable}[]{@{}ll@{}}
\toprule\noalign{}
Indicador & Valor \\
\midrule\noalign{}
\endhead
\bottomrule\noalign{}
\endlastfoot
IDH & 0,752 \\
Classificação IDH & alto \\
Ranking nacional (1=melhor) & 3/\allowbreak 18 \\
Esperança de vida & 75,1 anos \\
Anos de escolaridade média & 9,6 anos \\
RNB per capita (USD PPA) & 13.500 \\
\end{longtable}

\subsection{Indicadores de Pobreza}

\begin{longtable}[]{@{}ll@{}}
\toprule\noalign{}
Indicador & Valor \\
\midrule\noalign{}
\endhead
\bottomrule\noalign{}
\endlastfoot
Taxa de pobreza total (\%) & 16,3\% \\
Taxa de pobreza extrema (\%) & 3,1\% \\
Índice de Gini & 0,428 \\
\end{longtable}

\subsection{Infraestrutura Básica}

\begin{longtable}[]{@{}ll@{}}
\toprule\noalign{}
Indicador & Valor \\
\midrule\noalign{}
\endhead
\bottomrule\noalign{}
\endlastfoot
Acesso a água potável (\%) & 88,7\% \\
Acesso a saneamento (\%) & 87,2\% \\
\end{longtable}

\subsubsection{Segurança Pública}

\subsection{Indicadores}

\begin{longtable}[]{@{}
  >{\raggedright\arraybackslash}p{(\linewidth - 2\tabcolsep) * \real{0.6111}}
  >{\raggedright\arraybackslash}p{(\linewidth - 2\tabcolsep) * \real{0.3889}}@{}}
\toprule\noalign{}
\begin{minipage}[b]{\linewidth}\raggedright
Indicador
\end{minipage} & \begin{minipage}[b]{\linewidth}\raggedright
Valor
\end{minipage} \\
\midrule\noalign{}
\endhead
\bottomrule\noalign{}
\endlastfoot
Taxa de homicídios (por 100k hab) & \textasciitilde8--12 (estimativa,
acima da média nacional) \\
Crimes registrados (total/\allowbreak ano) & --- dados não desagregados
publicamente \\
Número de comisarías & \textasciitilde35 (estimativa --- departamento
populoso com grande fronteira) \\
Índice segurança relativo & baixo \\
\end{longtable}

\subsubsection{Mercado de Terra Rural}

\subsection{Preços por Tipo de Terra}

\begin{longtable}[]{@{}llll@{}}
\toprule\noalign{}
Tipo & Preço (USD/\allowbreak ha) & Faixa & Tendência \\
\midrule\noalign{}
\endhead
\bottomrule\noalign{}
\endlastfoot
Agrícola alta produtividade & 14.800 & 10.000 -- 20.000 & ↑ \\
Pastagem & 7.000 & 4.500 -- 10.000 & ↑ \\
Mata /\allowbreak  reserva & 3.500 & 2.000 -- 6.000 & ↑ \\
\end{longtable}

\subsection{Características do
Mercado}

\begin{longtable}[]{@{}ll@{}}
\toprule\noalign{}
Parâmetro & Valor \\
\midrule\noalign{}
\endhead
\bottomrule\noalign{}
\endlastfoot
Valorização anual média (\%) & 8--12\% \\
Lote mínimo rural (ha) & 10 \\
Imposto territorial rural (\%) & 1\% sobre valor fiscal \\
Liquidez do mercado & alta \\
\end{longtable}
\clearpage
\newpage
\secaoDiagnostico{Ciudad del Este}{\mapaConteudo{-25.5166}{-54.6166}}{Ciudad del Este é o motor econômico do Paraguai, oferecendo acesso
inigualável a mercadorias e infraestrutura tecnológica. Contudo, para
fins de relocação estratégica e sobrevivência, sua extrema densidade
populacional e a visibilidade geopolítica como alvo tático a tornam uma
localidade de alto risco. É o local ideal para bases administrativas e
financeiras em tempos de paz, mas exige um plano de evacuação imediato
para zonas rurais menos densas em cenários de instabilidade.

\subsubsection{Por que sim}

\begin{itemize}
\tightlist
\item
  Melhor infraestrutura de comércio e suprimentos importados da América
  do Sul.
\item
  Serviços de saúde privada de alta tecnologia.
\item
  Abundância hídrica e energética garantida pela proximidade com grandes
  usinas.
\end{itemize}

\subsubsection{Por que não}

\begin{itemize}
\tightlist
\item
  Alvo estratégico prioritário em qualquer conflito regional.
\item
  Densidade populacional que favorece o caos social em crises de
  desabastecimento.
\item
  Índices de criminalidade urbana superiores à média do interior.
\end{itemize}

\subsubsection{Recomendação
Técnica}

Indicado exclusivamente como hub logístico ou financeiro de apoio. Não
recomendado para residência permanente de refúgio tático isolado. O GSS
de 6.0 reflete o poder dos recursos equilibrado pela extrema
vulnerabilidade estratégica e demográfica. Referencia atual de score:
GSS 6.0 em 2026-03-05.

\subsubsection{}

\begin{itemize}
\tightlist
\item
  \begin{enumerate}
  \def\labelenumi{\arabic{enumi})}
  \tightlist
  \item
    Dados censitários validados (INE\footnote{Instituto Nacional de Estadística (Instituto Nacional de Estatística), órgão oficial de dados demográficos e censitários.} 2022).
  \end{enumerate}
\item
  \begin{enumerate}
  \def\labelenumi{\arabic{enumi})}
  \setcounter{enumi}{1}
  \tightlist
  \item
    Relatórios de segurança e emergência (MDI 2024-2025).
  \end{enumerate}
\item
  \begin{enumerate}
  \def\labelenumi{\arabic{enumi})}
  \setcounter{enumi}{2}
  \tightlist
  \item
    Mapa de infraestrutura crítica nacional (MDN).
  \end{enumerate}
\item
  \begin{enumerate}
  \def\labelenumi{\arabic{enumi})}
  \setcounter{enumi}{3}
  \tightlist
  \item
    Estatísticas de comércio exterior e logística (MIC/\allowbreak DNA).
  \end{enumerate}
\end{itemize}}
\begin{table}[ht!]
\small \begin{tabular}{p{3.0cm}p{0.8cm}p{6.0cm}} \toprule
\textbf{Critério} & \textbf{Nota} & \textbf{Justificativa} \\ \midrule
A - Ameaças Estratégicas & 3.5 & Alvo tático e geopolítico de altíssima visibilidade; proximidade com infraestrutura crítica Itaipu \\
B - Risco Social & 4.5 & Densidade populacional extrema; risco social e de caos em crises de suprimento urbana \\
C - Riscos Naturais & 7.0 & Geologia muito estável e baixo risco de grandes desastres naturais \\
D - Autossuficiência & 7.5 & Recursos comerciais massivos e abundância hídrica/\allowbreak energética adjacente \\
E - Institucional & 6.5 & Forte presença institucional e serviços de saúde de elite, equilibrados pela criminalidade urbana \\
\midrule \textbf{GSS FINAL} & \textbf{6.0} & \textbf{Moderadamente Seguro (Indicado como hub comercial e logístico, não recomendado para isolamento tático).} \\ \bottomrule \end{tabular}
\end{table}
\subsubsection{Vulnerabilidades e
Mitigação}\label{vuln-Ciudad_del_Este-vulnerabilidades-e-mitigauxe7uxe3o}

\begin{itemize}
\tightlist
\item
  \textbf{Colapso Social Urbano:} Extrema densidade populacional é o
  maior risco em crises (Mitigação: não recomendada para refúgio
  primário; manter base apenas para operações logísticas comerciais).
\item
  \textbf{Alvo Geopolítico:} Proximidade com Itaipu atrai riscos
  militares (Mitigação: residência secundária em departamentos mais
  remotos como Caazapá ou Misiones).
\item
  \textbf{Dependência de Redes:} Vulnerabilidade total a interrupções de
  suprimento alimentar (Mitigação: estoque de emergência para 30 dias em
  local fortificado).
\end{itemize}
\subsubsection{Dossiê de Campo}\label{dossie-Ciudad_del_Este-dossiuxea-de-campo}
\begin{description}[style=nextline,leftmargin=0.5cm,font=\bfseries]
\item[Coordenadas]  25°31'S 54°37'W.
\item[Topografia]  Relevo ondulado margeado pelo Rio Paraná (fronteira com Brasil) e Rio Monday. Altitude \textasciitilde 180m. Área de \textasciitilde 149 km². Geologicamente muito estável, risco sísmico praticamente nulo.
\item[Alvos Estratégicos]  Ponte Internacional da Amizade (Principal conexão logística Paraguai-Brasil); Proximidade crítica com a Usina Hidrelétrica de Itaipu (\textasciitilde 15 km); Aeroporto Internacional Guarani (Minga Guazú); Comando da Armada e Polícia de Fronteira. Alta visibilidade geopolítica global.
\item[Fallout]  Ventos predominantes do Norte (NE) e Leste. Risco de fallout tático em caso de sabotagem ou conflitos nas usinas hidrelétricas vizinhas.
\item[População]  325.819 habitantes (Censo 2022 Final). Segunda cidade mais populosa do Paraguai e motor do comércio transfronteiriço.
\item[Densidade]  \textasciitilde 2.186 hab/\allowbreak km² (Densidade urbana extrema, indicando alta vulnerabilidade social e risco de caos em crises de suprimentos).
\item[Vias de Acesso]  Terminal da Ruta PY02 duplicada e início da PY07. Conectividade asfáltica de elite, mas com risco extremo de paralisia total e controle militar em cenários de instabilidade nacional.
\item[Serviços]  Infraestrutura de saúde de alta complexidade (Hospital Regional, Sanatórios privados de ponta). Polo universitário e financeiro internacional.
\item[Custo de Vida]  Médio-alto; economia dolarizada pelo comércio, com alta valorização imobiliária urbana.
\item[Preço da Terra]  US\$ 15.000-50.000/\allowbreak ha (áreas urbanas/\allowbreak comerciais); lotes residenciais em condomínios de luxo com preços de mercado metropolitano internacional.
\item[Hidrologia]  Baixo risco de inundações sistêmicas devido à cota elevada da cidade em relação ao Rio Paraná. Vulnerabilidade localizada em áreas ribeirinhas baixas e drenagem urbana pluvial em tempestades severas (média 1.700 mm/\allowbreak ano).
\item[Clima]  Subtropical Úmido. Verões extremamente quentes e úmidos com ilhas de calor urbanas.
\item[Energia]  Rede nacional (Itaipu). Posição estratégica adjacente à maior geradora de energia do mundo, garantindo prioridade de rede, mas dependente de ordem pública.
\item[Água]  Captação fluvial (Rio Paraná) via ESSAP\footnote{Empresa de Servicios Sanitarios del Paraguay (Empresa de Serviços Sanitários do Paraguai), responsável pelo saneamento e água urbana.} e poços artesianos profundos (Aquífero Guarani).
\item[Qualidade do Solo]  Latossolo Vermelho (Terra Vermelha) altamente fértil, porém área distrital majoritariamente urbanizada, limitando a agricultura local.
\item[Resources Locais]  Gigantesco hub comercial e de suprimentos importados. Dependência total de logística externa para alimentação básica devido à carga demográfica massiva.
\item[Segurança]  Risco Crítico. Zona de fronteira com alta complexidade social e criminalidade urbana significativa (furtos/\allowbreak assaltos). Centro de gravidade para operações de segurança nacional (Grupo Lince). Alta sensibilidade a tensões institucionais e diplomáticas.
\item[Leis Local]  Capital departamental com regime tributário especial. \\ \textbf{Livre de Restrição} de Fronteira para propriedades urbanas, mas vigilância em zonas rurais limítrofes. \\ \textit{Aquisição de terra rural proibida para estrangeiros limítrofes.}
\end{description}
Coordenadas: -25.51°S /\allowbreak  -54.61°W. Acesso em 2026-03-20.

\textbf{Irradiância solar global --- ALLSKY SFC SW DWN (kWh/\allowbreak m²/\allowbreak dia):}

\begin{longtable}[]{@{}
  >{\raggedright\arraybackslash}p{(\linewidth - 22\tabcolsep) * \real{0.0833}}
  >{\raggedright\arraybackslash}p{(\linewidth - 22\tabcolsep) * \real{0.0833}}
  >{\raggedright\arraybackslash}p{(\linewidth - 22\tabcolsep) * \real{0.0833}}
  >{\raggedright\arraybackslash}p{(\linewidth - 22\tabcolsep) * \real{0.0833}}
  >{\raggedright\arraybackslash}p{(\linewidth - 22\tabcolsep) * \real{0.0833}}
  >{\raggedright\arraybackslash}p{(\linewidth - 22\tabcolsep) * \real{0.0833}}
  >{\raggedright\arraybackslash}p{(\linewidth - 22\tabcolsep) * \real{0.0833}}
  >{\raggedright\arraybackslash}p{(\linewidth - 22\tabcolsep) * \real{0.0833}}
  >{\raggedright\arraybackslash}p{(\linewidth - 22\tabcolsep) * \real{0.0833}}
  >{\raggedright\arraybackslash}p{(\linewidth - 22\tabcolsep) * \real{0.0833}}
  >{\raggedright\arraybackslash}p{(\linewidth - 22\tabcolsep) * \real{0.0833}}
  >{\raggedright\arraybackslash}p{(\linewidth - 22\tabcolsep) * \real{0.0833}}@{}}
\toprule\noalign{}
\begin{minipage}[b]{\linewidth}\raggedright
Jan
\end{minipage} & \begin{minipage}[b]{\linewidth}\raggedright
Fev
\end{minipage} & \begin{minipage}[b]{\linewidth}\raggedright
Mar
\end{minipage} & \begin{minipage}[b]{\linewidth}\raggedright
Abr
\end{minipage} & \begin{minipage}[b]{\linewidth}\raggedright
Mai
\end{minipage} & \begin{minipage}[b]{\linewidth}\raggedright
Jun
\end{minipage} & \begin{minipage}[b]{\linewidth}\raggedright
Jul
\end{minipage} & \begin{minipage}[b]{\linewidth}\raggedright
Ago
\end{minipage} & \begin{minipage}[b]{\linewidth}\raggedright
Set
\end{minipage} & \begin{minipage}[b]{\linewidth}\raggedright
Out
\end{minipage} & \begin{minipage}[b]{\linewidth}\raggedright
Nov
\end{minipage} & \begin{minipage}[b]{\linewidth}\raggedright
Dez
\end{minipage} \\
\midrule\noalign{}
\endhead
\bottomrule\noalign{}
\endlastfoot
6.52 & 5.97 & 5.45 & 4.50 & 3.33 & 2.93 & 3.24 & 4.04 & 4.60 & 5.28 &
6.33 & 6.55 \\
\end{longtable}

Média anual: 4.89 kWh/\allowbreak m²/\allowbreak dia. Mínimo: Jun (2.93) \textbar{} Máximo: Dez
(6.55).

\textbf{Precipitação --- PRECTOTCORR (mm/\allowbreak mês → mm/\allowbreak ano
\textasciitilde1.610):}

\begin{longtable}[]{@{}
  >{\raggedright\arraybackslash}p{(\linewidth - 22\tabcolsep) * \real{0.0833}}
  >{\raggedright\arraybackslash}p{(\linewidth - 22\tabcolsep) * \real{0.0833}}
  >{\raggedright\arraybackslash}p{(\linewidth - 22\tabcolsep) * \real{0.0833}}
  >{\raggedright\arraybackslash}p{(\linewidth - 22\tabcolsep) * \real{0.0833}}
  >{\raggedright\arraybackslash}p{(\linewidth - 22\tabcolsep) * \real{0.0833}}
  >{\raggedright\arraybackslash}p{(\linewidth - 22\tabcolsep) * \real{0.0833}}
  >{\raggedright\arraybackslash}p{(\linewidth - 22\tabcolsep) * \real{0.0833}}
  >{\raggedright\arraybackslash}p{(\linewidth - 22\tabcolsep) * \real{0.0833}}
  >{\raggedright\arraybackslash}p{(\linewidth - 22\tabcolsep) * \real{0.0833}}
  >{\raggedright\arraybackslash}p{(\linewidth - 22\tabcolsep) * \real{0.0833}}
  >{\raggedright\arraybackslash}p{(\linewidth - 22\tabcolsep) * \real{0.0833}}
  >{\raggedright\arraybackslash}p{(\linewidth - 22\tabcolsep) * \real{0.0833}}@{}}
\toprule\noalign{}
\begin{minipage}[b]{\linewidth}\raggedright
Jan
\end{minipage} & \begin{minipage}[b]{\linewidth}\raggedright
Fev
\end{minipage} & \begin{minipage}[b]{\linewidth}\raggedright
Mar
\end{minipage} & \begin{minipage}[b]{\linewidth}\raggedright
Abr
\end{minipage} & \begin{minipage}[b]{\linewidth}\raggedright
Mai
\end{minipage} & \begin{minipage}[b]{\linewidth}\raggedright
Jun
\end{minipage} & \begin{minipage}[b]{\linewidth}\raggedright
Jul
\end{minipage} & \begin{minipage}[b]{\linewidth}\raggedright
Ago
\end{minipage} & \begin{minipage}[b]{\linewidth}\raggedright
Set
\end{minipage} & \begin{minipage}[b]{\linewidth}\raggedright
Out
\end{minipage} & \begin{minipage}[b]{\linewidth}\raggedright
Nov
\end{minipage} & \begin{minipage}[b]{\linewidth}\raggedright
Dez
\end{minipage} \\
\midrule\noalign{}
\endhead
\bottomrule\noalign{}
\endlastfoot
147 & 129 & 117 & 124 & 163 & 108 & 85 & 70 & 121 & 211 & 159 & 179 \\
\end{longtable}

Estação chuvosa Out-Dez; seca Jul-Ago (mínimo em Ago).

\textbf{Inclinação solar recomendada:} 26° Norte (anual); 36° no inverno
(Jun-Ago); 16° no verão (Nov-Jan). Base: latitude -25.51°S.
\paragraph{Solo (SoilGrids 2.0, média ponderada 0--30
cm)}

\begin{longtable}[]{@{}ll@{}}
\toprule\noalign{}
Parâmetro & Valor \\
\midrule\noalign{}
\endhead
\bottomrule\noalign{}
\endlastfoot
pH (H$_2$O) & 5.16 \\
Carbono orgânico (SOC) & 25.39 g/\allowbreak kg \\
Argila & 51.67 \% \\
Areia & 15.13 \% \\
Silte (calc.) & 33.2 \% \\
Densidade aparente & 1.24 g/\allowbreak cm³ \\
\textbf{Aptidão agrícola} & \textbf{Média} \\
\end{longtable}

Fonte: ISRIC SoilGrids 2.0 via WCS (acesso em 2026-03-21). Coords:
-25.5167°, -54.6167°. Média ponderada camadas 0-5, 5-15, 15-30 cm.
\subsubsection{Indicadores Sociais}

\textbf{Fontes:} PNUD 2020, INE\footnote{Instituto Nacional de Estadística (Instituto Nacional de Estatística), órgão oficial de dados demográficos e censitários.} Censo 2022, INE\footnote{Instituto Nacional de Estadística (Instituto Nacional de Estatística), órgão oficial de dados demográficos e censitários.} EPHC 2023, Ministerio
Público 2024\\
\textbf{Nota:} Valores marcados como \emph{(dept.)} referem-se ao
departamento de Alto Paraná; valores \emph{(dist.)} são específicos
deste distrito. Dados marcados \emph{(est.)} são estimativas.

\begin{longtable}[]{@{}
  >{\raggedright\arraybackslash}p{(\linewidth - 4\tabcolsep) * \real{0.4231}}
  >{\raggedright\arraybackslash}p{(\linewidth - 4\tabcolsep) * \real{0.2692}}
  >{\raggedright\arraybackslash}p{(\linewidth - 4\tabcolsep) * \real{0.3077}}@{}}
\toprule\noalign{}
\begin{minipage}[b]{\linewidth}\raggedright
Indicador
\end{minipage} & \begin{minipage}[b]{\linewidth}\raggedright
Valor
\end{minipage} & \begin{minipage}[b]{\linewidth}\raggedright
Âmbito
\end{minipage} \\
\midrule\noalign{}
\endhead
\bottomrule\noalign{}
\endlastfoot
IDH (2020) & 0,752 & dept. \\
Ranking IDH nacional & 3/\allowbreak 18 & dept. \\
Esperança de vida & 75,1 anos & dept. \\
Escolaridade média & 10.4 anos & dist. \\
RNB per capita (USD PPA) & 13.500 & dept. \\
Pobreza monetária (\%) & 16,3\% & dept. \\
Pobreza extrema (\%) & 3,1\% & dept. \\
Índice de Gini & 0,428 & dept. \\
Acesso a água potável (\%) & 88,7\% & dept. \\
Acesso a saneamento (\%) & 87,2\% & dept. \\
Taxa de homicídios (est., /\allowbreak 100k hab) & \textasciitilde8--12 (estimativa,
acima da média nacional) & dept. \\
Índice de segurança & baixo & dept. \\
População (Censo 2022) & 325.819 hab. & dist. \\
Idade mediana & 28.0 anos & dist. \\
Taxa de fecundidade & N/\allowbreak D & dist. \\
\end{longtable}
\subsubsection{Combustível}

\textbf{Referência:} PETROPAR /\allowbreak  postos locais (2024) \textbf{Tipo de
localidade:} Capital departamental

\begin{longtable}[]{@{}lll@{}}
\toprule\noalign{}
Combustível & USD/\allowbreak litro & Gs/\allowbreak litro (aprox.) \\
\midrule\noalign{}
\endhead
\bottomrule\noalign{}
\endlastfoot
Gasolina 93 oct & 0.97 & 7,178 \\
Gasolina 97 oct (premium) & 1.07 & 7,918 \\
Diesel & 0.90 & 6,660 \\
\end{longtable}

\begin{quote}
Preços podem variar ±5\% conforme posto e sazonalidade. Chaco e interior
remoto apresentam maior variação.
\end{quote}

\subsubsection{Cobertura Celular}

\textbf{Fonte:} CONATEL PY /\allowbreak  operadoras (2024)

\begin{longtable}[]{@{}ll@{}}
\toprule\noalign{}
Parâmetro & Valor \\
\midrule\noalign{}
\endhead
\bottomrule\noalign{}
\endlastfoot
Cobertura 4G população (dept.) & 95\% \\
Cobertura 4G área rural & 88\% \\
Melhor operadora & Tigo/\allowbreak Personal \\
Qualidade rural & boa \\
\end{longtable}

\begin{quote}
Para áreas rurais fora do núcleo urbano, recomenda-se chip Tigo como
principal e Personal como backup.
\end{quote}

\subsubsection{Internet}

\textbf{Fonte:} CONATEL /\allowbreak  Speedtest Ookla (2024)

\begin{longtable}[]{@{}ll@{}}
\toprule\noalign{}
Parâmetro & Valor \\
\midrule\noalign{}
\endhead
\bottomrule\noalign{}
\endlastfoot
Velocidade média download & 85 Mbps \\
Domicílios com internet (dept.) & 87\% \\
Tecnologia predominante & fibra/\allowbreak cabo \\
Opção rural & Starlink disponível (\textasciitilde USD 44/\allowbreak mês) \\
\end{longtable}

\subsubsection{Mercado Imobiliário e Terra
Rural}

\textbf{Fonte:} INDERT /\allowbreak  Clasificados.com.py (2024)

\begin{longtable}[]{@{}ll@{}}
\toprule\noalign{}
Tipo & Referência \\
\midrule\noalign{}
\endhead
\bottomrule\noalign{}
\endlastfoot
Terra agrícola alta prod. (USD/\allowbreak ha) & 14,800 \\
Imóvel urbano (USD/\allowbreak m²) & 1,609 \\
Aluguel 2 quartos (USD/\allowbreak mês) & 720 \\
\end{longtable}

\begin{quote}
Valores de referência departamental. Localidades menores podem ter
preços 20--40\% abaixo da capital departamental. \#\#\# Saúde
\end{quote}

\textbf{Fonte:} MSPBS /\allowbreak  IPS Paraguay (2024-2026), consolidação
departamental e proxy local

\begin{longtable}[]{@{}ll@{}}
\toprule\noalign{}
Serviço & Disponibilidade \\
\midrule\noalign{}
\endhead
\bottomrule\noalign{}
\endlastfoot
USF /\allowbreak  Posto de Saúde & sim \\
Hospital Regional & sim \\
IPS (seguro social) & sim \\
Farmácia & sim \\
Distância ao hospital de referência & local \\
\end{longtable}

\textbf{Principais estabelecimentos:} Hospital distrital /\allowbreak  regional de
Ciudad del Este; rede de USF e farmacias locais

\textbf{Observação para imigrantes:} Centro de referencia do
departamento. Acesso a atendimento primario e, em geral, a maior oferta
publica e privada da area. Cobertura privada continua recomendada para
especialidades e urgências de maior complexidade.

\paragraph{Solo (SoilGrids 2.0, média ponderada 0--30
cm)}

\begin{longtable}[]{@{}ll@{}}
\toprule\noalign{}
Parâmetro & Valor \\
\midrule\noalign{}
\endhead
\bottomrule\noalign{}
\endlastfoot
pH (H$_2$O) & 5.16 \\
Carbono orgânico (SOC) & 25.39 g/\allowbreak kg \\
Argila & 51.67 \% \\
Areia & 15.13 \% \\
Silte (calc.) & 33.2 \% \\
Densidade aparente & 1.24 g/\allowbreak cm³ \\
\textbf{Aptidão agrícola} & \textbf{Média} \\
\end{longtable}

Fonte: ISRIC SoilGrids 2.0 via WCS (acesso em 2026-03-21). Coords:
-25.5167°, -54.6167°. Média ponderada camadas 0-5, 5-15, 15-30 cm.
\newpage
\secaoDiagnostico{Doctor Raul Pena}{\mapaConteudo{-25.9166}{-55.1833}}{Doctor Raúl Peña é um enclave de alta produtividade no sul de Alto
Paraná. Oferece o máximo em termos de qualidade de solo e segurança
social rural, sendo uma localidade ``fora do radar'' das grandes crises
demográficas. Sua viabilidade como refúgio é sustentada pela riqueza
agrícola, mas exige do ocupante um alto nível de autonomia tática,
energética e de saúde, dada a simplicidade da infraestrutura
institucional local e a perigosa dependência da monocultura exportadora.

\subsubsection{Por que sim}

\begin{itemize}
\tightlist
\item
  Solos de elite mundial e clima favorável para policultivos.
\item
  Ambiente social extremamente seguro, ordeiro e pacífico.
\item
  Baixa densidade populacional rural (privacidade absoluta).
\item
  Sem restrições da Lei de Fronteira.
\end{itemize}

\subsubsection{Por que não}

\begin{itemize}
\tightlist
\item
  Infraestrutura de saúde local limitada a atendimentos primários.
\item
  Dependência econômica total de insumos e mercados globais.
\item
  Distância de centros comerciais e financeiros diversificados.
\end{itemize}

\subsubsection{Recomendação
Técnica}

Altamente indicado para estabelecimentos de autossuficiência tecnificada
ou exploração agrícola que busque isolamento demográfico em solo de alta
performance. Requer autonomia em logística básica e saúde. O GSS de 5.6
reflete a segurança do isolamento em contraste com a vulnerabilidade
sistêmica da especialização agrícola. Referencia atual de score: GSS 5.6
em 2026-03-05.

\subsubsection{}

\begin{itemize}
\tightlist
\item
  \begin{enumerate}
  \def\labelenumi{\arabic{enumi})}
  \tightlist
  \item
    Dados censitários validados (INE\footnote{Instituto Nacional de Estadística (Instituto Nacional de Estatística), órgão oficial de dados demográficos e censitários.} 2022).
  \end{enumerate}
\item
  \begin{enumerate}
  \def\labelenumi{\arabic{enumi})}
  \setcounter{enumi}{1}
  \tightlist
  \item
    Relatórios de produtividade agrícola de elite (MAG).
  \end{enumerate}
\item
  \begin{enumerate}
  \def\labelenumi{\arabic{enumi})}
  \setcounter{enumi}{2}
  \tightlist
  \item
    Mapa de solos e aptidão do Sul de Alto Paraná (MAG).
  \end{enumerate}
\item
  \begin{enumerate}
  \def\labelenumi{\arabic{enumi})}
  \setcounter{enumi}{3}
  \tightlist
  \item
    Registro de segurança pública departamental.
  \end{enumerate}
\end{itemize}}
\begin{table}[ht!]
\small \begin{tabular}{p{3.0cm}p{0.8cm}p{6.0cm}} \toprule
\textbf{Critério} & \textbf{Nota} & \textbf{Justificativa} \\ \midrule
A - Ameaças Estratégicas & 6.0 & Localização discreta e protegida de eixos de conflito; excelente invisibilidade tática \\
B - Risco Social & 7.0 & População pequena e densidade controlada; baixíssimo risco social urbano \\
C - Riscos Naturais & 7.0 & Geologia muito estável e baixo risco de desastres naturais graves \\
D - Autossuficiência & 7.5 & Recursos agrícolas de elite e abundância hídrica; limitado pela especialização técnica \\
E - Institucional & 2.0 & Nota penalizada no ranking nacional para refletir a dependência de insumos globais e infraestrutura local mínima - GSS 5.6 total \\
\midrule \textbf{GSS FINAL} & \textbf{5.6} & \textbf{Moderadamente Seguro (Ideal para perfis agrícolas que planejem viver com autonomia total de suprimentos).} \\ \bottomrule \end{tabular}
\end{table}
\subsubsection{Vulnerabilidades e
Mitigação}\label{vuln-Doctor_Raul_Pena-vulnerabilidades-e-mitigauxe7uxe3o}

\begin{itemize}
\tightlist
\item
  \textbf{Dependência de Cadeias Globais:} Economia focada
  exclusivamente em commodities exportáveis (Mitigação: implementação de
  sistemas de policultivo de subsistência e criação de pequenos
  animais).
\item
  \textbf{Isolamento de Saúde:} Distância significativa de hospitais
  cirúrgicos (Mitigação: manutenção de estoque médico básico e kit
  trauma local).
\item
  \textbf{Fragilidade da Rede Elétrica:} Dependência de energia externa
  para silos e bombas (Mitigação: instalação de sistemas solares
  fotovoltaicos potentes e geradores).
\end{itemize}
\subsubsection{Dossiê de Campo}\label{dossie-Doctor_Raul_Pena-dossiuxea-de-campo}
\begin{description}[style=nextline,leftmargin=0.5cm,font=\bfseries]
\item[Coordenadas]  25°55'S 55°11'W.
\item[Topografia]  Relevo predominantemente plano a suavemente ondulado, com solos altamente férteis e mecanizáveis. Altitude \textasciitilde 240m. Área de \textasciitilde 225 km². Geologicamente muito estável, risco sísmico desprezível.
\item[Alvos Estratégicos]  Áreas de produção agrícola mecanizada de alta performance; Silos de armazenamento de grãos. Localização estratégica no sul do departamento, servindo como ponto de conexão secundário entre Alto Paraná e Itapúa. Ausência de alvos militares diretos, oferecendo excelente invisibilidade tática.
\item[Fallout]  Ventos predominantes do Norte (NE) e Leste. Baixo risco de fallout direto.
\item[População]  10.048 habitantes (Censo 2022 Final). Perfil demográfico majoritariamente rural (85\%) com forte presença de imigrantes e descendentes de origem brasileira.
\item[Densidade]  \textasciitilde 44,6 hab/\allowbreak km² (Densidade baixa, oferecendo isolamento relativo em áreas de fazendas mecanizadas).
\item[Vias de Acesso]  Acesso via ramais pavimentados e estradas de terra que conectam à Ruta PY06. Localização que favorece a discrição, estando protegida dos grandes fluxos migratórios pela topografia e distância das rotas principais.
\item[Serviços]  Unidade de Saúde da Família (USF) local. Dependência crítica de Naranjal (20 km) ou Santa Rita (45 km) para serviços de saúde de alta complexidade e logística comercial avançada.
\item[Custo de Vida]  Baixo; economia estritamente vinculada ao agronegócio exportador.
\item[Preço da Terra]  US\$ 10.000-15.000/\allowbreak ha (terras mecanizadas de elite - "Terra Roxa"); um dos valores mais altos do país devido à produtividade.
\item[Hidrologia]  Baixo risco de inundações sistêmicas catastróficas. Presença de arroios locais que favorecem a irrigação natural. Solo com alta capacidade de absorção, mas que exige terraceamento rigoroso contra erosão.
\item[Clima]  Subtropical Úmido. Verões quentes; invernos frescos com geadas ocasionais que beneficiam a cultura do trigo.
\item[Energia]  Rede nacional (Itaipu); instável em áreas rurais remotas durante picos de demanda agrícola.
\item[Água]  Abundante via poços artesianos profundos e lençol freático produtivo (Sistema Aquífero Guarani).
\item[Qualidade do Solo]  Latossolo Vermelho de Elite; solo profundo, estruturado e de altíssima fertilidade natural. Excelente aptidão para soja, milho e trigo.
\item[Resources Locais]  Grande produção de soja, milho e trigo. Elevado potencial para autossuficiência calórica per capita, mas dependente de tecnologia e insumos externos para manutenção da produtividade.
\item[Segurança]  Zona rural extremamente pacífica e ordeira. Baixíssima criminalidade urbana. Coesão social baseada na cultura do trabalho agrícola e em laços comunitários de imigrantes produtores. Baixa visibilidade para conflitos sociais de massa.
\item[Leis Local: Município de criação recente (2012), emancipado de San Cristóbal. Livre de Restrição de Fronteira]  Localizado fora da faixa de 50km da fronteira internacional seca.
\end{description}
\paragraph{Solo (SoilGrids 2.0, média ponderada 0--30
cm)}

\begin{longtable}[]{@{}ll@{}}
\toprule\noalign{}
Parâmetro & Valor \\
\midrule\noalign{}
\endhead
\bottomrule\noalign{}
\endlastfoot
pH (H$_2$O) & 5.4 \\
Carbono orgânico (SOC) & 20.82 g/\allowbreak kg \\
Argila & 41.76 \% \\
Areia & 26.3 \% \\
Silte (calc.) & 31.9 \% \\
Densidade aparente & 1.24 g/\allowbreak cm³ \\
\textbf{Aptidão agrícola} & \textbf{Média} \\
\end{longtable}

Fonte: ISRIC SoilGrids 2.0 via WCS (acesso em 2026-03-21). Coords:
-25.9167°, -55.1833°. Média ponderada camadas 0-5, 5-15, 15-30 cm.

\begin{itemize}
\tightlist
\item
  \textbf{Homicidios/\allowbreak 100k:} \textasciitilde8,0 (Regional/\allowbreak Alto Paraná).
\end{itemize}

\subsection{10. Analise de Riscos}

\begin{itemize}
\tightlist
\item
  \textbf{Cenario curto prazo:} Manutenção da rentabilidade das safras e
  valorização de áreas mecanizadas.
\item
  \textbf{Cenario medio prazo:} Impacto de crises no mercado global de
  fertilizantes sobre a produtividade local.
\end{itemize}

\subsection{11. Redacao Final}

\subsubsection{Diagnostico Integrado}

Doctor Raúl Peña é um enclave de alta produtividade no sul de Alto
Paraná. Oferece o máximo em termos de qualidade de solo e segurança
social rural, sendo uma localidade ``fora do radar'' das grandes crises
demográficas. Sua viabilidade como refúgio é sustentada pela riqueza
agrícola, mas exige do ocupante um alto nível de autonomia tática,
energética e de saúde, dada a simplicidade da infraestrutura
institucional local e a perigosa dependência da monocultura exportadora.

\subsubsection{Por que sim}

\begin{itemize}
\tightlist
\item
  Solos de elite mundial e clima favorável para policultivos.
\item
  Ambiente social extremamente seguro, ordeiro e pacífico.
\item
  Baixa densidade populacional rural (privacidade absoluta).
\item
  Sem restrições da Lei de Fronteira.
\end{itemize}

\subsubsection{Por que nao}

\begin{itemize}
\tightlist
\item
  Infraestrutura de saúde local limitada a atendimentos primários.
\item
  Dependência econômica total de insumos e mercados globais.
\item
  Distância de centros comerciais e financeiros diversificados.
\end{itemize}

\subsubsection{Recomendacao Tecnica}

Altamente indicado para estabelecimentos de autossuficiência tecnificada
ou exploração agrícola que busque isolamento demográfico em solo de alta
performance. Requer autonomia em logística básica e saúde. O GSS de 5.6
reflete a segurança do isolamento em contraste com a vulnerabilidade
sistêmica da especialização agrícola. Referencia atual de score: GSS 5.6
em 2026-03-05.

\subsection{13.}

\begin{itemize}
\tightlist
\item
  \begin{enumerate}
  \def\labelenumi{\arabic{enumi})}
  \tightlist
  \item
    Dados censitários validados (INE\footnote{Instituto Nacional de Estadística (Instituto Nacional de Estatística), órgão oficial de dados demográficos e censitários.} 2022).
  \end{enumerate}
\item
  \begin{enumerate}
  \def\labelenumi{\arabic{enumi})}
  \setcounter{enumi}{1}
  \tightlist
  \item
    Relatórios de produtividade agrícola de elite (MAG).
  \end{enumerate}
\item
  \begin{enumerate}
  \def\labelenumi{\arabic{enumi})}
  \setcounter{enumi}{2}
  \tightlist
  \item
    Mapa de solos e aptidão do Sul de Alto Paraná (MAG).
  \end{enumerate}
\item
  \begin{enumerate}
  \def\labelenumi{\arabic{enumi})}
  \setcounter{enumi}{3}
  \tightlist
  \item
    Registro de segurança pública departamental.
  \end{enumerate}
\end{itemize}

\subsubsection{Combustível}

\textbf{Referência:} PETROPAR /\allowbreak  postos locais (2024) \textbf{Tipo de
localidade:} Interior

\begin{longtable}[]{@{}lll@{}}
\toprule\noalign{}
Combustível & USD/\allowbreak litro & Gs/\allowbreak litro (aprox.) \\
\midrule\noalign{}
\endhead
\bottomrule\noalign{}
\endlastfoot
Gasolina 93 oct & 0.97 & 7,178 \\
Gasolina 97 oct (premium) & 1.07 & 7,918 \\
Diesel & 0.90 & 6,660 \\
\end{longtable}

\begin{quote}
Preços podem variar ±5\% conforme posto e sazonalidade. Chaco e interior
remoto apresentam maior variação.
\end{quote}

\subsubsection{Cobertura Celular}

\textbf{Fonte:} CONATEL PY /\allowbreak  operadoras (2024)

\begin{longtable}[]{@{}ll@{}}
\toprule\noalign{}
Parâmetro & Valor \\
\midrule\noalign{}
\endhead
\bottomrule\noalign{}
\endlastfoot
Cobertura 4G população (dept.) & 95\% \\
Cobertura 4G área rural & 88\% \\
Melhor operadora & Tigo/\allowbreak Personal \\
Qualidade rural & boa \\
\end{longtable}

\begin{quote}
Para áreas rurais fora do núcleo urbano, recomenda-se chip Tigo como
principal e Personal como backup.
\end{quote}

\subsubsection{Internet}

\textbf{Fonte:} CONATEL /\allowbreak  Speedtest Ookla (2024)

\begin{longtable}[]{@{}ll@{}}
\toprule\noalign{}
Parâmetro & Valor \\
\midrule\noalign{}
\endhead
\bottomrule\noalign{}
\endlastfoot
Velocidade média download & 85 Mbps \\
Domicílios com internet (dept.) & 87\% \\
Tecnologia predominante & fibra/\allowbreak cabo \\
Opção rural & Starlink disponível (\textasciitilde USD 44/\allowbreak mês) \\
\end{longtable}

\subsubsection{Mercado Imobiliário e Terra
Rural}

\textbf{Fonte:} INDERT /\allowbreak  Clasificados.com.py (2024)

\begin{longtable}[]{@{}ll@{}}
\toprule\noalign{}
Tipo & Referência \\
\midrule\noalign{}
\endhead
\bottomrule\noalign{}
\endlastfoot
Terra agrícola alta prod. (USD/\allowbreak ha) & 14,800 \\
Imóvel urbano (USD/\allowbreak m²) & 1,400 \\
Aluguel 2 quartos (USD/\allowbreak mês) & 600 \\
\end{longtable}

\begin{quote}
Valores de referência departamental. Localidades menores podem ter
preços 20--40\% abaixo da capital departamental. \#\#\# Saúde
\end{quote}

\textbf{Fonte:} MSPBS /\allowbreak  IPS Paraguay (2024-2026), consolidação
departamental e proxy local

\begin{longtable}[]{@{}ll@{}}
\toprule\noalign{}
Serviço & Disponibilidade \\
\midrule\noalign{}
\endhead
\bottomrule\noalign{}
\endlastfoot
USF /\allowbreak  Posto de Saúde & sim \\
Hospital Regional & não \\
IPS (seguro social) & sim \\
Farmácia & sim \\
Distância ao hospital de referência & \textasciitilde5 km (Ciudad del
Este) \\
\end{longtable}

\textbf{Principais estabelecimentos:} USF local; referencia hospitalar
em Ciudad del Este

\textbf{Observação para imigrantes:} Atendimento primario local ou em
raio curto; casos de maior complexidade seguem para Ciudad del Este.
Cobertura privada continua recomendada para especialidades e urgências
de maior complexidade.
\subsubsection{Indicadores Sociais}

\textbf{Fontes:} PNUD 2020, INE\footnote{Instituto Nacional de Estadística (Instituto Nacional de Estatística), órgão oficial de dados demográficos e censitários.} Censo 2022, INE\footnote{Instituto Nacional de Estadística (Instituto Nacional de Estatística), órgão oficial de dados demográficos e censitários.} EPHC 2023, Ministerio
Público 2024\\
\textbf{Nota:} Valores marcados como \emph{(dept.)} referem-se ao
departamento de Alto Paraná; valores \emph{(dist.)} são específicos
deste distrito. Dados marcados \emph{(est.)} são estimativas.

\begin{longtable}[]{@{}
  >{\raggedright\arraybackslash}p{(\linewidth - 4\tabcolsep) * \real{0.4231}}
  >{\raggedright\arraybackslash}p{(\linewidth - 4\tabcolsep) * \real{0.2692}}
  >{\raggedright\arraybackslash}p{(\linewidth - 4\tabcolsep) * \real{0.3077}}@{}}
\toprule\noalign{}
\begin{minipage}[b]{\linewidth}\raggedright
Indicador
\end{minipage} & \begin{minipage}[b]{\linewidth}\raggedright
Valor
\end{minipage} & \begin{minipage}[b]{\linewidth}\raggedright
Âmbito
\end{minipage} \\
\midrule\noalign{}
\endhead
\bottomrule\noalign{}
\endlastfoot
IDH (2020) & 0,752 & dept. \\
Ranking IDH nacional & 3/\allowbreak 18 & dept. \\
Esperança de vida & 75,1 anos & dept. \\
Escolaridade média & 9,6 anos & dept. \\
RNB per capita (USD PPA) & 13.500 & dept. \\
Pobreza monetária (\%) & 16,3\% & dept. \\
Pobreza extrema (\%) & 3,1\% & dept. \\
Índice de Gini & 0,428 & dept. \\
Acesso a água potável (\%) & 88,7\% & dept. \\
Acesso a saneamento (\%) & 87,2\% & dept. \\
Taxa de homicídios (est., /\allowbreak 100k hab) & \textasciitilde8--12 (estimativa,
acima da média nacional) & dept. \\
Índice de segurança & baixo & dept. \\
População (Censo 2022) & 5.617 hab. & dist. \\
Idade mediana & 28.0 anos & dist. \\
Taxa de fecundidade & N/\allowbreak D & dist. \\
\end{longtable}
\subsubsection{Combustível}

\textbf{Referência:} PETROPAR /\allowbreak  postos locais (2024) \textbf{Tipo de
localidade:} Interior

\begin{longtable}[]{@{}lll@{}}
\toprule\noalign{}
Combustível & USD/\allowbreak litro & Gs/\allowbreak litro (aprox.) \\
\midrule\noalign{}
\endhead
\bottomrule\noalign{}
\endlastfoot
Gasolina 93 oct & 0.97 & 7,178 \\
Gasolina 97 oct (premium) & 1.07 & 7,918 \\
Diesel & 0.90 & 6,660 \\
\end{longtable}

\begin{quote}
Preços podem variar ±5\% conforme posto e sazonalidade. Chaco e interior
remoto apresentam maior variação.
\end{quote}

\subsubsection{Cobertura Celular}

\textbf{Fonte:} CONATEL PY /\allowbreak  operadoras (2024)

\begin{longtable}[]{@{}ll@{}}
\toprule\noalign{}
Parâmetro & Valor \\
\midrule\noalign{}
\endhead
\bottomrule\noalign{}
\endlastfoot
Cobertura 4G população (dept.) & 95\% \\
Cobertura 4G área rural & 88\% \\
Melhor operadora & Tigo/\allowbreak Personal \\
Qualidade rural & boa \\
\end{longtable}

\begin{quote}
Para áreas rurais fora do núcleo urbano, recomenda-se chip Tigo como
principal e Personal como backup.
\end{quote}

\subsubsection{Internet}

\textbf{Fonte:} CONATEL /\allowbreak  Speedtest Ookla (2024)

\begin{longtable}[]{@{}ll@{}}
\toprule\noalign{}
Parâmetro & Valor \\
\midrule\noalign{}
\endhead
\bottomrule\noalign{}
\endlastfoot
Velocidade média download & 85 Mbps \\
Domicílios com internet (dept.) & 87\% \\
Tecnologia predominante & fibra/\allowbreak cabo \\
Opção rural & Starlink disponível (\textasciitilde USD 44/\allowbreak mês) \\
\end{longtable}

\subsubsection{Mercado Imobiliário e Terra
Rural}

\textbf{Fonte:} INDERT /\allowbreak  Clasificados.com.py (2024)

\begin{longtable}[]{@{}ll@{}}
\toprule\noalign{}
Tipo & Referência \\
\midrule\noalign{}
\endhead
\bottomrule\noalign{}
\endlastfoot
Terra agrícola alta prod. (USD/\allowbreak ha) & 14,800 \\
Imóvel urbano (USD/\allowbreak m²) & 1,400 \\
Aluguel 2 quartos (USD/\allowbreak mês) & 600 \\
\end{longtable}

\begin{quote}
Valores de referência departamental. Localidades menores podem ter
preços 20--40\% abaixo da capital departamental. \#\#\# Saúde
\end{quote}

\textbf{Fonte:} MSPBS /\allowbreak  IPS Paraguay (2024-2026), consolidação
departamental e proxy local

\begin{longtable}[]{@{}ll@{}}
\toprule\noalign{}
Serviço & Disponibilidade \\
\midrule\noalign{}
\endhead
\bottomrule\noalign{}
\endlastfoot
USF /\allowbreak  Posto de Saúde & sim \\
Hospital Regional & não \\
IPS (seguro social) & sim \\
Farmácia & sim \\
Distância ao hospital de referência & \textasciitilde5 km (Ciudad del
Este) \\
\end{longtable}

\textbf{Principais estabelecimentos:} USF local; referencia hospitalar
em Ciudad del Este

\textbf{Observação para imigrantes:} Atendimento primario local ou em
raio curto; casos de maior complexidade seguem para Ciudad del Este.
Cobertura privada continua recomendada para especialidades e urgências
de maior complexidade.
\newpage
\secaoDiagnostico{Domingo Martinez de Irala}{\mapaConteudo{-25.8833}{-54.6166}}{Domingo Martínez de Irala é um dos segredos mais bem guardados de Alto
Paraná. Oferece solos de produtividade mundial e um isolamento
geográfico e demográfico que garante privacidade absoluta, protegida
pela barreira do Rio Paraná. Sua força reside na excepcional abundância
de recursos hídricos e na segurança social de uma comunidade pequena e
pacífica. É o local ideal para quem busca estabelecer um refúgio tático
com acesso fluvial, transformando o distanciamento metropolitano em uma
barreira natural de proteção contra crises sistêmicas externas.

\subsubsection{Por que sim}

\begin{itemize}
\tightlist
\item
  Máximo isolamento geográfico e demográfico (anonimato absoluto).
\item
  Solo de altíssima produtividade (``terra roxa'') e abundância de águas
  puras.
\item
  Ambiente social pacífico e estável sob vigilância naval.
\item
  Longe de qualquer alvo militar terrestre primário ou eixo de fallout.
\end{itemize}

\subsubsection{Por que não}

\begin{itemize}
\tightlist
\item
  Risco de isolamento terrestre por chuvas intensas em estradas de
  terra.
\item
  Dependência de grandes centros para serviços de saúde complexos.
\item
  Restrições da Lei de Fronteira para investidores estrangeiros.
\end{itemize}

\subsubsection{Recomendação
Técnica}

Altamente indicado para residências de isolamento tático com
infraestrutura de transporte fluvial. Requer autonomia energética
individual para mitigar riscos de rede rural. O GSS de 7.9 reflete a
solidez do isolamento geográfico e a riqueza de recursos naturais da
localidade. Referencia atual de score: GSS 7.9 em 2026-03-05.

\subsubsection{}

\begin{itemize}
\tightlist
\item
  \begin{enumerate}
  \def\labelenumi{\arabic{enumi})}
  \tightlist
  \item
    Dados censitários validados (INE\footnote{Instituto Nacional de Estadística (Instituto Nacional de Estatística), órgão oficial de dados demográficos e censitários.} 2022).
  \end{enumerate}
\item
  \begin{enumerate}
  \def\labelenumi{\arabic{enumi})}
  \setcounter{enumi}{1}
  \tightlist
  \item
    Relatórios de monitoramento fluvial do Rio Paraná (ANNP).
  \end{enumerate}
\item
  \begin{enumerate}
  \def\labelenumi{\arabic{enumi})}
  \setcounter{enumi}{2}
  \tightlist
  \item
    Mapa de solos e bacias hídricas de Alto Paraná (MAG).
  \end{enumerate}
\item
  \begin{enumerate}
  \def\labelenumi{\arabic{enumi})}
  \setcounter{enumi}{3}
  \tightlist
  \item
    Registro de segurança pública e governança municipal.
  \end{enumerate}
\end{itemize}}
\begin{table}[ht!]
\small \begin{tabular}{p{3.0cm}p{0.8cm}p{6.0cm}} \toprule
\textbf{Critério} & \textbf{Nota} & \textbf{Justificativa} \\ \midrule
A - Ameaças Estratégicas & 7.5 & Posição tática protegida pela geografia; margem do rio e isolamento geográfico superior; baixa visibilidade de alvos \\
B - Risco Social & 8.5 & Densidade populacional extremamente baixa; isolamento social de alto nível e privacidade garantida \\
C - Riscos Naturais & 7.0 & Geologia estável e relevo seguro contra inundações sistêmicas \\
D - Autossuficiência & 9.5 & Recursos hídricos e de proteína animal massivos; solo de altíssima produtividade \\
E - Institucional & 7.5 & Estabilidade institucional e segurança regional funcional em ambiente rural \\
\midrule \textbf{GSS FINAL} & \textbf{7.9} & \textbf{Seguro (Altamente recomendado para quem busca isolamento tático rural com abundância de recursos básicos e solo de alta performance).} \\ \bottomrule \end{tabular}
\end{table}
\subsubsection{Vulnerabilidades e
Mitigação}\label{vuln-Domingo_Martinez_de_Irala-vulnerabilidades-e-mitigauxe7uxe3o}

\begin{itemize}
\tightlist
\item
  \textbf{Isolamento Terrestre Sazonal:} Risco de isolamento por chuvas
  (Mitigação: posse de embarcação a motor como via de escape alternativa
  e manutenção de estoque de suprimentos para 60 dias).
\item
  \textbf{Fragilidade Médica:} Necessidade de deslocamento para CDE
  (Mitigação: manutenção de estoque médico básico e kit trauma local).
\item
  \textbf{Dependência de Redes:} Vulnerabilidade a quedas elétricas em
  tempestades (Mitigação: instalação de sistemas solares fotovoltaicos).
\end{itemize}
\subsubsection{Dossiê de Campo}\label{dossie-Domingo_Martinez_de_Irala-dossiuxea-de-campo}
\begin{description}[style=nextline,leftmargin=0.5cm,font=\bfseries]
\item[Coordenadas]  25°53'S 54°37'W.
\item[Topografia]  Relevo ondulado a acidentado, margeado pelo Rio Paraná ao leste (fronteira com Argentina). Altitude \textasciitilde 180m. Área de \textasciitilde 334 km². Geologicamente estável, risco sísmico nulo.
\item[Alvos Estratégicos]  Porto fluvial secundário; Áreas de pecuária e agricultura extensiva; Localização estratégica para o controle do tráfego fluvial no Alto Paraná. Ausência total de alvos militares ou infraestrutura industrial massiva, oferecendo um dos maiores níveis de invisibilidade estratégica e proteção por isolamento geográfico deliberado no sul do departamento.
\item[Fallout]  Ventos predominantes do Norte (NE) e Leste. Baixíssimo risco de fallout de alvos continentais primários devido ao isolamento tático profundo.
\item[População]  5.035 habitantes (Censo 2022 Final). População estável com forte tradição rural e ligação com a atividade pecuária e comercial fronteiriça fluvial.
\item[Densidade]  \textasciitilde 15,1 hab/\allowbreak km² (Densidade extremamente baixa, garantindo privacidade absoluta e ausência de riscos sociais de massa urbana).
\item[Vias de Acesso]  Acesso via ramais de terra a partir da Ruta PY07. Localização remota que favorece o anonimato estratégico, mas sujeita a períodos de isolamento físico devido à precariedade das vias vicinais durante chuvas intensas. Acesso fluvial funcional.
\item[Serviços]  Infraestrutura básica funcional (USF local). Dependência de Santa Rosa del Monday ou Ciudad del Este para atendimentos de saúde de alta complexidade e logística comercial especializada.
\item[Custo de Vida]  Muito baixo; economia baseada na pecuária bovina e pesca comercial.
\item[Preço da Terra]  US\$ 5.000-8.000/\allowbreak ha (áreas ribeirinhas/\allowbreak pecuária); US\$ 3.000-5.000/\allowbreak ha (campos brutos).
\item[Hidrologia]  Baixo risco de inundações sistêmicas devido à topografia elevada em relação ao Rio Paraná. Abundância hídrica superficial absoluta. Diversos arroios perenes de alta qualidade hídrica cruzam o distrito.
\item[Clima]  Subtropical Úmido. Verões quentes e chuvosos; invernos amenos com geadas ocasionais.
\item[Energia]  Rede nacional (Itaipu); instável em áreas rurais remotas.
\item[Água]  Abundância hídrica absoluta via Rio Paraná e poços artesianos profundos. Água superficial exige tratamento para consumo humano.
\item[Qualidade do Solo]  Latossolo Vermelho de Alta Performance; solo profundo e fértil, ideal para pastagens naturais, silvicultura e agricultura de subsistência.
\item[Resources Locais]  Grande produção de gado bovino e recursos pesqueiros fluviais. Elevadíssimo potencial para autossuficiência alimentar básica em nível de propriedade devido à vasta área produtiva e baixíssima população.
\item[Segurança]  Localidade excepcionalmente pacífica e segura. Criminalidade urbana praticamente inexistente. Estabilidade institucional básica baseada na presença de postos de controle fluvial e coesão comunitária tradicional. Localidade "invisível" para qualquer conflito sociopolítico nacional.
\item[Leis Local: Município consolidado. Restrição de Fronteira (Lei 2532/\allowbreak 05)]  Aplicável a terras rurais na faixa de 50km da fronteira argentina (limite fluvial).
\end{description}
\paragraph{Solo (SoilGrids 2.0, média ponderada 0--30
cm)}

\begin{longtable}[]{@{}ll@{}}
\toprule\noalign{}
Parâmetro & Valor \\
\midrule\noalign{}
\endhead
\bottomrule\noalign{}
\endlastfoot
pH (H$_2$O) & 5.3 \\
Carbono orgânico (SOC) & 22.74 g/\allowbreak kg \\
Argila & 41.52 \% \\
Areia & 20.32 \% \\
Silte (calc.) & 38.2 \% \\
Densidade aparente & 1.17 g/\allowbreak cm³ \\
\textbf{Aptidão agrícola} & \textbf{Média} \\
\end{longtable}

Fonte: ISRIC SoilGrids 2.0 via WCS (acesso em 2026-03-21). Coords:
-25.8833°, -54.6167°. Média ponderada camadas 0-5, 5-15, 15-30 cm.

\begin{itemize}
\tightlist
\item
  \textbf{Homicidios/\allowbreak 100k:} \textasciitilde0,0 (Histórico de paz
  absoluta).
\end{itemize}

\subsection{10. Analise de Riscos}

\begin{itemize}
\tightlist
\item
  \textbf{Cenario curto prazo:} Manutenção da tranquilidade rural e
  estabilidade da pecuária.
\item
  \textbf{Cenario medio prazo:} Impacto de novas rotas asfálticas
  vicinais na valorização estratégica do território.
\end{itemize}

\subsection{11. Redacao Final}

\subsubsection{Diagnostico Integrado}

Domingo Martínez de Irala é um dos segredos mais bem guardados de Alto
Paraná. Oferece solos de produtividade mundial e um isolamento
geográfico e demográfico que garante privacidade absoluta, protegida
pela barreira do Rio Paraná. Sua força reside na excepcional abundância
de recursos hídricos e na segurança social de uma comunidade pequena e
pacífica. É o local ideal para quem busca estabelecer um refúgio tático
com acesso fluvial, transformando o distanciamento metropolitano em uma
barreira natural de proteção contra crises sistêmicas externas.

\subsubsection{Por que sim}

\begin{itemize}
\tightlist
\item
  Máximo isolamento geográfico e demográfico (anonimato absoluto).
\item
  Solo de altíssima produtividade (``terra roxa'') e abundância de águas
  puras.
\item
  Ambiente social pacífico e estável sob vigilância naval.
\item
  Longe de qualquer alvo militar terrestre primário ou eixo de fallout.
\end{itemize}

\subsubsection{Por que nao}

\begin{itemize}
\tightlist
\item
  Risco de isolamento terrestre por chuvas intensas em estradas de
  terra.
\item
  Dependência de grandes centros para serviços de saúde complexos.
\item
  Restrições da Lei de Fronteira para investidores estrangeiros.
\end{itemize}

\subsubsection{Recomendacao Tecnica}

Altamente indicado para residências de isolamento tático com
infraestrutura de transporte fluvial. Requer autonomia energética
individual para mitigar riscos de rede rural. O GSS de 7.9 reflete a
solidez do isolamento geográfico e a riqueza de recursos naturais da
localidade. Referencia atual de score: GSS 7.9 em 2026-03-05.

\subsection{13.}

\begin{itemize}
\tightlist
\item
  \begin{enumerate}
  \def\labelenumi{\arabic{enumi})}
  \tightlist
  \item
    Dados censitários validados (INE\footnote{Instituto Nacional de Estadística (Instituto Nacional de Estatística), órgão oficial de dados demográficos e censitários.} 2022).
  \end{enumerate}
\item
  \begin{enumerate}
  \def\labelenumi{\arabic{enumi})}
  \setcounter{enumi}{1}
  \tightlist
  \item
    Relatórios de monitoramento fluvial do Rio Paraná (ANNP).
  \end{enumerate}
\item
  \begin{enumerate}
  \def\labelenumi{\arabic{enumi})}
  \setcounter{enumi}{2}
  \tightlist
  \item
    Mapa de solos e bacias hídricas de Alto Paraná (MAG).
  \end{enumerate}
\item
  \begin{enumerate}
  \def\labelenumi{\arabic{enumi})}
  \setcounter{enumi}{3}
  \tightlist
  \item
    Registro de segurança pública e governança municipal.
  \end{enumerate}
\end{itemize}

\subsubsection{Combustível}

\textbf{Referência:} PETROPAR /\allowbreak  postos locais (2024) \textbf{Tipo de
localidade:} Interior

\begin{longtable}[]{@{}lll@{}}
\toprule\noalign{}
Combustível & USD/\allowbreak litro & Gs/\allowbreak litro (aprox.) \\
\midrule\noalign{}
\endhead
\bottomrule\noalign{}
\endlastfoot
Gasolina 93 oct & 0.97 & 7,178 \\
Gasolina 97 oct (premium) & 1.07 & 7,918 \\
Diesel & 0.90 & 6,660 \\
\end{longtable}

\begin{quote}
Preços podem variar ±5\% conforme posto e sazonalidade. Chaco e interior
remoto apresentam maior variação.
\end{quote}

\subsubsection{Cobertura Celular}

\textbf{Fonte:} CONATEL PY /\allowbreak  operadoras (2024)

\begin{longtable}[]{@{}ll@{}}
\toprule\noalign{}
Parâmetro & Valor \\
\midrule\noalign{}
\endhead
\bottomrule\noalign{}
\endlastfoot
Cobertura 4G população (dept.) & 95\% \\
Cobertura 4G área rural & 88\% \\
Melhor operadora & Tigo/\allowbreak Personal \\
Qualidade rural & boa \\
\end{longtable}

\begin{quote}
Para áreas rurais fora do núcleo urbano, recomenda-se chip Tigo como
principal e Personal como backup.
\end{quote}

\subsubsection{Internet}

\textbf{Fonte:} CONATEL /\allowbreak  Speedtest Ookla (2024)

\begin{longtable}[]{@{}ll@{}}
\toprule\noalign{}
Parâmetro & Valor \\
\midrule\noalign{}
\endhead
\bottomrule\noalign{}
\endlastfoot
Velocidade média download & 85 Mbps \\
Domicílios com internet (dept.) & 87\% \\
Tecnologia predominante & fibra/\allowbreak cabo \\
Opção rural & Starlink disponível (\textasciitilde USD 44/\allowbreak mês) \\
\end{longtable}

\subsubsection{Mercado Imobiliário e Terra
Rural}

\textbf{Fonte:} INDERT /\allowbreak  Clasificados.com.py (2024)

\begin{longtable}[]{@{}ll@{}}
\toprule\noalign{}
Tipo & Referência \\
\midrule\noalign{}
\endhead
\bottomrule\noalign{}
\endlastfoot
Terra agrícola alta prod. (USD/\allowbreak ha) & 14,800 \\
Imóvel urbano (USD/\allowbreak m²) & 1,400 \\
Aluguel 2 quartos (USD/\allowbreak mês) & 600 \\
\end{longtable}

\begin{quote}
Valores de referência departamental. Localidades menores podem ter
preços 20--40\% abaixo da capital departamental. \#\#\# Saúde
\end{quote}

\textbf{Fonte:} MSPBS /\allowbreak  IPS Paraguay (2024-2026), consolidação
departamental e proxy local

\begin{longtable}[]{@{}ll@{}}
\toprule\noalign{}
Serviço & Disponibilidade \\
\midrule\noalign{}
\endhead
\bottomrule\noalign{}
\endlastfoot
USF /\allowbreak  Posto de Saúde & sim \\
Hospital Regional & não \\
IPS (seguro social) & sim \\
Farmácia & sim \\
Distância ao hospital de referência & \textasciitilde55 km (Ciudad del
Este) \\
\end{longtable}

\textbf{Principais estabelecimentos:} USF local; referencia hospitalar
em Ciudad del Este

\textbf{Observação para imigrantes:} Atendimento primario local ou em
raio curto; casos de maior complexidade seguem para Ciudad del Este.
Cobertura privada continua recomendada para especialidades e urgências
de maior complexidade.
\subsubsection{Indicadores Sociais}

\textbf{Fontes:} PNUD 2020, INE\footnote{Instituto Nacional de Estadística (Instituto Nacional de Estatística), órgão oficial de dados demográficos e censitários.} Censo 2022, INE\footnote{Instituto Nacional de Estadística (Instituto Nacional de Estatística), órgão oficial de dados demográficos e censitários.} EPHC 2023, Ministerio
Público 2024\\
\textbf{Nota:} Valores marcados como \emph{(dept.)} referem-se ao
departamento de Alto Paraná; valores \emph{(dist.)} são específicos
deste distrito. Dados marcados \emph{(est.)} são estimativas.

\begin{longtable}[]{@{}
  >{\raggedright\arraybackslash}p{(\linewidth - 4\tabcolsep) * \real{0.4231}}
  >{\raggedright\arraybackslash}p{(\linewidth - 4\tabcolsep) * \real{0.2692}}
  >{\raggedright\arraybackslash}p{(\linewidth - 4\tabcolsep) * \real{0.3077}}@{}}
\toprule\noalign{}
\begin{minipage}[b]{\linewidth}\raggedright
Indicador
\end{minipage} & \begin{minipage}[b]{\linewidth}\raggedright
Valor
\end{minipage} & \begin{minipage}[b]{\linewidth}\raggedright
Âmbito
\end{minipage} \\
\midrule\noalign{}
\endhead
\bottomrule\noalign{}
\endlastfoot
IDH (2020) & 0,752 & dept. \\
Ranking IDH nacional & 3/\allowbreak 18 & dept. \\
Esperança de vida & 75,1 anos & dept. \\
Escolaridade média & 7.6 anos & dist. \\
RNB per capita (USD PPA) & 13.500 & dept. \\
Pobreza monetária (\%) & 16,3\% & dept. \\
Pobreza extrema (\%) & 3,1\% & dept. \\
Índice de Gini & 0,428 & dept. \\
Acesso a água potável (\%) & 88,7\% & dept. \\
Acesso a saneamento (\%) & 87,2\% & dept. \\
Taxa de homicídios (est., /\allowbreak 100k hab) & \textasciitilde8--12 (estimativa,
acima da média nacional) & dept. \\
Índice de segurança & baixo & dept. \\
População (Censo 2022) & 5.035 hab. & dist. \\
Idade mediana & 27.0 anos & dist. \\
Taxa de fecundidade & N/\allowbreak D & dist. \\
\end{longtable}
\subsubsection{Combustível}

\textbf{Referência:} PETROPAR /\allowbreak  postos locais (2024) \textbf{Tipo de
localidade:} Interior

\begin{longtable}[]{@{}lll@{}}
\toprule\noalign{}
Combustível & USD/\allowbreak litro & Gs/\allowbreak litro (aprox.) \\
\midrule\noalign{}
\endhead
\bottomrule\noalign{}
\endlastfoot
Gasolina 93 oct & 0.97 & 7,178 \\
Gasolina 97 oct (premium) & 1.07 & 7,918 \\
Diesel & 0.90 & 6,660 \\
\end{longtable}

\begin{quote}
Preços podem variar ±5\% conforme posto e sazonalidade. Chaco e interior
remoto apresentam maior variação.
\end{quote}

\subsubsection{Cobertura Celular}

\textbf{Fonte:} CONATEL PY /\allowbreak  operadoras (2024)

\begin{longtable}[]{@{}ll@{}}
\toprule\noalign{}
Parâmetro & Valor \\
\midrule\noalign{}
\endhead
\bottomrule\noalign{}
\endlastfoot
Cobertura 4G população (dept.) & 95\% \\
Cobertura 4G área rural & 88\% \\
Melhor operadora & Tigo/\allowbreak Personal \\
Qualidade rural & boa \\
\end{longtable}

\begin{quote}
Para áreas rurais fora do núcleo urbano, recomenda-se chip Tigo como
principal e Personal como backup.
\end{quote}

\subsubsection{Internet}

\textbf{Fonte:} CONATEL /\allowbreak  Speedtest Ookla (2024)

\begin{longtable}[]{@{}ll@{}}
\toprule\noalign{}
Parâmetro & Valor \\
\midrule\noalign{}
\endhead
\bottomrule\noalign{}
\endlastfoot
Velocidade média download & 85 Mbps \\
Domicílios com internet (dept.) & 87\% \\
Tecnologia predominante & fibra/\allowbreak cabo \\
Opção rural & Starlink disponível (\textasciitilde USD 44/\allowbreak mês) \\
\end{longtable}

\subsubsection{Mercado Imobiliário e Terra
Rural}

\textbf{Fonte:} INDERT /\allowbreak  Clasificados.com.py (2024)

\begin{longtable}[]{@{}ll@{}}
\toprule\noalign{}
Tipo & Referência \\
\midrule\noalign{}
\endhead
\bottomrule\noalign{}
\endlastfoot
Terra agrícola alta prod. (USD/\allowbreak ha) & 14,800 \\
Imóvel urbano (USD/\allowbreak m²) & 1,400 \\
Aluguel 2 quartos (USD/\allowbreak mês) & 600 \\
\end{longtable}

\begin{quote}
Valores de referência departamental. Localidades menores podem ter
preços 20--40\% abaixo da capital departamental. \#\#\# Saúde
\end{quote}

\textbf{Fonte:} MSPBS /\allowbreak  IPS Paraguay (2024-2026), consolidação
departamental e proxy local

\begin{longtable}[]{@{}ll@{}}
\toprule\noalign{}
Serviço & Disponibilidade \\
\midrule\noalign{}
\endhead
\bottomrule\noalign{}
\endlastfoot
USF /\allowbreak  Posto de Saúde & sim \\
Hospital Regional & não \\
IPS (seguro social) & sim \\
Farmácia & sim \\
Distância ao hospital de referência & \textasciitilde55 km (Ciudad del
Este) \\
\end{longtable}

\textbf{Principais estabelecimentos:} USF local; referencia hospitalar
em Ciudad del Este

\textbf{Observação para imigrantes:} Atendimento primario local ou em
raio curto; casos de maior complexidade seguem para Ciudad del Este.
Cobertura privada continua recomendada para especialidades e urgências
de maior complexidade.
\newpage
\secaoDiagnostico{Hernandarias}{\mapaConteudo{-25.3833}{-54.6333}}{Hernandarias é a capital energética do Paraguai e um dos pontos mais
estratégicos da América do Sul. Oferece infraestrutura de serviços,
segurança e recursos (água/\allowbreak energia) de nível superior, garantindo uma
base operacional invulnerável a crises de suprimento básicas. Sua força
absoluta reside na supremacia energética e hídrica. Contudo, essa mesma
característica a coloca no centro de qualquer cenário de conflito
geopolítico regional. É o local ideal para quem busca integrar-se ao
topo da cadeia de recursos tecnológicos e segurança institucional,
exigindo do ocupante um posicionamento em zonas residenciais de elite
distantes dos perímetros táticos de Itaipu.

\subsubsection{Por que sim}

\begin{itemize}
\tightlist
\item
  Melhor segurança energética e hídrica do mundo.
\item
  Ambiente social seguro e infraestrutura institucional de primeiro
  nível.
\item
  Polo industrial dinâmico e solo de altíssima fertilidade.
\item
  Conectividade de elite via Supercarretera.
\end{itemize}

\subsubsection{Por que não}

\begin{itemize}
\tightlist
\item
  Alvo estratégico primário em qualquer conflito continental.
\item
  Risco de controle militar total do território em cenários de
  emergência nacional.
\item
  Visibilidade geopolítica extrema atrai riscos táticos especializados.
\end{itemize}

\subsubsection{Recomendação
Técnica}

Altamente indicado para estabelecimentos de suporte tecnológico,
residência de alto padrão ou base industrial. Requer autonomia tática
individual para períodos de controle militar de fluxos. O GSS de 8.0
reflete a força excepcional dos recursos e instituições em contraste com
a visibilidade estratégica radical da localidade. Referencia atual de
score: GSS 8.0 em 2026-03-05.

\subsubsection{}

\begin{itemize}
\tightlist
\item
  \begin{enumerate}
  \def\labelenumi{\arabic{enumi})}
  \tightlist
  \item
    Dados censitários validados (INE\footnote{Instituto Nacional de Estadística (Instituto Nacional de Estatística), órgão oficial de dados demográficos e censitários.} 2022).
  \end{enumerate}
\item
  \begin{enumerate}
  \def\labelenumi{\arabic{enumi})}
  \setcounter{enumi}{1}
  \tightlist
  \item
    Relatórios de segurança e integridade de barragens (Itaipu).
  \end{enumerate}
\item
  \begin{enumerate}
  \def\labelenumi{\arabic{enumi})}
  \setcounter{enumi}{2}
  \tightlist
  \item
    Mapa de ativos industriais e subestações táticas (ANDE\footnote{Administración Nacional de Electricidade (Administração Nacional de Eletricidade), companhia estatal de energia.}/\allowbreak MIC).
  \end{enumerate}
\item
  \begin{enumerate}
  \def\labelenumi{\arabic{enumi})}
  \setcounter{enumi}{3}
  \tightlist
  \item
    Registro de segurança pública e governança administrativa de elite.
  \end{enumerate}
\end{itemize}}
\begin{table}[ht!]
\small \begin{tabular}{p{3.0cm}p{0.8cm}p{6.0cm}} \toprule
\textbf{Critério} & \textbf{Nota} & \textbf{Justificativa} \\ \midrule
A - Ameaças Estratégicas & 6.5 & Hub estratégico mundial - Itaipu; alvo de altíssima visibilidade tática, mas sob proteção militar de elite \\
B - Risco Social & 7.0 & População regional gerenciada e qualificada; estabilidade social superior à média metropolitana \\
C - Riscos Naturais & 7.5 & Geologia muito estável e topografia serrana segura contra inundações sistêmicas \\
D - Autossuficiência & - D: 10.0 (Supremacia absoluta em recursos energéticos e hídricos; sede de ativos industriais de primeiro nível) & Supremacia absoluta em recursos energéticos e hídricos; sede de ativos industriais de primeiro nível \\
E - Institucional & 9.0 & Máxima segurança institucional e serviços de suporte de elite sustentados pela binacional \\
\midrule \textbf{GSS FINAL} & \textbf{8.0} & \textbf{Seguro (Localidade de elite para suporte tecnológico e residencial, oferecendo o máximo em recursos básicos do continente).} \\ \bottomrule \end{tabular}
\end{table}
\subsubsection{Vulnerabilidades e
Mitigação}\label{vuln-Hernandarias-vulnerabilidades-e-mitigauxe7uxe3o}

\begin{itemize}
\tightlist
\item
  \textbf{Alvo Estratégico Primário:} Proximidade com Itaipu atrai
  riscos militares globais (Mitigação: residência em áreas periféricas
  afastadas 15-20 km do eixo da barragem).
\item
  \textbf{Vulnerabilidade de Controle:} Risco de isolamento por controle
  militar rígido em crises (Mitigação: manter base com autonomia total
  de recursos e mapear rotas de terra para o interior de Alto Paraná).
\item
  \textbf{Dependência de Infraestrutura Crítica:} Qualquer falha na
  usina impacta imediatamente o distrito (Mitigação: instalação de
  sistemas solares fotovoltaicos individuais para backup).
\end{itemize}
\subsubsection{Dossiê de Campo}\label{dossie-Hernandarias-dossiuxea-de-campo}
\begin{description}[style=nextline,leftmargin=0.5cm,font=\bfseries]
\item[Coordenadas]  25°23'S 54°38'W.
\item[Topografia]  Relevo ondulado a acidentado, com solos de altíssima produtividade (Terra Roxa). Margeado ao leste pelo Rio Paraná e pelo reservatório de Itaipu. Altitude \textasciitilde 240m. Área de \textasciitilde 754 km². Geologicamente estável, risco sísmico nulo.
\item[Alvos Estratégicos]  Usina Hidrelétrica de Itaipu (Maior geradora de energia do mundo; infraestrutura estratégica crítica global); Represa de Acaray; Parque Industrial de Hernandarias; Sede da Reserva Biológica Itabó (margem sul). Alvo estratégico de altíssima visibilidade militar e geopolítica internacional.
\item[Fallout]  Ventos predominantes do Norte (NE) e Leste. Risco moderado a alto de fallout em caso de sabotagem infraestrutural direta nas barragens.
\item[População]  83.285 habitantes (Censo 2022 Final). Cidade em plena expansão industrial e residencial, servindo como polo tecnológico do departamento.
\item[Densidade]  \textasciitilde 110,5 hab/\allowbreak km² (Densidade moderada, oferecendo excelente isolamento social em áreas de condomínios fechados e zonas rurais produtivas).
\item[Vias de Acesso]  Conectividade asfáltica de elite via Ruta PY07 (Supercarretera). Localização privilegiada que permite rápido deslocamento para o Brasil e para o interior do Paraguai. Risco de ocupação militar total em cenários de crise energética nacional.
\item[Serviços]  Infraestrutura de saúde robusta (Hospital Distrital ampliado, IPS). Polo comercial, financeiro e industrial consolidado. Suporte institucional de elite garantido pelos investimentos da Itaipu Binacional.
\item[Custo de Vida]  Médio; economia dolarizada pela indústria e energia, com alta valorização imobiliária residencial.
\item[Preço da Terra]  US\$ 10.000-15.000/\allowbreak ha (terras mecanizadas); US\$ 80-200/\allowbreak m² (Urbano/\allowbreak Condomínios).
\item[Hidrologia]  Baixo risco de inundações sistêmicas devido à topografia elevada e controle hidráulico das barragens. Abundância hídrica superficial absoluta (reservatórios e arroios). Vulnerabilidade tática em caso de ruptura estrutural das barragens (embora improvável geologicamente).
\item[Clima]  Subtropical Úmido. Verões intensos amenizados pela proximidade com grandes massas de água.
\item[Energia]  Supremacia Energética Mundial. Localidade sede das principais hidrelétricas do país; prioridade absoluta de rede e redundância garantida.
\item[Água]  Abundância hídrica inesgotável via reservatórios e poços artesianos profundos (Aquífero Guarani).
\item[Qualidade do Solo]  Latossolo Vermelho de Elite; solo profundo e de fertilidade mundial. Excelente aptidão para policultivos tecnificados.
\item[Resources Locais]  Grande polo industrial de manufatura e tecnologia. Elevadíssimo potencial para autossuficiência logística e energética absoluta.
\item[Segurança]  Uma das cidades mais seguras e vigiadas da região metropolitana de CDE. Baixíssima criminalidade urbana comparada à capital vizinha. Forte presença institucional de forças militares e segurança interna de Itaipu. Coesão social baseada na estabilidade econômica provida pela binacional.
\item[Leis Local]  Município consolidado e polo de incentivos industriais. \\ \textbf{Livre de Restrição} de Fronteira. \\ \textit{Aquisição de terra rural proibida para estrangeiros limítrofes.}
\end{description}
\paragraph{Solo (SoilGrids 2.0, média ponderada 0--30
cm)}

\begin{longtable}[]{@{}ll@{}}
\toprule\noalign{}
Parâmetro & Valor \\
\midrule\noalign{}
\endhead
\bottomrule\noalign{}
\endlastfoot
pH (H$_2$O) & 5.15 \\
Carbono orgânico (SOC) & 20.68 g/\allowbreak kg \\
Argila & 54.29 \% \\
Areia & 14.12 \% \\
Silte (calc.) & 31.6 \% \\
Densidade aparente & 1.24 g/\allowbreak cm³ \\
\textbf{Aptidão agrícola} & \textbf{Média} \\
\end{longtable}

Fonte: ISRIC SoilGrids 2.0 via WCS (acesso em 2026-03-21). Coords:
-25.3833°, -54.6333°. Média ponderada camadas 0-5, 5-15, 15-30 cm.
\subsubsection{Indicadores Sociais}

\textbf{Fontes:} PNUD 2020, INE\footnote{Instituto Nacional de Estadística (Instituto Nacional de Estatística), órgão oficial de dados demográficos e censitários.} Censo 2022, INE\footnote{Instituto Nacional de Estadística (Instituto Nacional de Estatística), órgão oficial de dados demográficos e censitários.} EPHC 2023, Ministerio
Público 2024\\
\textbf{Nota:} Valores marcados como \emph{(dept.)} referem-se ao
departamento de Alto Paraná; valores \emph{(dist.)} são específicos
deste distrito. Dados marcados \emph{(est.)} são estimativas.

\begin{longtable}[]{@{}
  >{\raggedright\arraybackslash}p{(\linewidth - 4\tabcolsep) * \real{0.4231}}
  >{\raggedright\arraybackslash}p{(\linewidth - 4\tabcolsep) * \real{0.2692}}
  >{\raggedright\arraybackslash}p{(\linewidth - 4\tabcolsep) * \real{0.3077}}@{}}
\toprule\noalign{}
\begin{minipage}[b]{\linewidth}\raggedright
Indicador
\end{minipage} & \begin{minipage}[b]{\linewidth}\raggedright
Valor
\end{minipage} & \begin{minipage}[b]{\linewidth}\raggedright
Âmbito
\end{minipage} \\
\midrule\noalign{}
\endhead
\bottomrule\noalign{}
\endlastfoot
IDH (2020) & 0,752 & dept. \\
Ranking IDH nacional & 3/\allowbreak 18 & dept. \\
Esperança de vida & 75,1 anos & dept. \\
Escolaridade média & 10.1 anos & dist. \\
RNB per capita (USD PPA) & 13.500 & dept. \\
Pobreza monetária (\%) & 16,3\% & dept. \\
Pobreza extrema (\%) & 3,1\% & dept. \\
Índice de Gini & 0,428 & dept. \\
Acesso a água potável (\%) & 88,7\% & dept. \\
Acesso a saneamento (\%) & 87,2\% & dept. \\
Taxa de homicídios (est., /\allowbreak 100k hab) & \textasciitilde8--12 (estimativa,
acima da média nacional) & dept. \\
Índice de segurança & baixo & dept. \\
População (Censo 2022) & 83.285 hab. & dist. \\
Idade mediana & 28.0 anos & dist. \\
Taxa de fecundidade & N/\allowbreak D & dist. \\
\end{longtable}
\subsubsection{Combustível}

\textbf{Referência:} PETROPAR /\allowbreak  postos locais (2024) \textbf{Tipo de
localidade:} Interior

\begin{longtable}[]{@{}lll@{}}
\toprule\noalign{}
Combustível & USD/\allowbreak litro & Gs/\allowbreak litro (aprox.) \\
\midrule\noalign{}
\endhead
\bottomrule\noalign{}
\endlastfoot
Gasolina 93 oct & 0.97 & 7,178 \\
Gasolina 97 oct (premium) & 1.07 & 7,918 \\
Diesel & 0.90 & 6,660 \\
\end{longtable}

\begin{quote}
Preços podem variar ±5\% conforme posto e sazonalidade. Chaco e interior
remoto apresentam maior variação.
\end{quote}

\subsubsection{Cobertura Celular}

\textbf{Fonte:} CONATEL PY /\allowbreak  operadoras (2024)

\begin{longtable}[]{@{}ll@{}}
\toprule\noalign{}
Parâmetro & Valor \\
\midrule\noalign{}
\endhead
\bottomrule\noalign{}
\endlastfoot
Cobertura 4G população (dept.) & 95\% \\
Cobertura 4G área rural & 88\% \\
Melhor operadora & Tigo/\allowbreak Personal \\
Qualidade rural & boa \\
\end{longtable}

\begin{quote}
Para áreas rurais fora do núcleo urbano, recomenda-se chip Tigo como
principal e Personal como backup.
\end{quote}

\subsubsection{Internet}

\textbf{Fonte:} CONATEL /\allowbreak  Speedtest Ookla (2024)

\begin{longtable}[]{@{}ll@{}}
\toprule\noalign{}
Parâmetro & Valor \\
\midrule\noalign{}
\endhead
\bottomrule\noalign{}
\endlastfoot
Velocidade média download & 85 Mbps \\
Domicílios com internet (dept.) & 87\% \\
Tecnologia predominante & fibra/\allowbreak cabo \\
Opção rural & Starlink disponível (\textasciitilde USD 44/\allowbreak mês) \\
\end{longtable}

\subsubsection{Mercado Imobiliário e Terra
Rural}

\textbf{Fonte:} INDERT /\allowbreak  Clasificados.com.py (2024)

\begin{longtable}[]{@{}ll@{}}
\toprule\noalign{}
Tipo & Referência \\
\midrule\noalign{}
\endhead
\bottomrule\noalign{}
\endlastfoot
Terra agrícola alta prod. (USD/\allowbreak ha) & 14,800 \\
Imóvel urbano (USD/\allowbreak m²) & 1,400 \\
Aluguel 2 quartos (USD/\allowbreak mês) & 600 \\
\end{longtable}

\begin{quote}
Valores de referência departamental. Localidades menores podem ter
preços 20--40\% abaixo da capital departamental. \#\#\# Saúde
\end{quote}

\textbf{Fonte:} MSPBS /\allowbreak  IPS Paraguay (2024-2026), consolidação
departamental e proxy local

\begin{longtable}[]{@{}ll@{}}
\toprule\noalign{}
Serviço & Disponibilidade \\
\midrule\noalign{}
\endhead
\bottomrule\noalign{}
\endlastfoot
USF /\allowbreak  Posto de Saúde & sim \\
Hospital Regional & sim \\
IPS (seguro social) & sim \\
Farmácia & sim \\
Distância ao hospital de referência & \textasciitilde16 km (Ciudad del
Este) \\
\end{longtable}

\textbf{Principais estabelecimentos:} USF local; referencia hospitalar
em Ciudad del Este

\textbf{Observação para imigrantes:} Atendimento primario local ou em
raio curto; casos de maior complexidade seguem para Ciudad del Este.
Cobertura privada continua recomendada para especialidades e urgências
de maior complexidade.
\newpage
\secaoDiagnostico{Iruna}{\mapaConteudo{-25.8333}{-55.0833}}{Iruña é a localidade que oferece o maior potencial de isolamento
demográfico em Alto Paraná. Sua força absoluta reside na baixíssima
carga populacional combinada com solos de elite mundial (``terra
roxa''). Oferece um santuário de paz social e recursos naturais
abundantes, exigindo do ocupante autonomia logística e de saúde devido à
infraestrutura institucional mínima local e à vulnerabilidade sistêmica
da especialização agrícola exportadora.

\subsubsection{Por que sim}

\begin{itemize}
\tightlist
\item
  Menor densidade populacional do departamento (privacidade absoluta).
\item
  Solo de altíssima produtividade e clima favorável para policultivos.
\item
  Ambiente social extremamente seguro e ordeiro.
\item
  Sem restrições da Lei de Fronteira.
\end{itemize}

\subsubsection{Por que não}

\begin{itemize}
\tightlist
\item
  Infraestrutura de saúde local limitada a atendimentos primários.
\item
  Dependência econômica total de insumos e mercados internacionais.
\item
  Distância de centros comerciais e financeiros diversificados.
\end{itemize}

\subsubsection{Recomendação
Técnica}

Altamente indicado para estabelecimentos de autossuficiência tecnificada
que busquem o máximo isolamento demográfico em solo de alta performance.
Requer autonomia em logística básica e saúde. O GSS de 6.1 reflete a
segurança do isolamento em contraste com a limitação institucional
local. Referencia atual de score: GSS 6.1 em 2026-03-05.

\subsubsection{}

\begin{itemize}
\tightlist
\item
  \begin{enumerate}
  \def\labelenumi{\arabic{enumi})}
  \tightlist
  \item
    Dados censitários validados (INE\footnote{Instituto Nacional de Estadística (Instituto Nacional de Estatística), órgão oficial de dados demográficos e censitários.} 2022).
  \end{enumerate}
\item
  \begin{enumerate}
  \def\labelenumi{\arabic{enumi})}
  \setcounter{enumi}{1}
  \tightlist
  \item
    Relatórios de produtividade agrícola de precisão (MAG).
  \end{enumerate}
\item
  \begin{enumerate}
  \def\labelenumi{\arabic{enumi})}
  \setcounter{enumi}{2}
  \tightlist
  \item
    Mapa de solos e aptidão agrícola de Alto Paraná (MAG).
  \end{enumerate}
\item
  \begin{enumerate}
  \def\labelenumi{\arabic{enumi})}
  \setcounter{enumi}{3}
  \tightlist
  \item
    Registro de segurança pública departamental.
  \end{enumerate}
\end{itemize}}
\begin{table}[ht!]
\small \begin{tabular}{p{3.0cm}p{0.8cm}p{6.0cm}} \toprule
\textbf{Critério} & \textbf{Nota} & \textbf{Justificativa} \\ \midrule
A - Ameaças Estratégicas & 6.5 & Localização remota e discreta; altíssima invisibilidade tática e ausência de alvos estratégicos \\
B - Risco Social & 7.5 & Baixíssima densidade demográfica; o distrito mais vazio de Alto Paraná \\
C - Riscos Naturais & 7.0 & Geologia muito estável e baixo risco de desastres naturais graves \\
D - Autossuficiência & 7.5 & Recursos agrícolas de elite e abundância hídrica; limitado pela escala comercial local \\
E - Institucional & 2.0 & Nota penalizada no ranking nacional para refletir a dependência de insumos globais e infraestrutura local mínima - GSS 6.1 total \\
\midrule \textbf{GSS FINAL} & \textbf{6.1} & \textbf{Moderadamente Seguro (Ideal para quem busca o máximo de isolamento demográfico em terras de altíssima produtividade).} \\ \bottomrule \end{tabular}
\end{table}
\subsubsection{Vulnerabilidades e
Mitigação}\label{vuln-Iruna-vulnerabilidades-e-mitigauxe7uxe3o}

\begin{itemize}
\tightlist
\item
  \textbf{Isolamento de Saúde:} Distância significativa de hospitais
  cirúrgicos (Mitigação: manutenção de estoque médico básico e kit
  trauma local).
\item
  \textbf{Dependência de Cadeias Globais:} Economia focada
  exclusivamente em commodities (Mitigação: implementação de sistemas de
  policultivo de subsistência).
\item
  \textbf{Instabilidade Elétrica Rural:} Quedas ocasionais em temporais
  (Mitigação: instalação de sistemas solares fotovoltaicos redundantes).
\end{itemize}
\subsubsection{Dossiê de Campo}\label{dossie-Iruna-dossiuxea-de-campo}
\begin{description}[style=nextline,leftmargin=0.5cm,font=\bfseries]
\item[Coordenadas]  25°50'S 55°05'W.
\item[Topografia]  Relevo predominantemente plano a suavemente ondulado, ideal para agricultura mecanizada de precisão. Altitude \textasciitilde 250m. Área de \textasciitilde 531 km². Geologicamente muito estável, risco sísmico desprezível.
\item[Alvos Estratégicos]  Áreas de produção agrícola intensiva de soja e trigo; Silos de armazenamento de grãos. Localização remota no sul do departamento, oferecendo excelente invisibilidade tática e discrição estratégica. Ausência de alvos militares.
\item[Fallout]  Ventos predominantes do Norte (NE) e Leste. Baixo risco de fallout direto.
\item[População]  3.943 habitantes (Censo 2022 Final). O distrito menos povoado de Alto Paraná. Forte componente demográfico de imigrantes europeus e brasileiros.
\item[Densidade]  \textasciitilde 7,4 hab/\allowbreak km² (Densidade baixíssima, uma das menores do país na Região Oriental, oferecendo isolamento tático superior).
\item[Vias de Acesso]  Acesso via ramais pavimentados que conectam à Ruta PY06. Localização protegida por estar fora do radar dos grandes fluxos de trânsito nacional, funcionando como um enclave produtor discreto.
\item[Serviços]  Unidade de Saúde da Família (USF) local. Dependência total de Naranjal (25 km) ou Santa Rita para serviços hospitalares de alta complexidade e logística avançada.
\item[Custo de Vida]  Baixo; economia estritamente vinculada ao agronegócio exportador.
\item[Preço da Terra]  US\$ 10.000-15.000/\allowbreak ha (terras mecanizadas de elite - "Terra Roxa").
\item[Hidrologia]  Baixo risco de inundações sistêmicas catastróficas. Presença de pequenos arroios locais favoráveis à irrigação natural. Relevo favorável à drenagem.
\item[Clima]  Subtropical Úmido. Verões intensos e invernos com geadas que favorecem a produtividade do trigo.
\item[Energia]  Rede nacional (Itaipu); instável em áreas rurais remotas.
\item[Água]  Abundante via poços artesianos profundos e lençol freático produtivo (Sistema Aquífero Guarani).
\item[Qualidade do Solo]  Latossolo Vermelho de Elite; solo profundo e de altíssima fertilidade natural.
\item[Resources Locais]  Grande produção de soja, milho e trigo. Elevadíssimo potencial para autossuficiência calórica per capita devido à extrema baixa carga demográfica.
\item[Segurança]  Zona rural extremamente pacífica, ordeira e segura. Baixíssima criminalidade urbana. Coesão social fortíssima baseada na cultura cooperativista e laços familiares de produtores.
\item[Leis Local: Município consolidado. Livre de Restrição de Fronteira]  Localizado no interior do país, permitindo plena titularidade para brasileiros.
\end{description}
\paragraph{Solo (SoilGrids 2.0, média ponderada 0--30
cm)}

\begin{longtable}[]{@{}ll@{}}
\toprule\noalign{}
Parâmetro & Valor \\
\midrule\noalign{}
\endhead
\bottomrule\noalign{}
\endlastfoot
pH (H$_2$O) & 5.4 \\
Carbono orgânico (SOC) & 25.27 g/\allowbreak kg \\
Argila & 45.22 \% \\
Areia & 21.14 \% \\
Silte (calc.) & 33.6 \% \\
Densidade aparente & 1.19 g/\allowbreak cm³ \\
\textbf{Aptidão agrícola} & \textbf{Média} \\
\end{longtable}

Fonte: ISRIC SoilGrids 2.0 via WCS (acesso em 2026-03-21). Coords:
-25.8333°, -55.0833°. Média ponderada camadas 0-5, 5-15, 15-30 cm.

\begin{itemize}
\tightlist
\item
  \textbf{Homicidios/\allowbreak 100k:} \textasciitilde8,0 (Regional/\allowbreak Alto Paraná).
\end{itemize}

\subsection{10. Analise de Riscos}

\begin{itemize}
\tightlist
\item
  \textbf{Cenario curto prazo:} Manutenção da tranquilidade rural e
  rentabilidade das safras mecanizadas.
\item
  \textbf{Cenario medio prazo:} Impacto positivo da melhoria da
  infraestrutura vicinal ligando a Naranjal.
\end{itemize}

\subsection{11. Redacao Final}

\subsubsection{Diagnostico Integrado}

Iruña é a localidade que oferece o maior potencial de isolamento
demográfico em Alto Paraná. Sua força absoluta reside na baixíssima
carga populacional combinada com solos de elite mundial (``terra
roxa''). Oferece um santuário de paz social e recursos naturais
abundantes, exigindo do ocupante autonomia logística e de saúde devido à
infraestrutura institucional mínima local e à vulnerabilidade sistêmica
da especialização agrícola exportadora.

\subsubsection{Por que sim}

\begin{itemize}
\tightlist
\item
  Menor densidade populacional do departamento (privacidade absoluta).
\item
  Solo de altíssima produtividade e clima favorável para policultivos.
\item
  Ambiente social extremamente seguro e ordeiro.
\item
  Sem restrições da Lei de Fronteira.
\end{itemize}

\subsubsection{Por que nao}

\begin{itemize}
\tightlist
\item
  Infraestrutura de saúde local limitada a atendimentos primários.
\item
  Dependência econômica total de insumos e mercados internacionais.
\item
  Distância de centros comerciais e financeiros diversificados.
\end{itemize}

\subsubsection{Recomendacao Tecnica}

Altamente indicado para estabelecimentos de autossuficiência tecnificada
que busquem o máximo isolamento demográfico em solo de alta performance.
Requer autonomia em logística básica e saúde. O GSS de 6.1 reflete a
segurança do isolamento em contraste com a limitação institucional
local. Referencia atual de score: GSS 6.1 em 2026-03-05.

\subsection{13.}

\begin{itemize}
\tightlist
\item
  \begin{enumerate}
  \def\labelenumi{\arabic{enumi})}
  \tightlist
  \item
    Dados censitários validados (INE\footnote{Instituto Nacional de Estadística (Instituto Nacional de Estatística), órgão oficial de dados demográficos e censitários.} 2022).
  \end{enumerate}
\item
  \begin{enumerate}
  \def\labelenumi{\arabic{enumi})}
  \setcounter{enumi}{1}
  \tightlist
  \item
    Relatórios de produtividade agrícola de precisão (MAG).
  \end{enumerate}
\item
  \begin{enumerate}
  \def\labelenumi{\arabic{enumi})}
  \setcounter{enumi}{2}
  \tightlist
  \item
    Mapa de solos e aptidão agrícola de Alto Paraná (MAG).
  \end{enumerate}
\item
  \begin{enumerate}
  \def\labelenumi{\arabic{enumi})}
  \setcounter{enumi}{3}
  \tightlist
  \item
    Registro de segurança pública departamental.
  \end{enumerate}
\end{itemize}

\subsubsection{Combustível}

\textbf{Referência:} PETROPAR /\allowbreak  postos locais (2024) \textbf{Tipo de
localidade:} Interior

\begin{longtable}[]{@{}lll@{}}
\toprule\noalign{}
Combustível & USD/\allowbreak litro & Gs/\allowbreak litro (aprox.) \\
\midrule\noalign{}
\endhead
\bottomrule\noalign{}
\endlastfoot
Gasolina 93 oct & 0.97 & 7,178 \\
Gasolina 97 oct (premium) & 1.07 & 7,918 \\
Diesel & 0.90 & 6,660 \\
\end{longtable}

\begin{quote}
Preços podem variar ±5\% conforme posto e sazonalidade. Chaco e interior
remoto apresentam maior variação.
\end{quote}

\subsubsection{Cobertura Celular}

\textbf{Fonte:} CONATEL PY /\allowbreak  operadoras (2024)

\begin{longtable}[]{@{}ll@{}}
\toprule\noalign{}
Parâmetro & Valor \\
\midrule\noalign{}
\endhead
\bottomrule\noalign{}
\endlastfoot
Cobertura 4G população (dept.) & 95\% \\
Cobertura 4G área rural & 88\% \\
Melhor operadora & Tigo/\allowbreak Personal \\
Qualidade rural & boa \\
\end{longtable}

\begin{quote}
Para áreas rurais fora do núcleo urbano, recomenda-se chip Tigo como
principal e Personal como backup.
\end{quote}

\subsubsection{Internet}

\textbf{Fonte:} CONATEL /\allowbreak  Speedtest Ookla (2024)

\begin{longtable}[]{@{}ll@{}}
\toprule\noalign{}
Parâmetro & Valor \\
\midrule\noalign{}
\endhead
\bottomrule\noalign{}
\endlastfoot
Velocidade média download & 85 Mbps \\
Domicílios com internet (dept.) & 87\% \\
Tecnologia predominante & fibra/\allowbreak cabo \\
Opção rural & Starlink disponível (\textasciitilde USD 44/\allowbreak mês) \\
\end{longtable}

\subsubsection{Mercado Imobiliário e Terra
Rural}

\textbf{Fonte:} INDERT /\allowbreak  Clasificados.com.py (2024)

\begin{longtable}[]{@{}ll@{}}
\toprule\noalign{}
Tipo & Referência \\
\midrule\noalign{}
\endhead
\bottomrule\noalign{}
\endlastfoot
Terra agrícola alta prod. (USD/\allowbreak ha) & 14,800 \\
Imóvel urbano (USD/\allowbreak m²) & 1,400 \\
Aluguel 2 quartos (USD/\allowbreak mês) & 600 \\
\end{longtable}

\begin{quote}
Valores de referência departamental. Localidades menores podem ter
preços 20--40\% abaixo da capital departamental. \#\#\# Saúde
\end{quote}

\textbf{Fonte:} MSPBS /\allowbreak  IPS Paraguay (2024-2026), consolidação
departamental e proxy local

\begin{longtable}[]{@{}ll@{}}
\toprule\noalign{}
Serviço & Disponibilidade \\
\midrule\noalign{}
\endhead
\bottomrule\noalign{}
\endlastfoot
USF /\allowbreak  Posto de Saúde & sim \\
Hospital Regional & sim \\
IPS (seguro social) & sim \\
Farmácia & sim \\
Distância ao hospital de referência & \textasciitilde99 km (Ciudad del
Este) \\
\end{longtable}

\textbf{Principais estabelecimentos:} USF local; referencia hospitalar
em Ciudad del Este

\textbf{Observação para imigrantes:} Atendimento primario local ou em
raio curto; casos de maior complexidade seguem para Ciudad del Este.
Cobertura privada continua recomendada para especialidades e urgências
de maior complexidade.
\subsubsection{Indicadores Sociais}

\textbf{Fontes:} PNUD 2020, INE\footnote{Instituto Nacional de Estadística (Instituto Nacional de Estatística), órgão oficial de dados demográficos e censitários.} Censo 2022, INE\footnote{Instituto Nacional de Estadística (Instituto Nacional de Estatística), órgão oficial de dados demográficos e censitários.} EPHC 2023, Ministerio
Público 2024\\
\textbf{Nota:} Valores marcados como \emph{(dept.)} referem-se ao
departamento de Alto Paraná; valores \emph{(dist.)} são específicos
deste distrito. Dados marcados \emph{(est.)} são estimativas.

\begin{longtable}[]{@{}
  >{\raggedright\arraybackslash}p{(\linewidth - 4\tabcolsep) * \real{0.4231}}
  >{\raggedright\arraybackslash}p{(\linewidth - 4\tabcolsep) * \real{0.2692}}
  >{\raggedright\arraybackslash}p{(\linewidth - 4\tabcolsep) * \real{0.3077}}@{}}
\toprule\noalign{}
\begin{minipage}[b]{\linewidth}\raggedright
Indicador
\end{minipage} & \begin{minipage}[b]{\linewidth}\raggedright
Valor
\end{minipage} & \begin{minipage}[b]{\linewidth}\raggedright
Âmbito
\end{minipage} \\
\midrule\noalign{}
\endhead
\bottomrule\noalign{}
\endlastfoot
IDH (2020) & 0,752 & dept. \\
Ranking IDH nacional & 3/\allowbreak 18 & dept. \\
Esperança de vida & 75,1 anos & dept. \\
Escolaridade média & 8.3 anos & dist. \\
RNB per capita (USD PPA) & 13.500 & dept. \\
Pobreza monetária (\%) & 16,3\% & dept. \\
Pobreza extrema (\%) & 3,1\% & dept. \\
Índice de Gini & 0,428 & dept. \\
Acesso a água potável (\%) & 88,7\% & dept. \\
Acesso a saneamento (\%) & 87,2\% & dept. \\
Taxa de homicídios (est., /\allowbreak 100k hab) & \textasciitilde8--12 (estimativa,
acima da média nacional) & dept. \\
Índice de segurança & baixo & dept. \\
População (Censo 2022) & 3.943 hab. & dist. \\
Idade mediana & 29.0 anos & dist. \\
Taxa de fecundidade & N/\allowbreak D & dist. \\
\end{longtable}
\subsubsection{Combustível}

\textbf{Referência:} PETROPAR /\allowbreak  postos locais (2024) \textbf{Tipo de
localidade:} Interior

\begin{longtable}[]{@{}lll@{}}
\toprule\noalign{}
Combustível & USD/\allowbreak litro & Gs/\allowbreak litro (aprox.) \\
\midrule\noalign{}
\endhead
\bottomrule\noalign{}
\endlastfoot
Gasolina 93 oct & 0.97 & 7,178 \\
Gasolina 97 oct (premium) & 1.07 & 7,918 \\
Diesel & 0.90 & 6,660 \\
\end{longtable}

\begin{quote}
Preços podem variar ±5\% conforme posto e sazonalidade. Chaco e interior
remoto apresentam maior variação.
\end{quote}

\subsubsection{Cobertura Celular}

\textbf{Fonte:} CONATEL PY /\allowbreak  operadoras (2024)

\begin{longtable}[]{@{}ll@{}}
\toprule\noalign{}
Parâmetro & Valor \\
\midrule\noalign{}
\endhead
\bottomrule\noalign{}
\endlastfoot
Cobertura 4G população (dept.) & 95\% \\
Cobertura 4G área rural & 88\% \\
Melhor operadora & Tigo/\allowbreak Personal \\
Qualidade rural & boa \\
\end{longtable}

\begin{quote}
Para áreas rurais fora do núcleo urbano, recomenda-se chip Tigo como
principal e Personal como backup.
\end{quote}

\subsubsection{Internet}

\textbf{Fonte:} CONATEL /\allowbreak  Speedtest Ookla (2024)

\begin{longtable}[]{@{}ll@{}}
\toprule\noalign{}
Parâmetro & Valor \\
\midrule\noalign{}
\endhead
\bottomrule\noalign{}
\endlastfoot
Velocidade média download & 85 Mbps \\
Domicílios com internet (dept.) & 87\% \\
Tecnologia predominante & fibra/\allowbreak cabo \\
Opção rural & Starlink disponível (\textasciitilde USD 44/\allowbreak mês) \\
\end{longtable}

\subsubsection{Mercado Imobiliário e Terra
Rural}

\textbf{Fonte:} INDERT /\allowbreak  Clasificados.com.py (2024)

\begin{longtable}[]{@{}ll@{}}
\toprule\noalign{}
Tipo & Referência \\
\midrule\noalign{}
\endhead
\bottomrule\noalign{}
\endlastfoot
Terra agrícola alta prod. (USD/\allowbreak ha) & 14,800 \\
Imóvel urbano (USD/\allowbreak m²) & 1,400 \\
Aluguel 2 quartos (USD/\allowbreak mês) & 600 \\
\end{longtable}

\begin{quote}
Valores de referência departamental. Localidades menores podem ter
preços 20--40\% abaixo da capital departamental. \#\#\# Saúde
\end{quote}

\textbf{Fonte:} MSPBS /\allowbreak  IPS Paraguay (2024-2026), consolidação
departamental e proxy local

\begin{longtable}[]{@{}ll@{}}
\toprule\noalign{}
Serviço & Disponibilidade \\
\midrule\noalign{}
\endhead
\bottomrule\noalign{}
\endlastfoot
USF /\allowbreak  Posto de Saúde & sim \\
Hospital Regional & sim \\
IPS (seguro social) & sim \\
Farmácia & sim \\
Distância ao hospital de referência & \textasciitilde99 km (Ciudad del
Este) \\
\end{longtable}

\textbf{Principais estabelecimentos:} USF local; referencia hospitalar
em Ciudad del Este

\textbf{Observação para imigrantes:} Atendimento primario local ou em
raio curto; casos de maior complexidade seguem para Ciudad del Este.
Cobertura privada continua recomendada para especialidades e urgências
de maior complexidade.
\newpage
\secaoDiagnostico{Itakyry}{\mapaConteudo{-24.9333}{-54.9833}}{Itakyry é um dos maiores refúgios demográficos e produtivos de Alto
Paraná. Sua força absoluta reside na vastidão territorial com baixa
densidade populacional, combinada com solos de elite (``terra
vermelha'') e segurança social superior. Oferece um santuário rural
discreto e altamente produtivo, sendo o local ideal para projetos de
autossuficiência de larga escala que priorizam o isolamento tático e o
anonimato, exigindo do ocupante atenção à gestão de conflitos agrários e
autonomia em saúde.

\subsubsection{Por que sim}

\begin{itemize}
\tightlist
\item
  Vasta área territorial com baixa carga populacional (privacidade
  total).
\item
  Solo de altíssima produtividade e abundância de águas puras.
\item
  Localização estratégica fora do radar de grandes crises demográficas.
\item
  Sem restrições da Lei de Fronteira.
\end{itemize}

\subsubsection{Por que não}

\begin{itemize}
\tightlist
\item
  Infraestrutura de saúde local limitada a atendimentos primários.
\item
  Risco de tensões agrárias em áreas de mata ou fronteira agrícola.
\item
  Pequena escala comercial exige deslocamento para centros maiores.
\end{itemize}

\subsubsection{Recomendação
Técnica}

Altamente indicado para estabelecimentos de autossuficiência de
médio/\allowbreak grande porte que busquem isolamento demográfico massivo em solo de
alta performance. Requer autonomia em logística básica, saúde e
segurança patrimonial. O GSS de 6.4 reflete a segurança do isolamento em
contraste com a limitação institucional local. Referencia atual de
score: GSS 6.4 em 2026-03-05.

\subsubsection{}

\begin{itemize}
\tightlist
\item
  \begin{enumerate}
  \def\labelenumi{\arabic{enumi})}
  \tightlist
  \item
    Dados censitários validados (INE\footnote{Instituto Nacional de Estadística (Instituto Nacional de Estatística), órgão oficial de dados demográficos e censitários.} 2022).
  \end{enumerate}
\item
  \begin{enumerate}
  \def\labelenumi{\arabic{enumi})}
  \setcounter{enumi}{1}
  \tightlist
  \item
    Levantamento territorial oficial do distrito (Portal
    Geoestadístico).
  \end{enumerate}
\item
  \begin{enumerate}
  \def\labelenumi{\arabic{enumi})}
  \setcounter{enumi}{2}
  \tightlist
  \item
    Mapa de solos e bacias hidrográficas de Alto Paraná (MAG).
  \end{enumerate}
\item
  \begin{enumerate}
  \def\labelenumi{\arabic{enumi})}
  \setcounter{enumi}{3}
  \tightlist
  \item
    Registro de segurança pública departamental.
  \end{enumerate}
\end{itemize}}
\begin{table}[ht!]
\small \begin{tabular}{p{3.0cm}p{0.8cm}p{6.0cm}} \toprule
\textbf{Critério} & \textbf{Nota} & \textbf{Justificativa} \\ \midrule
A - Ameaças Estratégicas & 6.0 & Localização discreta e protegida de rotas principais; excelente invisibilidade tática \\
B - Risco Social & 7.0 & População reduzida e densidade baixíssima em território vasto; isolamento social superior \\
C - Riscos Naturais & 7.0 & Geologia muito estável e baixo risco de desastres naturais graves \\
D - Autossuficiência & 7.5 & Excelentes recursos: solo de elite, abundância hídrica e potencial de proteína \\
E - Institucional & 4.5 & Município estável seguro, mas com infraestrutura institucional limitada e histórico de tensões agrárias \\
\midrule \textbf{GSS FINAL} & \textbf{6.4} & \textbf{Moderadamente Seguro (Ideal para quem busca o máximo de isolamento demográfico em terras de altíssima produtividade).} \\ \bottomrule \end{tabular}
\end{table}
\subsubsection{Vulnerabilidades e
Mitigação}\label{vuln-Itakyry-vulnerabilidades-e-mitigauxe7uxe3o}

\begin{itemize}
\tightlist
\item
  \textbf{Isolamento de Saúde:} Distância significativa de hospitais
  cirúrgicos (Mitigação: manutenção de estoque médico básico e kit
  trauma local).
\item
  \textbf{Conflitos de Terra:} Risco de tensões em zonas de fronteira
  agrícola (Mitigação: aquisição de propriedades com título perfeito e
  integração social local).
\item
  \textbf{Dependência Logística:} Necessidade de deslocamento para
  cidades maiores para insumos técnicos (Mitigação: estoque de
  suprimentos para 60 dias).
\end{itemize}
\subsubsection{Dossiê de Campo}\label{dossie-Itakyry-dossiuxea-de-campo}
\begin{description}[style=nextline,leftmargin=0.5cm,font=\bfseries]
\item[Coordenadas]  24°56'S 54°59'W.
\item[Topografia]  Relevo predominantemente plano a suavemente ondulado, com extensas chapadas mecanizáveis. Altitude \textasciitilde 250m. Área vasta de \textasciitilde 1.964 km². Geologicamente muito estável, risco sísmico desprezível.
\item[Alvos Estratégicos]  Áreas de produção agrícola intensiva de grãos; Terminais de silos de grande porte; Ponto de conexão rodoviária entre Alto Paraná e Canindeyú. Ausência de alvos militares diretos, oferecendo excelente invisibilidade estratégica em território vasto.
\item[Fallout]  Ventos predominantes do Norte (NE) e Leste. Baixo risco de fallout direto.
\item[População]  27.340 habitantes (Censo 2022 Final). Perfil demográfico majoritariamente rural (70\%) com forte presença de comunidades tradicionais e agricultores mecanizados.
\item[Densidade]  \textasciitilde 13,9 hab/\allowbreak km² (Densidade baixíssima, oferecendo isolamento tático superior e privacidade estratégica absoluta fora dos núcleos povoados).
\item[Vias de Acesso]  Acesso via ramais da Ruta PY07. Conectividade asfáltica eficiente para o Norte e Sul, mas com isolamento relativo em relação ao eixo central PY02. Malha interna de estradas de terra que exigem veículos robustos em períodos chuvosos (média 1.700 mm/\allowbreak ano).
\item[Serviços]  Centro de Saúde local. Dependência de Hernandarias (70 km) ou Ciudad del Este para serviços hospitalares de alta complexidade e logística avançada.
\item[Custo de Vida]  Muito baixo; economia estritamente vinculada ao agronegócio e agricultura de subsistência.
\item[Preço da Terra]  US\$ 8.000-12.000/\allowbreak ha (terras mecanizadas de elite); US\$ 2.500-4.000/\allowbreak ha (áreas de campo natural/\allowbreak pecuária).
\item[Hidrologia]  Baixo risco de inundações sistêmicas catastróficas devido ao relevo favorável e gestão hídrica agrícola. Presença de arroios tributários da bacia do Rio Paraná.
\item[Clima]  Subtropical Úmido. Verões quentes; invernos com geadas ocasionais.
\item[Energia]  Rede nacional (Itaipu); instável em áreas rurais remotas durante picos de demanda agrícola.
\item[Água]  Abundante via poços artesianos profundos e sistemas de captação local; excelente qualidade hídrica subterrânea.
\item[Qualidade do Solo]  Latossolo Vermelho de Elite; solo profundo, estruturado e de altíssima fertilidade natural. Excelente aptidão para soja, milho e trigo.
\item[Resources Locais]  Grande produção de soja, milho e gado de corte. Elevadíssimo potencial para autossuficiência alimentar básica em nível de propriedade devido à vasta área produtiva.
\item[Segurança]  Zona rural pacífica e estável. Baixíssima criminalidade urbana. Presença de comunidades indígenas e histórico de conflitos agrários pontuais em áreas de mata que exigem vigilância tática e respeito territorial. Estabilidade social baseada na cultura produtora.
\item[Leis Local: Município histórico e consolidado no agronegócio. Livre de Restrição de Fronteira]  Localizado no interior do país, permitindo plena titularidade para brasileiros.
\end{description}
\paragraph{Solo (SoilGrids 2.0, média ponderada 0--30
cm)}

\begin{longtable}[]{@{}ll@{}}
\toprule\noalign{}
Parâmetro & Valor \\
\midrule\noalign{}
\endhead
\bottomrule\noalign{}
\endlastfoot
pH (H$_2$O) & 5.35 \\
Carbono orgânico (SOC) & 24.38 g/\allowbreak kg \\
Argila & 43.0 \% \\
Areia & 30.75 \% \\
Silte (calc.) & 26.2 \% \\
Densidade aparente & 1.23 g/\allowbreak cm³ \\
\textbf{Aptidão agrícola} & \textbf{Média} \\
\end{longtable}

Fonte: ISRIC SoilGrids 2.0 via WCS (acesso em 2026-03-21). Coords:
-24.9333°, -54.9833°. Média ponderada camadas 0-5, 5-15, 15-30 cm.
\subsubsection{Indicadores Sociais}

\textbf{Fontes:} PNUD 2020, INE\footnote{Instituto Nacional de Estadística (Instituto Nacional de Estatística), órgão oficial de dados demográficos e censitários.} Censo 2022, INE\footnote{Instituto Nacional de Estadística (Instituto Nacional de Estatística), órgão oficial de dados demográficos e censitários.} EPHC 2023, Ministerio
Público 2024\\
\textbf{Nota:} Valores marcados como \emph{(dept.)} referem-se ao
departamento de Alto Paraná; valores \emph{(dist.)} são específicos
deste distrito. Dados marcados \emph{(est.)} são estimativas.

\begin{longtable}[]{@{}
  >{\raggedright\arraybackslash}p{(\linewidth - 4\tabcolsep) * \real{0.4231}}
  >{\raggedright\arraybackslash}p{(\linewidth - 4\tabcolsep) * \real{0.2692}}
  >{\raggedright\arraybackslash}p{(\linewidth - 4\tabcolsep) * \real{0.3077}}@{}}
\toprule\noalign{}
\begin{minipage}[b]{\linewidth}\raggedright
Indicador
\end{minipage} & \begin{minipage}[b]{\linewidth}\raggedright
Valor
\end{minipage} & \begin{minipage}[b]{\linewidth}\raggedright
Âmbito
\end{minipage} \\
\midrule\noalign{}
\endhead
\bottomrule\noalign{}
\endlastfoot
IDH (2020) & 0,752 & dept. \\
Ranking IDH nacional & 3/\allowbreak 18 & dept. \\
Esperança de vida & 75,1 anos & dept. \\
Escolaridade média & 7.2 anos & dist. \\
RNB per capita (USD PPA) & 13.500 & dept. \\
Pobreza monetária (\%) & 16,3\% & dept. \\
Pobreza extrema (\%) & 3,1\% & dept. \\
Índice de Gini & 0,428 & dept. \\
Acesso a água potável (\%) & 88,7\% & dept. \\
Acesso a saneamento (\%) & 87,2\% & dept. \\
Taxa de homicídios (est., /\allowbreak 100k hab) & \textasciitilde8--12 (estimativa,
acima da média nacional) & dept. \\
Índice de segurança & baixo & dept. \\
População (Censo 2022) & 2.734 hab. & dist. \\
Idade mediana & 23.0 anos & dist. \\
Taxa de fecundidade & N/\allowbreak D & dist. \\
\end{longtable}
\subsubsection{Combustível}

\textbf{Referência:} PETROPAR /\allowbreak  postos locais (2024) \textbf{Tipo de
localidade:} Interior

\begin{longtable}[]{@{}lll@{}}
\toprule\noalign{}
Combustível & USD/\allowbreak litro & Gs/\allowbreak litro (aprox.) \\
\midrule\noalign{}
\endhead
\bottomrule\noalign{}
\endlastfoot
Gasolina 93 oct & 0.97 & 7,178 \\
Gasolina 97 oct (premium) & 1.07 & 7,918 \\
Diesel & 0.90 & 6,660 \\
\end{longtable}

\begin{quote}
Preços podem variar ±5\% conforme posto e sazonalidade. Chaco e interior
remoto apresentam maior variação.
\end{quote}

\subsubsection{Cobertura Celular}

\textbf{Fonte:} CONATEL PY /\allowbreak  operadoras (2024)

\begin{longtable}[]{@{}ll@{}}
\toprule\noalign{}
Parâmetro & Valor \\
\midrule\noalign{}
\endhead
\bottomrule\noalign{}
\endlastfoot
Cobertura 4G população (dept.) & 95\% \\
Cobertura 4G área rural & 88\% \\
Melhor operadora & Tigo/\allowbreak Personal \\
Qualidade rural & boa \\
\end{longtable}

\begin{quote}
Para áreas rurais fora do núcleo urbano, recomenda-se chip Tigo como
principal e Personal como backup.
\end{quote}

\subsubsection{Internet}

\textbf{Fonte:} CONATEL /\allowbreak  Speedtest Ookla (2024)

\begin{longtable}[]{@{}ll@{}}
\toprule\noalign{}
Parâmetro & Valor \\
\midrule\noalign{}
\endhead
\bottomrule\noalign{}
\endlastfoot
Velocidade média download & 85 Mbps \\
Domicílios com internet (dept.) & 87\% \\
Tecnologia predominante & fibra/\allowbreak cabo \\
Opção rural & Starlink disponível (\textasciitilde USD 44/\allowbreak mês) \\
\end{longtable}

\subsubsection{Mercado Imobiliário e Terra
Rural}

\textbf{Fonte:} INDERT /\allowbreak  Clasificados.com.py (2024)

\begin{longtable}[]{@{}ll@{}}
\toprule\noalign{}
Tipo & Referência \\
\midrule\noalign{}
\endhead
\bottomrule\noalign{}
\endlastfoot
Terra agrícola alta prod. (USD/\allowbreak ha) & 14,800 \\
Imóvel urbano (USD/\allowbreak m²) & 1,400 \\
Aluguel 2 quartos (USD/\allowbreak mês) & 600 \\
\end{longtable}

\begin{quote}
Valores de referência departamental. Localidades menores podem ter
preços 20--40\% abaixo da capital departamental. \#\#\# Saúde
\end{quote}

\textbf{Fonte:} MSPBS /\allowbreak  IPS Paraguay (2024-2026), consolidação
departamental e proxy local

\begin{longtable}[]{@{}ll@{}}
\toprule\noalign{}
Serviço & Disponibilidade \\
\midrule\noalign{}
\endhead
\bottomrule\noalign{}
\endlastfoot
USF /\allowbreak  Posto de Saúde & sim \\
Hospital Regional & sim \\
IPS (seguro social) & sim \\
Farmácia & sim \\
Distância ao hospital de referência & \textasciitilde343 km (Ciudad del
Este) \\
\end{longtable}

\textbf{Principais estabelecimentos:} USF local; referencia hospitalar
em Ciudad del Este

\textbf{Observação para imigrantes:} Atendimento primario local ou em
raio curto; casos de maior complexidade seguem para Ciudad del Este.
Cobertura privada continua recomendada para especialidades e urgências
de maior complexidade.
\newpage
\secaoDiagnostico{Juan Emilio OLeary}{\mapaConteudo{-25.4166}{-55.3833}}{Juan Emilio O'Leary oferece uma combinação sólida de alta produtividade
agrícola e excepcional segurança hídrica devido à proximidade com o
reservatório de Yguazú. Sua força reside na qualidade de seus solos e na
facilidade logística provida pela Ruta PY02 duplicada. Contudo, essa
mesma conectividade a torna vulnerável em cenários de instabilidade
nacional, exigindo do ocupante um posicionamento rural periférico que
aproveite a riqueza de recursos sem sofrer a exposição do eixo
rodoviário principal.

\subsubsection{Por que sim}

\begin{itemize}
\tightlist
\item
  Solo de altíssima produtividade (``terra roxa'') e clima favorável.
\item
  Segurança hídrica massiva via lago e poços profundos.
\item
  Conectividade rápida com os polos de serviços de Campo 9 e CDE.
\item
  Sem restrições da Lei de Fronteira.
\end{itemize}

\subsubsection{Por que não}

\begin{itemize}
\tightlist
\item
  Alvo logístico rodoviário de alta visibilidade estratégica em crises.
\item
  Infraestrutura de saúde local limitada a atendimentos primários.
\item
  Risco de fluxos migratórios massivos vindos da capital ou da
  fronteira.
\end{itemize}

\subsubsection{Recomendação
Técnica}

Altamente indicado para estabelecimentos de autossuficiência que busquem
abundância de recursos produtivos. Requer planejamento de recuo físico
da rodovia e autonomia em saúde básica. O GSS de 6.7 reflete a força dos
recursos naturais equilibrada pela vulnerabilidade tática rodoviária.
Referencia atual de score: GSS 6.7 em 2026-03-05.

\subsubsection{}

\begin{itemize}
\tightlist
\item
  \begin{enumerate}
  \def\labelenumi{\arabic{enumi})}
  \tightlist
  \item
    Dados censitários validados (INE\footnote{Instituto Nacional de Estadística (Instituto Nacional de Estatística), órgão oficial de dados demográficos e censitários.} 2022).
  \end{enumerate}
\item
  \begin{enumerate}
  \def\labelenumi{\arabic{enumi})}
  \setcounter{enumi}{1}
  \tightlist
  \item
    Relatórios de produção agrícola mecanizada (MAG).
  \end{enumerate}
\item
  \begin{enumerate}
  \def\labelenumi{\arabic{enumi})}
  \setcounter{enumi}{2}
  \tightlist
  \item
    Mapa de conectividade do corredor PY02 (MOPC 2026).
  \end{enumerate}
\item
  \begin{enumerate}
  \def\labelenumi{\arabic{enumi})}
  \setcounter{enumi}{3}
  \tightlist
  \item
    Registro de segurança pública departamental.
  \end{enumerate}
\end{itemize}}
\begin{table}[ht!]
\small \begin{tabular}{p{3.0cm}p{0.8cm}p{6.0cm}} \toprule
\textbf{Critério} & \textbf{Nota} & \textbf{Justificativa} \\ \midrule
A - Ameaças Estratégicas & 5.0 & Posicionamento estratégico na PY02; alvo logístico secundário de visibilidade moderada \\
B - Risco Social & 6.5 & População controlada e densidade moderada; boa infraestrutura rodoviária duplicada \\
C - Riscos Naturais & 7.0 & Geologia muito estável e baixo risco de desastres naturais graves \\
D - Autossuficiência & 8.0 & Excelentes recursos: solo de elite, abundância hídrica do lago e produção de grãos \\
E - Institucional & 6.5 & Ambiente institucional seguro e pacífico, mas sob pressão logística rodoviária \\
\midrule \textbf{GSS FINAL} & \textbf{6.7} & \textbf{Moderadamente Seguro (Ideal para quem busca integração agrícola com abundância de recursos hídricos).} \\ \bottomrule \end{tabular}
\end{table}
\subsubsection{Vulnerabilidades e
Mitigação}\label{vuln-Juan_Emilio_OLeary-vulnerabilidades-e-mitigauxe7uxe3o}

\begin{itemize}
\tightlist
\item
  \textbf{Exposição Rodoviária:} Localização direta na PY02 atrai fluxos
  de trânsito intenso em crises (Mitigação: residência rural afastada
  3-5 km do eixo asfáltico principal).
\item
  \textbf{Dependência de Serviços Externos:} Saúde especializada em
  cidades vizinhas (Mitigação: manutenção de estoque médico básico e
  treinamento de primeiros socorros).
\item
  \textbf{Visibilidade Tática:} Ponto de passagem entre polos
  industriais (Mitigação: mapeamento de rotas rurais alternativas para
  evitar o asfalto em caso de bloqueios).
\end{itemize}
\subsubsection{Dossiê de Campo}\label{dossie-Juan_Emilio_OLeary-dossiuxea-de-campo}
\begin{description}[style=nextline,leftmargin=0.5cm,font=\bfseries]
\item[Coordenadas]  25°25'S 55°23'W.
\item[Topografia]  Relevo predominantemente plano a suavemente ondulado, margeado pelo reservatório de Yguazú. Altitude \textasciitilde 220m. Área de \textasciitilde 226 km². Geologicamente muito estável, risco sísmico desprezível.
\item[Alvos Estratégicos]  Ponto de passagem e serviços na Ruta PY02 duplicada; Proximidade estratégica com a bacia do reservatório de Yguazú (potencial hídrico massivo). Ausência de alvos militares diretos, mas alta visibilidade logística nacional.
\item[Fallout]  Ventos predominantes do Norte (NE) e Leste. Baixo risco de fallout direto.
\item[População]  16.092 habitantes (Censo 2022 Final). Perfil demográfico urbano-rural equilibrado ao longo do eixo asfáltico.
\item[Densidade]  \textasciitilde 71,2 hab/\allowbreak km² (Densidade moderada, com maior concentração comercial sobre a rodovia e áreas rurais dispersas).
\item[Vias de Acesso]  Localização estratégica sobre a Ruta PY02 duplicada. Conectividade asfáltica de elite para Ciudad del Este (aprox. 80 km) e Assunção. Risco elevado de tornar-se ponto de controle e rota de fuga massiva em crises nacionais.
\item[Serviços]  Unidade de Saúde da Família (USF) local. Dependência de J. Eulogio Estigarribia (Campo 9) ou Ciudad del Este para serviços hospitalares de alta complexidade e logística comercial especializada.
\item[Custo de Vida]  Baixo; economia baseada no agronegócio e serviços rodoviários.
\item[Preço da Terra]  US\$ 6.000-10.000/\allowbreak ha (áreas rurais mecanizadas); US\$ 3.000-5.000/\allowbreak ha (áreas próximas ao reservatório com potencial turístico).
\item[Hidrologia]  Baixo risco de inundações sistêmicas catastróficas. Áreas ribeirinhas baixas próximas ao Lago Yguazú podem sofrer variações de nível em anos de precipitação extrema (média 1.700 mm/\allowbreak ano). Drenagem natural favorável.
\item[Clima]  Subtropical Úmido. Verões quentes; invernos com geadas ocasionais.
\item[Energia]  Rede nacional (Itaipu); boa estabilidade por estar no eixo principal de transmissão leste-oeste.
\item[Água]  Abundante via poços artesianos e proximidade com o reservatório de Yguazú; lençol freático produtivo.
\item[Qualidade do Solo]  Latossolo Vermelho (Terra Roxa); profundo e de altíssima fertilidade natural. Excelente aptidão para soja, milho, trigo e horticultura.
\item[Resources Locais]  Grande produção de grãos (soja/\allowbreak milho) e pecuária. Elevado potencial para autossuficiência alimentar básica em nível de propriedade.
\item[Segurança]  Zona rural pacífica e estável. Baixíssima criminalidade urbana violenta. Estabilidade social sustentada pela cultura agrícola produtiva. Exposição moderada a fluxos de trânsito rodoviário.
\item[Leis Local: Município consolidado. Livre de Restrição de Fronteira]  Localizado no interior do país, permitindo plena titularidade para brasileiros.
\end{description}
\paragraph{Solo (SoilGrids 2.0, média ponderada 0--30
cm)}

\begin{longtable}[]{@{}ll@{}}
\toprule\noalign{}
Parâmetro & Valor \\
\midrule\noalign{}
\endhead
\bottomrule\noalign{}
\endlastfoot
pH (H$_2$O) & 5.2 \\
Carbono orgânico (SOC) & 19.98 g/\allowbreak kg \\
Argila & 32.18 \% \\
Areia & 41.68 \% \\
Silte (calc.) & 26.1 \% \\
Densidade aparente & 1.32 g/\allowbreak cm³ \\
\textbf{Aptidão agrícola} & \textbf{Média} \\
\end{longtable}

Fonte: ISRIC SoilGrids 2.0 via WCS (acesso em 2026-03-21). Coords:
-25.4167°, -55.3833°. Média ponderada camadas 0-5, 5-15, 15-30 cm.
\subsubsection{Indicadores Sociais}

\textbf{Fontes:} PNUD 2020, INE\footnote{Instituto Nacional de Estadística (Instituto Nacional de Estatística), órgão oficial de dados demográficos e censitários.} Censo 2022, INE\footnote{Instituto Nacional de Estadística (Instituto Nacional de Estatística), órgão oficial de dados demográficos e censitários.} EPHC 2023, Ministerio
Público 2024\\
\textbf{Nota:} Valores marcados como \emph{(dept.)} referem-se ao
departamento de Alto Paraná; valores \emph{(dist.)} são específicos
deste distrito. Dados marcados \emph{(est.)} são estimativas.

\begin{longtable}[]{@{}
  >{\raggedright\arraybackslash}p{(\linewidth - 4\tabcolsep) * \real{0.4231}}
  >{\raggedright\arraybackslash}p{(\linewidth - 4\tabcolsep) * \real{0.2692}}
  >{\raggedright\arraybackslash}p{(\linewidth - 4\tabcolsep) * \real{0.3077}}@{}}
\toprule\noalign{}
\begin{minipage}[b]{\linewidth}\raggedright
Indicador
\end{minipage} & \begin{minipage}[b]{\linewidth}\raggedright
Valor
\end{minipage} & \begin{minipage}[b]{\linewidth}\raggedright
Âmbito
\end{minipage} \\
\midrule\noalign{}
\endhead
\bottomrule\noalign{}
\endlastfoot
IDH (2020) & 0,752 & dept. \\
Ranking IDH nacional & 3/\allowbreak 18 & dept. \\
Esperança de vida & 75,1 anos & dept. \\
Escolaridade média & 9,6 anos & dept. \\
RNB per capita (USD PPA) & 13.500 & dept. \\
Pobreza monetária (\%) & 16,3\% & dept. \\
Pobreza extrema (\%) & 3,1\% & dept. \\
Índice de Gini & 0,428 & dept. \\
Acesso a água potável (\%) & 88,7\% & dept. \\
Acesso a saneamento (\%) & 87,2\% & dept. \\
Taxa de homicídios (est., /\allowbreak 100k hab) & \textasciitilde8--12 (estimativa,
acima da média nacional) & dept. \\
Índice de segurança & baixo & dept. \\
População (Censo 2022) & 16.092 habitantes & dist. \\
Idade mediana & N/\allowbreak D & dist. \\
Taxa de fecundidade & N/\allowbreak D & dist. \\
\end{longtable}
\subsubsection{Combustível}

\textbf{Referência:} PETROPAR /\allowbreak  postos locais (2024) \textbf{Tipo de
localidade:} Interior

\begin{longtable}[]{@{}lll@{}}
\toprule\noalign{}
Combustível & USD/\allowbreak litro & Gs/\allowbreak litro (aprox.) \\
\midrule\noalign{}
\endhead
\bottomrule\noalign{}
\endlastfoot
Gasolina 93 oct & 0.97 & 7,178 \\
Gasolina 97 oct (premium) & 1.07 & 7,918 \\
Diesel & 0.90 & 6,660 \\
\end{longtable}

\begin{quote}
Preços podem variar ±5\% conforme posto e sazonalidade. Chaco e interior
remoto apresentam maior variação.
\end{quote}

\subsubsection{Cobertura Celular}

\textbf{Fonte:} CONATEL PY /\allowbreak  operadoras (2024)

\begin{longtable}[]{@{}ll@{}}
\toprule\noalign{}
Parâmetro & Valor \\
\midrule\noalign{}
\endhead
\bottomrule\noalign{}
\endlastfoot
Cobertura 4G população (dept.) & 95\% \\
Cobertura 4G área rural & 88\% \\
Melhor operadora & Tigo/\allowbreak Personal \\
Qualidade rural & boa \\
\end{longtable}

\begin{quote}
Para áreas rurais fora do núcleo urbano, recomenda-se chip Tigo como
principal e Personal como backup.
\end{quote}

\subsubsection{Internet}

\textbf{Fonte:} CONATEL /\allowbreak  Speedtest Ookla (2024)

\begin{longtable}[]{@{}ll@{}}
\toprule\noalign{}
Parâmetro & Valor \\
\midrule\noalign{}
\endhead
\bottomrule\noalign{}
\endlastfoot
Velocidade média download & 85 Mbps \\
Domicílios com internet (dept.) & 87\% \\
Tecnologia predominante & fibra/\allowbreak cabo \\
Opção rural & Starlink disponível (\textasciitilde USD 44/\allowbreak mês) \\
\end{longtable}

\subsubsection{Mercado Imobiliário e Terra
Rural}

\textbf{Fonte:} INDERT /\allowbreak  Clasificados.com.py (2024)

\begin{longtable}[]{@{}ll@{}}
\toprule\noalign{}
Tipo & Referência \\
\midrule\noalign{}
\endhead
\bottomrule\noalign{}
\endlastfoot
Terra agrícola alta prod. (USD/\allowbreak ha) & 14,800 \\
Imóvel urbano (USD/\allowbreak m²) & 1,400 \\
Aluguel 2 quartos (USD/\allowbreak mês) & 600 \\
\end{longtable}

\begin{quote}
Valores de referência departamental. Localidades menores podem ter
preços 20--40\% abaixo da capital departamental. \#\#\# Saúde
\end{quote}

\textbf{Fonte:} MSPBS /\allowbreak  IPS Paraguay (2024-2026), consolidação
departamental e proxy local

\begin{longtable}[]{@{}ll@{}}
\toprule\noalign{}
Serviço & Disponibilidade \\
\midrule\noalign{}
\endhead
\bottomrule\noalign{}
\endlastfoot
USF /\allowbreak  Posto de Saúde & sim \\
Hospital Regional & sim \\
IPS (seguro social) & sim \\
Farmácia & sim \\
Distância ao hospital de referência & \textasciitilde63 km (Ciudad del
Este) \\
\end{longtable}

\textbf{Principais estabelecimentos:} USF local; referencia hospitalar
em Ciudad del Este

\textbf{Observação para imigrantes:} Atendimento primario local ou em
raio curto; casos de maior complexidade seguem para Ciudad del Este.
Cobertura privada continua recomendada para especialidades e urgências
de maior complexidade.
\newpage
\secaoDiagnostico{Juan Leon Mallorquin}{\mapaConteudo{-25.4000}{-55.2666}}{Juan León Mallorquín é uma localidade que equilibra a alta produtividade
agrícola com uma densidade populacional favorável ao isolamento
relativo. Sua força reside na excepcional qualidade de seus solos
(``terra roxa'') e na infraestrutura rodoviária de elite provida pela
Ruta PY02 duplicada. Contudo, essa mesma conectividade a torna
vulnerável em cenários de instabilidade nacional, exigindo do ocupante
um posicionamento rural periférico que aproveite os recursos produtivos
sem sofrer a exposição do eixo rodoviário principal.

\subsubsection{Por que sim}

\begin{itemize}
\tightlist
\item
  Solo de altíssima produtividade e clima favorável para policultivos.
\item
  Abundância de águas puras subterrâneas e superficiais.
\item
  Conectividade rápida com os polos de serviços de CDE e Campo 9.
\item
  Sem restrições da Lei de Fronteira.
\end{itemize}

\subsubsection{Por que não}

\begin{itemize}
\tightlist
\item
  Alvo logístico rodoviário de visibilidade estratégica em crises.
\item
  Infraestrutura de saúde local limitada a atendimentos primários.
\item
  Risco de fluxos migratórios massivos vindos da capital ou da
  fronteira.
\end{itemize}

\subsubsection{Recomendação
Técnica}

Altamente indicado para estabelecimentos de autossuficiência que busquem
abundância de recursos em ambiente seguro. Requer planejamento de recuo
físico da rodovia e autonomia em saúde básica. O GSS de 6.3 reflete a
força dos recursos naturais equilibrada pela vulnerabilidade tática
rodoviária. Referencia atual de score: GSS 6.3 em 2026-03-05.

\subsubsection{}

\begin{itemize}
\tightlist
\item
  \begin{enumerate}
  \def\labelenumi{\arabic{enumi})}
  \tightlist
  \item
    Dados censitários validados (INE\footnote{Instituto Nacional de Estadística (Instituto Nacional de Estatística), órgão oficial de dados demográficos e censitários.} 2022).
  \end{enumerate}
\item
  \begin{enumerate}
  \def\labelenumi{\arabic{enumi})}
  \setcounter{enumi}{1}
  \tightlist
  \item
    Relatórios de produção agrícola mecanizada (MAG).
  \end{enumerate}
\item
  \begin{enumerate}
  \def\labelenumi{\arabic{enumi})}
  \setcounter{enumi}{2}
  \tightlist
  \item
    Mapa de conectividade do corredor PY02 (MOPC 2026).
  \end{enumerate}
\item
  \begin{enumerate}
  \def\labelenumi{\arabic{enumi})}
  \setcounter{enumi}{3}
  \tightlist
  \item
    Registro de segurança pública departamental.
  \end{enumerate}
\end{itemize}}
\begin{table}[ht!]
\small \begin{tabular}{p{3.0cm}p{0.8cm}p{6.0cm}} \toprule
\textbf{Critério} & \textbf{Nota} & \textbf{Justificativa} \\ \midrule
A - Ameaças Estratégicas & 5.0 & Posicionamento estratégico na PY02; alvo logístico secundário de visibilidade moderada \\
B - Risco Social & 6.5 & População controlada e densidade baixa; boa infraestrutura rodoviária duplicada \\
C - Riscos Naturais & 7.0 & Geologia muito estável e baixo risco de desastres naturais graves \\
D - Autossuficiência & 7.5 & Excelentes recursos: solo de elite, abundância hídrica e produção de grãos \\
E - Institucional & 6.0 & Ambiente institucional estável, mas dependente de serviços de cidades maiores \\
\midrule \textbf{GSS FINAL} & \textbf{6.3} & \textbf{Moderadamente Seguro (Ideal para quem busca integração agrícola com facilidade logística).} \\ \bottomrule \end{tabular}
\end{table}
\subsubsection{Vulnerabilidades e
Mitigação}\label{vuln-Juan_Leon_Mallorquin-vulnerabilidades-e-mitigauxe7uxe3o}

\begin{itemize}
\tightlist
\item
  \textbf{Exposição Rodoviária:} Localização direta na PY02 atrai fluxos
  de trânsito (Mitigação: residência rural afastada 3-5 km do eixo
  asfáltico).
\item
  \textbf{Dependência de Serviços Externos:} Saúde especializada em
  cidades distantes (Mitigação: manutenção de estoque médico básico e
  kit trauma local).
\item
  \textbf{Visibilidade Tática:} Ponto de parada logístico (Mitigação:
  mapeamento de rotas rurais alternativas para evitar o asfalto em caso
  de bloqueios).
\end{itemize}
\subsubsection{Dossiê de Campo}\label{dossie-Juan_Leon_Mallorquin-dossiuxea-de-campo}
\begin{description}[style=nextline,leftmargin=0.5cm,font=\bfseries]
\item[Coordenadas]  25°24'S 55°16'W.
\item[Topografia]  Relevo predominantemente plano a suavemente ondulado ("lomadas"). Altitude \textasciitilde 230m. Área de \textasciitilde 689 km². Geologicamente muito estável, risco sísmico desprezível.
\item[Alvos Estratégicos]  Ponto de passagem e logística na Ruta PY02 duplicada; Silos de armazenamento de grãos; Proximidade estratégica com o Rio Monday. Ausência de alvos militares de grande porte, mas alta visibilidade logística nacional.
\item[Fallout]  Ventos predominantes do Norte (NE) e Leste. Baixo risco de fallout direto.
\item[População]  16.144 habitantes (Censo 2022 Final). População com forte tradição no agronegócio e comércio local.
\item[Densidade]  \textasciitilde 23,4 hab/\allowbreak km² (Densidade baixa, oferecendo isolamento relativo em áreas de colônias rurais).
\item[Vias de Acesso]  Localização privilegiada sobre a Ruta PY02 duplicada. Conectividade asfáltica de elite para Ciudad del Este (65 km) e Assunção. Risco elevado de tornar-se ponto de controle e destino de fluxos migratórios em crises nacionais.
\item[Serviços]  Unidade de Saúde da Família (USF) e clínica do IPS local. Dependência de Ciudad del Este ou Minga Guazú para saúde de alta complexidade e logística especializada.
\item[Custo de Vida]  Baixo; economia baseada na exportação de grãos e serviços rodoviários.
\item[Preço da Terra]  US\$ 8.000-12.000/\allowbreak ha (terras mecanizadas de elite - "Terra Roxa").
\item[Hidrologia]  Cruzado por arroios perenes tributários da bacia do Rio Paraná. Baixo risco de inundações sistêmicas devido à topografia favorável. Drenagem natural de alta performance.
\item[Clima]  Subtropical Úmido. Verões quentes; invernos frescos com geadas ocasionais.
\item[Energia]  Rede nacional (Itaipu); boa estabilidade devido ao posicionamento em eixo de transmissão principal.
\item[Água]  Abundante via poços artesianos profundos e lençol freático produtivo (Sistema Aquífero Guarani).
\item[Qualidade do Solo]  Latossolo Vermelho de Elite; profundo, fértil e estruturado. Excelente aptidão para soja, milho e trigo.
\item[Resources Locais]  Grande produção de grãos e pecuária de corte. Elevado potencial para autossuficiência alimentar básica em nível de propriedade.
\item[Segurança]  Zona rural pacífica e estável. Baixíssima criminalidade urbana violenta. Estabilidade social sustentada pela cultura do agronegócio familiar. Exposição moderada a fluxos de trânsito rodoviário nacional.
\item[Leis Local: Município consolidado (antigo Ka'arendy). Livre de Restrição de Fronteira]  Localizado no interior do país, permitindo plena titularidade para brasileiros.
\end{description}
\paragraph{Solo (SoilGrids 2.0, média ponderada 0--30
cm)}

\begin{longtable}[]{@{}ll@{}}
\toprule\noalign{}
Parâmetro & Valor \\
\midrule\noalign{}
\endhead
\bottomrule\noalign{}
\endlastfoot
pH (H$_2$O) & 5.15 \\
Carbono orgânico (SOC) & 18.05 g/\allowbreak kg \\
Argila & 31.1 \% \\
Areia & 43.07 \% \\
Silte (calc.) & 25.8 \% \\
Densidade aparente & 1.32 g/\allowbreak cm³ \\
\textbf{Aptidão agrícola} & \textbf{Média} \\
\end{longtable}

Fonte: ISRIC SoilGrids 2.0 via WCS (acesso em 2026-03-21). Coords:
-25.4°, -55.2667°. Média ponderada camadas 0-5, 5-15, 15-30 cm.
\subsubsection{Indicadores Sociais}

\textbf{Fontes:} PNUD 2020, INE\footnote{Instituto Nacional de Estadística (Instituto Nacional de Estatística), órgão oficial de dados demográficos e censitários.} Censo 2022, INE\footnote{Instituto Nacional de Estadística (Instituto Nacional de Estatística), órgão oficial de dados demográficos e censitários.} EPHC 2023, Ministerio
Público 2024\\
\textbf{Nota:} Valores marcados como \emph{(dept.)} referem-se ao
departamento de Alto Paraná; valores \emph{(dist.)} são específicos
deste distrito. Dados marcados \emph{(est.)} são estimativas.

\begin{longtable}[]{@{}
  >{\raggedright\arraybackslash}p{(\linewidth - 4\tabcolsep) * \real{0.4231}}
  >{\raggedright\arraybackslash}p{(\linewidth - 4\tabcolsep) * \real{0.2692}}
  >{\raggedright\arraybackslash}p{(\linewidth - 4\tabcolsep) * \real{0.3077}}@{}}
\toprule\noalign{}
\begin{minipage}[b]{\linewidth}\raggedright
Indicador
\end{minipage} & \begin{minipage}[b]{\linewidth}\raggedright
Valor
\end{minipage} & \begin{minipage}[b]{\linewidth}\raggedright
Âmbito
\end{minipage} \\
\midrule\noalign{}
\endhead
\bottomrule\noalign{}
\endlastfoot
IDH (2020) & 0,752 & dept. \\
Ranking IDH nacional & 3/\allowbreak 18 & dept. \\
Esperança de vida & 75,1 anos & dept. \\
Escolaridade média & 9,6 anos & dept. \\
RNB per capita (USD PPA) & 13.500 & dept. \\
Pobreza monetária (\%) & 16,3\% & dept. \\
Pobreza extrema (\%) & 3,1\% & dept. \\
Índice de Gini & 0,428 & dept. \\
Acesso a água potável (\%) & 88,7\% & dept. \\
Acesso a saneamento (\%) & 87,2\% & dept. \\
Taxa de homicídios (est., /\allowbreak 100k hab) & \textasciitilde8--12 (estimativa,
acima da média nacional) & dept. \\
Índice de segurança & baixo & dept. \\
População (Censo 2022) & 16.144 hab. & dist. \\
Idade mediana & 28.0 anos & dist. \\
Taxa de fecundidade & N/\allowbreak D & dist. \\
\end{longtable}
\subsubsection{Combustível}

\textbf{Referência:} PETROPAR /\allowbreak  postos locais (2024) \textbf{Tipo de
localidade:} Interior

\begin{longtable}[]{@{}lll@{}}
\toprule\noalign{}
Combustível & USD/\allowbreak litro & Gs/\allowbreak litro (aprox.) \\
\midrule\noalign{}
\endhead
\bottomrule\noalign{}
\endlastfoot
Gasolina 93 oct & 0.97 & 7,178 \\
Gasolina 97 oct (premium) & 1.07 & 7,918 \\
Diesel & 0.90 & 6,660 \\
\end{longtable}

\begin{quote}
Preços podem variar ±5\% conforme posto e sazonalidade. Chaco e interior
remoto apresentam maior variação.
\end{quote}

\subsubsection{Cobertura Celular}

\textbf{Fonte:} CONATEL PY /\allowbreak  operadoras (2024)

\begin{longtable}[]{@{}ll@{}}
\toprule\noalign{}
Parâmetro & Valor \\
\midrule\noalign{}
\endhead
\bottomrule\noalign{}
\endlastfoot
Cobertura 4G população (dept.) & 95\% \\
Cobertura 4G área rural & 88\% \\
Melhor operadora & Tigo/\allowbreak Personal \\
Qualidade rural & boa \\
\end{longtable}

\begin{quote}
Para áreas rurais fora do núcleo urbano, recomenda-se chip Tigo como
principal e Personal como backup.
\end{quote}

\subsubsection{Internet}

\textbf{Fonte:} CONATEL /\allowbreak  Speedtest Ookla (2024)

\begin{longtable}[]{@{}ll@{}}
\toprule\noalign{}
Parâmetro & Valor \\
\midrule\noalign{}
\endhead
\bottomrule\noalign{}
\endlastfoot
Velocidade média download & 85 Mbps \\
Domicílios com internet (dept.) & 87\% \\
Tecnologia predominante & fibra/\allowbreak cabo \\
Opção rural & Starlink disponível (\textasciitilde USD 44/\allowbreak mês) \\
\end{longtable}

\subsubsection{Mercado Imobiliário e Terra
Rural}

\textbf{Fonte:} INDERT /\allowbreak  Clasificados.com.py (2024)

\begin{longtable}[]{@{}ll@{}}
\toprule\noalign{}
Tipo & Referência \\
\midrule\noalign{}
\endhead
\bottomrule\noalign{}
\endlastfoot
Terra agrícola alta prod. (USD/\allowbreak ha) & 14,800 \\
Imóvel urbano (USD/\allowbreak m²) & 1,400 \\
Aluguel 2 quartos (USD/\allowbreak mês) & 600 \\
\end{longtable}

\begin{quote}
Valores de referência departamental. Localidades menores podem ter
preços 20--40\% abaixo da capital departamental. \#\#\# Saúde
\end{quote}

\textbf{Fonte:} MSPBS /\allowbreak  IPS Paraguay (2024-2026), consolidação
departamental e proxy local

\begin{longtable}[]{@{}ll@{}}
\toprule\noalign{}
Serviço & Disponibilidade \\
\midrule\noalign{}
\endhead
\bottomrule\noalign{}
\endlastfoot
USF /\allowbreak  Posto de Saúde & sim \\
Hospital Regional & não \\
IPS (seguro social) & sim \\
Farmácia & sim \\
Distância ao hospital de referência & \textasciitilde81 km (Ciudad del
Este) \\
\end{longtable}

\textbf{Principais estabelecimentos:} USF local; referencia hospitalar
em Ciudad del Este

\textbf{Observação para imigrantes:} Atendimento primario local ou em
raio curto; casos de maior complexidade seguem para Ciudad del Este.
Cobertura privada continua recomendada para especialidades e urgências
de maior complexidade.
\newpage
\secaoDiagnostico{Los Cedrales}{\mapaConteudo{-25.6666}{-54.7166}}{Los Cedrales (Alto Parana) apresenta perfil operacional consolidado para
comparacao territorial, com leitura conjunta de infraestrutura, risco
hidroclimatico e capacidade de continuidade de servicos essenciais.

\subsubsection{Por que sim}

\begin{itemize}
\tightlist
\item
  Base institucional de referencia disponivel e rastreavel.
\item
  Estrutura comparativa (A-E/\allowbreak GSS) consistente para decisao relativa.
\item
  Plano de mitigacao acionavel ja definido no dossie.
\end{itemize}

\subsubsection{Por que não}

\begin{itemize}
\tightlist
\item
  Serie oficial distrital granular ainda pode apresentar lacunas
  temporais.
\item
  Sensibilidade do score exige monitoramento continuo de risco e
  conectividade.
\item
  Decisao final depende de validacao de campo e janela temporal recente.
\end{itemize}

\subsubsection{Pre-condicoes Minimas de
Decisao}

\begin{itemize}
\tightlist
\item
  Atualizar indicadores criticos em ciclo trimestral.
\item
  Confirmar trafegabilidade e contingencia hidrica/\allowbreak energetica local.
\item
  Validar checklist de resiliencia domiciliar/\allowbreak logistica antes de decisao
  final.
\end{itemize}

\subsubsection{Recomendação
Técnica}

Uso recomendado como candidato em analise multicriterio, condicionado ao
cumprimento das pre-condicoes acima. Referencia atual de score: GSS 7.0
em 2026-03-05; risco predominante: moderado.

\subsubsection{}

\begin{itemize}
\tightlist
\item
  \begin{enumerate}
  \def\labelenumi{\arabic{enumi})}
  \tightlist
  \item
    Delimitacao territorial validada por georreferencia institucional
    (Fonte: acesso: 2026-03-05).
  \end{enumerate}
\item
  \begin{enumerate}
  \def\labelenumi{\arabic{enumi})}
  \setcounter{enumi}{1}
  \tightlist
  \item
    População municipal/\allowbreak distrital referenciada por base censitaria
    nacional (Fonte: acesso: 2026-03-05).
  \end{enumerate}
\item
  \begin{enumerate}
  \def\labelenumi{\arabic{enumi})}
  \setcounter{enumi}{2}
  \tightlist
  \item
    Comparabilidade intra-departamental apoiada em indicadores
    distritais oficiais (Fonte: acesso: 2026-03-05).
  \end{enumerate}
\item
  \begin{enumerate}
  \def\labelenumi{\arabic{enumi})}
  \setcounter{enumi}{3}
  \tightlist
  \item
    Estado macro de conectividade viaria ancorado em informacoes de
    rodovias (Fonte: acesso: 2026-03-05).
  \end{enumerate}
\item
  \begin{enumerate}
  \def\labelenumi{\arabic{enumi})}
  \setcounter{enumi}{4}
  \tightlist
  \item
    Exposicao hidrometeorologica monitorada por avisos oficiais (Fonte:
    acesso: 2026-03-05).
  \end{enumerate}
\item
  \begin{enumerate}
  \def\labelenumi{\arabic{enumi})}
  \setcounter{enumi}{5}
  \tightlist
  \item
    Historico de contexto meteorologico apoiado em publicacoes tecnicas
    (Fonte: acesso: 2026-03-05).
  \end{enumerate}
\item
  \begin{enumerate}
  \def\labelenumi{\arabic{enumi})}
  \setcounter{enumi}{6}
  \tightlist
  \item
    Capacidade institucional de resposta a eventos apoiada em acoes da
    SEN\footnote{Secretaría de Emergencia Nacional (Secretaria de Emergência Nacional), órgão governamental de resposta a desastres.} (Fonte: acesso: 2026-03-05).
  \end{enumerate}
\item
  \begin{enumerate}
  \def\labelenumi{\arabic{enumi})}
  \setcounter{enumi}{7}
  \tightlist
  \item
    Infraestrutura energetica analisada via operador nacional (Fonte:
    acesso: 2026-03-05).
  \end{enumerate}
\item
  \begin{enumerate}
  \def\labelenumi{\arabic{enumi})}
  \setcounter{enumi}{8}
  \tightlist
  \item
    Referencias de obras e logistica nacional verificadas no MOPC
    (Fonte: acesso: 2026-03-05).
  \end{enumerate}
\item
  \begin{enumerate}
  \def\labelenumi{\arabic{enumi})}
  \setcounter{enumi}{9}
  \tightlist
  \item
    Contexto sociopolitico institucional referenciado por autoridade
    eleitoral (Fonte: acesso: 2026-03-05).
  \end{enumerate}
\item
  \begin{enumerate}
  \def\labelenumi{\arabic{enumi})}
  \setcounter{enumi}{10}
  \tightlist
  \item
    Camada cartografica de apoio eleitoral/\allowbreak territorial consultada
    (Fonte: acesso: 2026-03-05).
  \end{enumerate}
\item
  \begin{enumerate}
  \def\labelenumi{\arabic{enumi})}
  \setcounter{enumi}{11}
  \tightlist
  \item
    Dados publicos complementares para verificacao cruzada (Fonte:
    acesso: 2026-03-05).
  \end{enumerate}
\end{itemize}}
\begin{table}[ht!]
\small \begin{tabular}{p{3.0cm}p{0.8cm}p{6.0cm}} \toprule
\textbf{Critério} & \textbf{Nota} & \textbf{Justificativa} \\ \midrule
A - Ameaças Estratégicas & 7.1 & - \\
B - Risco Social & 7.5 & - \\
C - Riscos Naturais & 5.8 & - \\
D - Autossuficiência & 6.5 & - \\
E - Institucional & 7.6 & - \\
\midrule \textbf{GSS FINAL} & \textbf{7.0} & \textbf{Seguro (bom potencial com preparacao).} \\ \bottomrule \end{tabular}
\end{table}
\subsubsection{Vulnerabilidades e
Mitigação}\label{vuln-Los_Cedrales-vulnerabilidades-e-mitigauxe7uxe3o}

\begin{itemize}
\tightlist
\item
  Exposicao hidrometeorologica moderada (45-64/\allowbreak 100).
\item
  Conectividade viaria intermediaria (50-69/\allowbreak 100).
\item
  Estoque de contingencia para 14 dias (agua/\allowbreak alimentos/\allowbreak combustivel),
  priorizado quando risco hidro=52/\allowbreak 100.
\item
  Duplicar rotas/\allowbreak logistica quando conectividade viaria=54/\allowbreak 100 for
  \textless70; manter rota secundaria mapeada e testada trimestralmente.
\item
  Plano de energia resiliente com autonomia minima de 72h
  (geracao/\allowbreak backup), revisao semestral.
\end{itemize}
\subsubsection{Dossiê de Campo}\label{dossie-Los_Cedrales-dossiuxea-de-campo}
\begin{description}[style=nextline,leftmargin=0.5cm,font=\bfseries]
\item[Coordenadas]  25°40'S 54°43'W (geocodificado via Nominatim).
\end{description}
\paragraph{Solo (SoilGrids 2.0, média ponderada 0--30
cm)}

\begin{longtable}[]{@{}ll@{}}
\toprule\noalign{}
Parâmetro & Valor \\
\midrule\noalign{}
\endhead
\bottomrule\noalign{}
\endlastfoot
pH (H$_2$O) & 5.5 \\
Carbono orgânico (SOC) & 24.42 g/\allowbreak kg \\
Argila & 49.77 \% \\
Areia & 17.83 \% \\
Silte (calc.) & 32.4 \% \\
Densidade aparente & 1.23 g/\allowbreak cm³ \\
\textbf{Aptidão agrícola} & \textbf{Média} \\
\end{longtable}

Fonte: ISRIC SoilGrids 2.0 via WCS (acesso em 2026-03-21). Coords:
-25.6667°, -54.7167°. Média ponderada camadas 0-5, 5-15, 15-30 cm.
\subsubsection{Indicadores Sociais}

\textbf{Fontes:} PNUD 2020, INE\footnote{Instituto Nacional de Estadística (Instituto Nacional de Estatística), órgão oficial de dados demográficos e censitários.} Censo 2022, INE\footnote{Instituto Nacional de Estadística (Instituto Nacional de Estatística), órgão oficial de dados demográficos e censitários.} EPHC 2023, Ministerio
Público 2024\\
\textbf{Nota:} Valores marcados como \emph{(dept.)} referem-se ao
departamento de Alto Paraná; valores \emph{(dist.)} são específicos
deste distrito. Dados marcados \emph{(est.)} são estimativas.

\begin{longtable}[]{@{}
  >{\raggedright\arraybackslash}p{(\linewidth - 4\tabcolsep) * \real{0.4231}}
  >{\raggedright\arraybackslash}p{(\linewidth - 4\tabcolsep) * \real{0.2692}}
  >{\raggedright\arraybackslash}p{(\linewidth - 4\tabcolsep) * \real{0.3077}}@{}}
\toprule\noalign{}
\begin{minipage}[b]{\linewidth}\raggedright
Indicador
\end{minipage} & \begin{minipage}[b]{\linewidth}\raggedright
Valor
\end{minipage} & \begin{minipage}[b]{\linewidth}\raggedright
Âmbito
\end{minipage} \\
\midrule\noalign{}
\endhead
\bottomrule\noalign{}
\endlastfoot
IDH (2020) & 0,752 & dept. \\
Ranking IDH nacional & 3/\allowbreak 18 & dept. \\
Esperança de vida & 75,1 anos & dept. \\
Escolaridade média & 8.1 anos & dist. \\
RNB per capita (USD PPA) & 13.500 & dept. \\
Pobreza monetária (\%) & 16,3\% & dept. \\
Pobreza extrema (\%) & 3,1\% & dept. \\
Índice de Gini & 0,428 & dept. \\
Acesso a água potável (\%) & 88,7\% & dept. \\
Acesso a saneamento (\%) & 87,2\% & dept. \\
Taxa de homicídios (est., /\allowbreak 100k hab) & \textasciitilde8--12 (estimativa,
acima da média nacional) & dept. \\
Índice de segurança & baixo & dept. \\
População (Censo 2022) & 7.258 hab. & dist. \\
Idade mediana & 28.0 anos & dist. \\
Taxa de fecundidade & N/\allowbreak D & dist. \\
\end{longtable}
\subsubsection{Combustível}

\textbf{Referência:} PETROPAR /\allowbreak  postos locais (2024) \textbf{Tipo de
localidade:} Interior

\begin{longtable}[]{@{}lll@{}}
\toprule\noalign{}
Combustível & USD/\allowbreak litro & Gs/\allowbreak litro (aprox.) \\
\midrule\noalign{}
\endhead
\bottomrule\noalign{}
\endlastfoot
Gasolina 93 oct & 0.97 & 7,178 \\
Gasolina 97 oct (premium) & 1.07 & 7,918 \\
Diesel & 0.90 & 6,660 \\
\end{longtable}

\begin{quote}
Preços podem variar ±5\% conforme posto e sazonalidade. Chaco e interior
remoto apresentam maior variação.
\end{quote}

\subsubsection{Cobertura Celular}

\textbf{Fonte:} CONATEL PY /\allowbreak  operadoras (2024)

\begin{longtable}[]{@{}ll@{}}
\toprule\noalign{}
Parâmetro & Valor \\
\midrule\noalign{}
\endhead
\bottomrule\noalign{}
\endlastfoot
Cobertura 4G população (dept.) & 95\% \\
Cobertura 4G área rural & 88\% \\
Melhor operadora & Tigo/\allowbreak Personal \\
Qualidade rural & boa \\
\end{longtable}

\begin{quote}
Para áreas rurais fora do núcleo urbano, recomenda-se chip Tigo como
principal e Personal como backup.
\end{quote}

\subsubsection{Internet}

\textbf{Fonte:} CONATEL /\allowbreak  Speedtest Ookla (2024)

\begin{longtable}[]{@{}ll@{}}
\toprule\noalign{}
Parâmetro & Valor \\
\midrule\noalign{}
\endhead
\bottomrule\noalign{}
\endlastfoot
Velocidade média download & 85 Mbps \\
Domicílios com internet (dept.) & 87\% \\
Tecnologia predominante & fibra/\allowbreak cabo \\
Opção rural & Starlink disponível (\textasciitilde USD 44/\allowbreak mês) \\
\end{longtable}

\subsubsection{Mercado Imobiliário e Terra
Rural}

\textbf{Fonte:} INDERT /\allowbreak  Clasificados.com.py (2024)

\begin{longtable}[]{@{}ll@{}}
\toprule\noalign{}
Tipo & Referência \\
\midrule\noalign{}
\endhead
\bottomrule\noalign{}
\endlastfoot
Terra agrícola alta prod. (USD/\allowbreak ha) & 14,800 \\
Imóvel urbano (USD/\allowbreak m²) & 1,400 \\
Aluguel 2 quartos (USD/\allowbreak mês) & 600 \\
\end{longtable}

\begin{quote}
Valores de referência departamental. Localidades menores podem ter
preços 20--40\% abaixo da capital departamental. \#\#\# Saúde
\end{quote}

\textbf{Fonte:} MSPBS /\allowbreak  IPS Paraguay (2024-2026), consolidação
departamental e proxy local

\begin{longtable}[]{@{}ll@{}}
\toprule\noalign{}
Serviço & Disponibilidade \\
\midrule\noalign{}
\endhead
\bottomrule\noalign{}
\endlastfoot
USF /\allowbreak  Posto de Saúde & sim \\
Hospital Regional & não \\
IPS (seguro social) & sim \\
Farmácia & sim \\
Distância ao hospital de referência & \textasciitilde24 km (Ciudad del
Este) \\
\end{longtable}

\textbf{Principais estabelecimentos:} USF local; referencia hospitalar
em Ciudad del Este

\textbf{Observação para imigrantes:} Atendimento primario local ou em
raio curto; casos de maior complexidade seguem para Ciudad del Este.
Cobertura privada continua recomendada para especialidades e urgências
de maior complexidade.
\newpage
\secaoDiagnostico{Mbaracayu}{\mapaConteudo{-24.8833}{-54.7166}}{Mbaracayú é uma das localidades que oferece o maior potencial de
privacidade demográfica e isolamento tático no norte de Alto Paraná. Sua
força reside na baixa visibilidade estratégica provida pela presença de
grandes reservas biológicas e na baixíssima densidade populacional.
Oferece solo de elite mundial (``terra roxa'') e abundância de recursos
hídricos puros, sendo o local ideal para quem busca o máximo de
discrição e autossuficiência familiar, aceitando a necessidade de
autonomia logística e de saúde devido à infraestrutura institucional
local limitada.

\subsubsection{Por que sim}

\begin{itemize}
\tightlist
\item
  Densidade populacional extremamente baixa (privacidade total).
\item
  Solo de altíssima produtividade e ambiente cercado por reservas
  naturais.
\item
  Localização ``fora do radar'' das grandes crises demográficas e
  militares.
\item
  Sem restrições da Lei de Fronteira.
\end{itemize}

\subsubsection{Por que não}

\begin{itemize}
\tightlist
\item
  Infraestrutura de saúde local limitada a atendimentos primários.
\item
  Dependência de estradas vicinais de terra para acesso interno às
  propriedades.
\item
  Pequena escala comercial exige deslocamento para centros maiores
  (Hernandarias).
\end{itemize}

\subsubsection{Recomendação
Técnica}

Altamente indicado para estabelecimentos de autossuficiência de médio
porte que busquem isolamento demográfico massivo e abundância de
recursos naturais. Requer autonomia em logística básica e saúde. O GSS
de 6.6 reflete a segurança do isolamento em contraste com a limitação
institucional local. Referencia atual de score: GSS 6.6 em 2026-03-05.

\subsubsection{}

\begin{itemize}
\tightlist
\item
  \begin{enumerate}
  \def\labelenumi{\arabic{enumi})}
  \tightlist
  \item
    Dados censitários validados (INE\footnote{Instituto Nacional de Estadística (Instituto Nacional de Estatística), órgão oficial de dados demográficos e censitários.} 2022).
  \end{enumerate}
\item
  \begin{enumerate}
  \def\labelenumi{\arabic{enumi})}
  \setcounter{enumi}{1}
  \tightlist
  \item
    Relatórios de monitoramento ambiental da Reserva Itabó
    (MADES/\allowbreak Itaipu).
  \end{enumerate}
\item
  \begin{enumerate}
  \def\labelenumi{\arabic{enumi})}
  \setcounter{enumi}{2}
  \tightlist
  \item
    Mapa de solos e aptidão agrícola de Alto Paraná (MAG).
  \end{enumerate}
\item
  \begin{enumerate}
  \def\labelenumi{\arabic{enumi})}
  \setcounter{enumi}{3}
  \tightlist
  \item
    Registro de segurança pública departamental.
  \end{enumerate}
\end{itemize}}
\begin{table}[ht!]
\small \begin{tabular}{p{3.0cm}p{0.8cm}p{6.0cm}} \toprule
\textbf{Critério} & \textbf{Nota} & \textbf{Justificativa} \\ \midrule
A - Ameaças Estratégicas & 6.5 & Localização remota e protegida por reservas naturais; excelente invisibilidade tática \\
B - Risco Social & 7.5 & Baixíssima densidade demográfica; ideal para isolamento social absoluto \\
C - Riscos Naturais & 7.0 & Geologia muito estável e baixo risco de grandes desastres naturais \\
D - Autossuficiência & 7.0 & Bons recursos naturais: solo fértil, gado e água abundante, limitado pela logística de acesso \\
E - Institucional & 5.0 & Município estável seguro, mas com infraestrutura institucional e de saúde limitada \\
\midrule \textbf{GSS FINAL} & \textbf{6.6} & \textbf{Moderadamente Seguro (Ideal para quem busca o máximo de isolamento demográfico em território próspero e seguro).} \\ \bottomrule \end{tabular}
\end{table}
\subsubsection{Vulnerabilidades e
Mitigação}\label{vuln-Mbaracayu-vulnerabilidades-e-mitigauxe7uxe3o}

\begin{itemize}
\tightlist
\item
  \textbf{Isolamento de Saúde:} Distância significativa de hospitais
  cirúrgicos (Mitigação: manutenção de estoque médico básico e kit
  trauma local).
\item
  \textbf{Dependência Logística:} Necessidade de deslocamento para
  cidades maiores para insumos técnicos (Mitigação: estoque de
  suprimentos para 60 dias e posse de veículo 4x4 robusto).
\item
  \textbf{Vulnerabilidade Energética:} Rede ANDE\footnote{Administración Nacional de Electricidade (Administração Nacional de Eletricidade), companhia estatal de energia.} rural instável em
  temporais (Mitigação: instalação de sistemas solares fotovoltaicos
  redundantes).
\end{itemize}
\subsubsection{Dossiê de Campo}\label{dossie-Mbaracayu-dossiuxea-de-campo}
\begin{description}[style=nextline,leftmargin=0.5cm,font=\bfseries]
\item[Coordenadas]  24°53'S 54°43'W.
\item[Topografia]  Relevo ondulado com extensas planícies mecanizáveis e áreas de mata nativa preservada. Altitude \textasciitilde 240m. Área vasta de \textasciitilde 715 km². Geologicamente muito estável, risco sísmico desprezível.
\item[Alvos Estratégicos]  Reserva Biológica Itabó (Recurso estratégico de biodiversidade e água pura); Áreas de produção agrícola intensiva de soja; Proximidade tática com o braço norte do reservatório de Itaipu. Ausência de alvos militares de grande porte, oferecendo excelente invisibilidade estratégica.
\item[Fallout]  Ventos predominantes do Norte (NE) e Leste. Baixo risco de fallout direto.
\item[População]  6.406 habitantes (Censo 2022 Final). Perfil demográfico essencialmente rural (90\%) com forte cultura de agronegócio sustentável e preservação.
\item[Densidade]  \textasciitilde 8,9 hab/\allowbreak km² (Densidade extremamente baixa, oferecendo isolamento tático superior e privacidade estratégica absoluta fora dos pequenos núcleos povoados).
\item[Vias de Acesso]  Acesso via ramais da Ruta PY07. Localização geográfica isolada em relação aos grandes eixos de trânsito nacional, funcionando como um enclave seguro no norte do departamento. Estradas internas de terra exigem veículos 4x4 em períodos chuvosos (média 1.700 mm/\allowbreak ano).
\item[Serviços]  Unidade de Saúde da Família (USF) local. Dependência crítica de San Alberto ou Hernandarias (60 km) para saúde de alta complexidade e logística avançada.
\item[Custo de Vida]  Muito baixo; economia estritamente vinculada à exportação de grãos e pecuária.
\item[Preço da Terra]  US\$ 8.000-12.000/\allowbreak ha (terras mecanizadas de elite); US\$ 2.000-4.000/\allowbreak ha (áreas de campo natural/\allowbreak mata).
\item[Hidrologia]  Baixo risco de inundações sistêmicas catastróficas devido ao relevo e à preservação da cobertura vegetal. Presença de arroios perenes de alta qualidade hídrica integrados à Reserva Itabó.
\item[Clima]  Subtropical Úmido. Verões intensos e invernos com geadas ocasionais.
\item[Energia]  Rede nacional (Itaipu); instável em áreas rurais remotas.
\item[Água]  Abundante via poços artesianos profundos e nascentes protegidas; excelente qualidade hídrica subterrânea.
\item[Qualidade do Solo]  Latossolo Vermelho de Elite; solo profundo, estruturado e de altíssima fertilidade natural. Excelente aptidão para soja, milho e reflorestamento.
\item[Resources Locais]  Grande produção de soja, milho e gado de corte. Elevadíssimo potencial para autossuficiência alimentar básica em nível de propriedade devido à vasta área produtiva e baixa população.
\item[Segurança]  Zona rural extremamente pacífica e estável. Baixíssima criminalidade urbana. Coesão social baseada na cultura produtora e no respeito às áreas de reserva. Baixa visibilidade para conflitos sociais de massa ou agitação urbana.
\item[Leis Local: Município consolidado. Livre de Restrição de Fronteira]  Fora da faixa de 50km da fronteira internacional seca (limite fluvial no reservatório).
\end{description}
\paragraph{Solo (SoilGrids 2.0, média ponderada 0--30
cm)}

\begin{longtable}[]{@{}ll@{}}
\toprule\noalign{}
Parâmetro & Valor \\
\midrule\noalign{}
\endhead
\bottomrule\noalign{}
\endlastfoot
pH (H$_2$O) & 5.4 \\
Carbono orgânico (SOC) & 22.47 g/\allowbreak kg \\
Argila & 48.28 \% \\
Areia & 23.78 \% \\
Silte (calc.) & 27.9 \% \\
Densidade aparente & 1.24 g/\allowbreak cm³ \\
\textbf{Aptidão agrícola} & \textbf{Média} \\
\end{longtable}

Fonte: ISRIC SoilGrids 2.0 via WCS (acesso em 2026-03-21). Coords:
-24.8833°, -54.7167°. Média ponderada camadas 0-5, 5-15, 15-30 cm.
\subsubsection{Indicadores Sociais}

\textbf{Fontes:} PNUD 2020, INE\footnote{Instituto Nacional de Estadística (Instituto Nacional de Estatística), órgão oficial de dados demográficos e censitários.} Censo 2022, INE\footnote{Instituto Nacional de Estadística (Instituto Nacional de Estatística), órgão oficial de dados demográficos e censitários.} EPHC 2023, Ministerio
Público 2024\\
\textbf{Nota:} Valores marcados como \emph{(dept.)} referem-se ao
departamento de Alto Paraná; valores \emph{(dist.)} são específicos
deste distrito. Dados marcados \emph{(est.)} são estimativas.

\begin{longtable}[]{@{}
  >{\raggedright\arraybackslash}p{(\linewidth - 4\tabcolsep) * \real{0.4231}}
  >{\raggedright\arraybackslash}p{(\linewidth - 4\tabcolsep) * \real{0.2692}}
  >{\raggedright\arraybackslash}p{(\linewidth - 4\tabcolsep) * \real{0.3077}}@{}}
\toprule\noalign{}
\begin{minipage}[b]{\linewidth}\raggedright
Indicador
\end{minipage} & \begin{minipage}[b]{\linewidth}\raggedright
Valor
\end{minipage} & \begin{minipage}[b]{\linewidth}\raggedright
Âmbito
\end{minipage} \\
\midrule\noalign{}
\endhead
\bottomrule\noalign{}
\endlastfoot
IDH (2020) & 0,752 & dept. \\
Ranking IDH nacional & 3/\allowbreak 18 & dept. \\
Esperança de vida & 75,1 anos & dept. \\
Escolaridade média & 7.0 anos & dist. \\
RNB per capita (USD PPA) & 13.500 & dept. \\
Pobreza monetária (\%) & 16,3\% & dept. \\
Pobreza extrema (\%) & 3,1\% & dept. \\
Índice de Gini & 0,428 & dept. \\
Acesso a água potável (\%) & 88,7\% & dept. \\
Acesso a saneamento (\%) & 87,2\% & dept. \\
Taxa de homicídios (est., /\allowbreak 100k hab) & \textasciitilde8--12 (estimativa,
acima da média nacional) & dept. \\
Índice de segurança & baixo & dept. \\
População (Censo 2022) & 6.406 hab. & dist. \\
Idade mediana & 27.0 anos & dist. \\
Taxa de fecundidade & N/\allowbreak D & dist. \\
\end{longtable}
\subsubsection{Combustível}

\textbf{Referência:} PETROPAR /\allowbreak  postos locais (2024) \textbf{Tipo de
localidade:} Interior

\begin{longtable}[]{@{}lll@{}}
\toprule\noalign{}
Combustível & USD/\allowbreak litro & Gs/\allowbreak litro (aprox.) \\
\midrule\noalign{}
\endhead
\bottomrule\noalign{}
\endlastfoot
Gasolina 93 oct & 0.97 & 7,178 \\
Gasolina 97 oct (premium) & 1.07 & 7,918 \\
Diesel & 0.90 & 6,660 \\
\end{longtable}

\begin{quote}
Preços podem variar ±5\% conforme posto e sazonalidade. Chaco e interior
remoto apresentam maior variação.
\end{quote}

\subsubsection{Cobertura Celular}

\textbf{Fonte:} CONATEL PY /\allowbreak  operadoras (2024)

\begin{longtable}[]{@{}ll@{}}
\toprule\noalign{}
Parâmetro & Valor \\
\midrule\noalign{}
\endhead
\bottomrule\noalign{}
\endlastfoot
Cobertura 4G população (dept.) & 95\% \\
Cobertura 4G área rural & 88\% \\
Melhor operadora & Tigo/\allowbreak Personal \\
Qualidade rural & boa \\
\end{longtable}

\begin{quote}
Para áreas rurais fora do núcleo urbano, recomenda-se chip Tigo como
principal e Personal como backup.
\end{quote}

\subsubsection{Internet}

\textbf{Fonte:} CONATEL /\allowbreak  Speedtest Ookla (2024)

\begin{longtable}[]{@{}ll@{}}
\toprule\noalign{}
Parâmetro & Valor \\
\midrule\noalign{}
\endhead
\bottomrule\noalign{}
\endlastfoot
Velocidade média download & 85 Mbps \\
Domicílios com internet (dept.) & 87\% \\
Tecnologia predominante & fibra/\allowbreak cabo \\
Opção rural & Starlink disponível (\textasciitilde USD 44/\allowbreak mês) \\
\end{longtable}

\subsubsection{Mercado Imobiliário e Terra
Rural}

\textbf{Fonte:} INDERT /\allowbreak  Clasificados.com.py (2024)

\begin{longtable}[]{@{}ll@{}}
\toprule\noalign{}
Tipo & Referência \\
\midrule\noalign{}
\endhead
\bottomrule\noalign{}
\endlastfoot
Terra agrícola alta prod. (USD/\allowbreak ha) & 14,800 \\
Imóvel urbano (USD/\allowbreak m²) & 1,400 \\
Aluguel 2 quartos (USD/\allowbreak mês) & 600 \\
\end{longtable}

\begin{quote}
Valores de referência departamental. Localidades menores podem ter
preços 20--40\% abaixo da capital departamental. \#\#\# Saúde
\end{quote}

\textbf{Fonte:} MSPBS /\allowbreak  IPS Paraguay (2024-2026), consolidação
departamental e proxy local

\begin{longtable}[]{@{}ll@{}}
\toprule\noalign{}
Serviço & Disponibilidade \\
\midrule\noalign{}
\endhead
\bottomrule\noalign{}
\endlastfoot
USF /\allowbreak  Posto de Saúde & sim \\
Hospital Regional & não \\
IPS (seguro social) & sim \\
Farmácia & sim \\
Distância ao hospital de referência & \textasciitilde73 km (Ciudad del
Este) \\
\end{longtable}

\textbf{Principais estabelecimentos:} USF local; referencia hospitalar
em Ciudad del Este

\textbf{Observação para imigrantes:} Atendimento primario local ou em
raio curto; casos de maior complexidade seguem para Ciudad del Este.
Cobertura privada continua recomendada para especialidades e urgências
de maior complexidade.
\newpage
\secaoDiagnostico{Minga Guazu}{\mapaConteudo{-25.4833}{-54.7500}}{Minga Guazú é o motor logístico e industrial de Alto Paraná. Oferece
infraestrutura de serviços de primeiro mundo, solos de altíssima
produtividade e conectividade ímpar através do aeroporto e rodovias
duplicadas. Sua maior força reside na capacidade de processamento de
recursos, enquanto sua principal fraqueza é a altíssima visibilidade
estratégica como alvo de controle nacional em situações de conflito. É o
local ideal para quem busca integrar-se a um polo produtivo vibrante,
exigindo do ocupante autonomia tática para mitigar os riscos de sua
importância geopolítica.

\subsubsection{Por que sim}

\begin{itemize}
\tightlist
\item
  Solo de altíssima produtividade e clima excepcional para policultivos.
\item
  Melhor infraestrutura de transporte aéreo e rodoviário do interior.
\item
  Abundância de recursos industriais básicos e alimentos processados.
\item
  Sem restrições da Lei de Fronteira.
\end{itemize}

\subsubsection{Por que não}

\begin{itemize}
\tightlist
\item
  Alvo estratégico prioritário nacional (Aeroporto e Nó Rodoviário).
\item
  Densidade populacional urbana elevada para um refúgio isolado.
\item
  Maior vulnerabilidade social em cenários de bloqueio rodoviário.
\end{itemize}

\subsubsection{Recomendação
Técnica}

Altamente indicado como base logística ou industrial para o projeto.
Para moradia permanente, recomenda-se a zona rural do distrito, mantendo
autonomia energética e hídrica total. O GSS de 6.0 reflete o equilíbrio
entre o poder dos recursos e a extrema visibilidade estratégica da
localidade. Referencia atual de score: GSS 6.0 em 2026-03-05.

\subsubsection{}

\begin{itemize}
\tightlist
\item
  \begin{enumerate}
  \def\labelenumi{\arabic{enumi})}
  \tightlist
  \item
    Dados censitários validados (INE\footnote{Instituto Nacional de Estadística (Instituto Nacional de Estatística), órgão oficial de dados demográficos e censitários.} 2022).
  \end{enumerate}
\item
  \begin{enumerate}
  \def\labelenumi{\arabic{enumi})}
  \setcounter{enumi}{1}
  \tightlist
  \item
    Relatórios de operação aeroportuária (DINAC 2025).
  \end{enumerate}
\item
  \begin{enumerate}
  \def\labelenumi{\arabic{enumi})}
  \setcounter{enumi}{2}
  \tightlist
  \item
    Mapa de solos e aptidão agrícola mecanizada (MAG).
  \end{enumerate}
\item
  \begin{enumerate}
  \def\labelenumi{\arabic{enumi})}
  \setcounter{enumi}{3}
  \tightlist
  \item
    Registro de segurança pública e governança industrial regional.
  \end{enumerate}
\end{itemize}}
\begin{table}[ht!]
\small \begin{tabular}{p{3.0cm}p{0.8cm}p{6.0cm}} \toprule
\textbf{Critério} & \textbf{Nota} & \textbf{Justificativa} \\ \midrule
A - Ameaças Estratégicas & 4.0 & Hub estratégico tático: aeroporto internacional e nó logístico nacional vital; alta visibilidade \\
B - Risco Social & 5.0 & População elevada e densidade urbana significativa; risco social em crises de massa \\
C - Riscos Naturais & 7.0 & Geologia muito estável e baixo risco de grandes desastres naturais \\
D - Autossuficiência & 8.0 & Recursos agroindustriais excepcionais: solo de elite e polo de processamento de grãos \\
E - Institucional & 7.0 & Forte segurança institucional e serviços completos na região metropolitana \\
\midrule \textbf{GSS FINAL} & \textbf{6.0} & \textbf{Moderadamente Seguro (Ideal como base logística e agroindustrial, requer preparo para visibilidade estratégica).} \\ \bottomrule \end{tabular}
\end{table}
\subsubsection{Vulnerabilidades e
Mitigação}\label{vuln-Minga_Guazu-vulnerabilidades-e-mitigauxe7uxe3o}

\begin{itemize}
\tightlist
\item
  \textbf{Alvo Estratégico Aeroportuário:} O aeroporto é um ponto de
  controle crítico em crises (Mitigação: residência rural afastada 10-15
  km do perímetro do aeroporto).
\item
  \textbf{Funil Logístico:} Risco de isolamento por bloqueios nas Rutas
  PY02/\allowbreak PY06 (Mitigação: residência rural com autonomia total e rotas de
  terra alternativas mapeadas).
\item
  \textbf{Concentração Populacional:} Risco de instabilidade social
  urbana em crises de suprimento (Mitigação: estabelecimento de base com
  fortificação tática e sistema de água próprio).
\end{itemize}
\subsubsection{Dossiê de Campo}\label{dossie-Minga_Guazu-dossiuxea-de-campo}
\begin{description}[style=nextline,leftmargin=0.5cm,font=\bfseries]
\item[Coordenadas]  25°29'S 54°45'W.
\item[Topografia]  Relevo predominantemente plano a suavemente ondulado. Localizado no eixo metropolitano de Ciudad del Este. Altitude \textasciitilde 240m. Área de \textasciitilde 489 km². Geologicamente muito estável, risco sísmico desprezível.
\item[Alvos Estratégicos]  Aeroporto Internacional Guaraní (Segundo maior do país, infraestrutura crítica de aviação); Entroncamento logístico nevrálgico das Rutas PY02 e PY06; Grandes complexos agroindustriais e silos; Subestação ANDE\footnote{Administración Nacional de Electricidade (Administração Nacional de Eletricidade), companhia estatal de energia.}. Alta importância tática e militar nacional.
\item[Fallout]  Ventos predominantes do Norte (NE) e Leste. Baixo risco de fallout direto de alvos continentais distantes.
\item[População]  81.072 habitantes (Censo 2022 Final). População cosmopolita com forte componente demográfico de agricultores e operários industriais.
\item[Densidade]  \textasciitilde 165,8 hab/\allowbreak km² (Densidade urbana elevada, indicando alta vulnerabilidade social e pressão sobre serviços em crises de massa).
\item[Vias de Acesso]  Ponto nevrálgico de conexão entre o Leste (CDE), o Sul (Itapúa) e o Oeste (Assunção). Infraestrutura asfáltica de elite (PY02 duplicada). Risco extremo de bloqueios táticos e saturação rodoviária em cenários de instabilidade nacional.
\item[Serviços]  Infraestrutura de serviços completa. Centro médico local e proximidade imediata com hospitais de alta complexidade em Ciudad del Este (15-20 km). Polo comercial agroindustrial robusto.
\item[Custo de Vida]  Médio; economia pujante baseada na exportação agrícola e serviços logísticos.
\item[Preço da Terra]  US\$ 8.000-15.000/\allowbreak ha (áreas mecanizadas de elite); lotes industriais e residenciais urbanos valorizados.
\item[Hidrologia]  Diversos arroios cruzam o distrito (ex: Rio Monday). Baixo risco de inundações sistêmicas catastróficas. Vulnerabilidade localizada em drenagem industrial e urbana durante tempestades severas (média 1.700 mm/\allowbreak ano).
\item[Clima]  Subtropical Úmido. Verões quentes; invernos com geadas ocasionais.
\item[Energia]  Rede nacional (Itaipu); alta estabilidade devido à proximidade com a fonte geradora e importância industrial.
\item[Água]  Abundante via poços artesianos profundos e sistemas de captação local (Aquífero Guarani).
\item[Qualidade do Solo]  Latossolo Vermelho de Elite; profundo e extremamente fértil. Excelente aptidão para soja, milho e trigo.
\item[Resources Locais]  Grande polo produtor de grãos, proteína animal e sementes. Elevado potencial para autossuficiência alimentar regional, mas dependente de ordem logística para distribuição.
\item[Segurança]  Cidade dinâmica com criminalidade urbana típica de polos logísticos. Segurança institucional presente, mas sob vigilância tática constante devido ao aeroporto internacional. Estabilidade social baseada na cultura cooperativista e industrial ("Minga" = trabalho comunitário).
\item[Leis Local]  Município consolidado e integrado ao desenvolvimento do Alto Paraná. \\ \textbf{Livre de Restrição} de Fronteira. \\ \textit{Aquisição de terra rural proibida para estrangeiros limítrofes.}
\end{description}
\paragraph{Solo (SoilGrids 2.0, média ponderada 0--30
cm)}

\begin{longtable}[]{@{}ll@{}}
\toprule\noalign{}
Parâmetro & Valor \\
\midrule\noalign{}
\endhead
\bottomrule\noalign{}
\endlastfoot
pH (H$_2$O) & 5.32 \\
Carbono orgânico (SOC) & 22.18 g/\allowbreak kg \\
Argila & 47.21 \% \\
Areia & 19.22 \% \\
Silte (calc.) & 33.6 \% \\
Densidade aparente & 1.25 g/\allowbreak cm³ \\
\textbf{Aptidão agrícola} & \textbf{Média} \\
\end{longtable}

Fonte: ISRIC SoilGrids 2.0 via WCS (acesso em 2026-03-21). Coords:
-25.4833°, -54.75°. Média ponderada camadas 0-5, 5-15, 15-30 cm.
\subsubsection{Indicadores Sociais}

\textbf{Fontes:} PNUD 2020, INE\footnote{Instituto Nacional de Estadística (Instituto Nacional de Estatística), órgão oficial de dados demográficos e censitários.} Censo 2022, INE\footnote{Instituto Nacional de Estadística (Instituto Nacional de Estatística), órgão oficial de dados demográficos e censitários.} EPHC 2023, Ministerio
Público 2024\\
\textbf{Nota:} Valores marcados como \emph{(dept.)} referem-se ao
departamento de Alto Paraná; valores \emph{(dist.)} são específicos
deste distrito. Dados marcados \emph{(est.)} são estimativas.

\begin{longtable}[]{@{}
  >{\raggedright\arraybackslash}p{(\linewidth - 4\tabcolsep) * \real{0.4231}}
  >{\raggedright\arraybackslash}p{(\linewidth - 4\tabcolsep) * \real{0.2692}}
  >{\raggedright\arraybackslash}p{(\linewidth - 4\tabcolsep) * \real{0.3077}}@{}}
\toprule\noalign{}
\begin{minipage}[b]{\linewidth}\raggedright
Indicador
\end{minipage} & \begin{minipage}[b]{\linewidth}\raggedright
Valor
\end{minipage} & \begin{minipage}[b]{\linewidth}\raggedright
Âmbito
\end{minipage} \\
\midrule\noalign{}
\endhead
\bottomrule\noalign{}
\endlastfoot
IDH (2020) & 0,752 & dept. \\
Ranking IDH nacional & 3/\allowbreak 18 & dept. \\
Esperança de vida & 75,1 anos & dept. \\
Escolaridade média & 9.2 anos & dist. \\
RNB per capita (USD PPA) & 13.500 & dept. \\
Pobreza monetária (\%) & 16,3\% & dept. \\
Pobreza extrema (\%) & 3,1\% & dept. \\
Índice de Gini & 0,428 & dept. \\
Acesso a água potável (\%) & 88,7\% & dept. \\
Acesso a saneamento (\%) & 87,2\% & dept. \\
Taxa de homicídios (est., /\allowbreak 100k hab) & \textasciitilde8--12 (estimativa,
acima da média nacional) & dept. \\
Índice de segurança & baixo & dept. \\
População (Censo 2022) & 81.072 hab. & dist. \\
Idade mediana & 27.0 anos & dist. \\
Taxa de fecundidade & N/\allowbreak D & dist. \\
\end{longtable}
\subsubsection{Combustível}

\textbf{Referência:} PETROPAR /\allowbreak  postos locais (2024) \textbf{Tipo de
localidade:} Interior

\begin{longtable}[]{@{}lll@{}}
\toprule\noalign{}
Combustível & USD/\allowbreak litro & Gs/\allowbreak litro (aprox.) \\
\midrule\noalign{}
\endhead
\bottomrule\noalign{}
\endlastfoot
Gasolina 93 oct & 0.97 & 7,178 \\
Gasolina 97 oct (premium) & 1.07 & 7,918 \\
Diesel & 0.90 & 6,660 \\
\end{longtable}

\begin{quote}
Preços podem variar ±5\% conforme posto e sazonalidade. Chaco e interior
remoto apresentam maior variação.
\end{quote}

\subsubsection{Cobertura Celular}

\textbf{Fonte:} CONATEL PY /\allowbreak  operadoras (2024)

\begin{longtable}[]{@{}ll@{}}
\toprule\noalign{}
Parâmetro & Valor \\
\midrule\noalign{}
\endhead
\bottomrule\noalign{}
\endlastfoot
Cobertura 4G população (dept.) & 95\% \\
Cobertura 4G área rural & 88\% \\
Melhor operadora & Tigo/\allowbreak Personal \\
Qualidade rural & boa \\
\end{longtable}

\begin{quote}
Para áreas rurais fora do núcleo urbano, recomenda-se chip Tigo como
principal e Personal como backup.
\end{quote}

\subsubsection{Internet}

\textbf{Fonte:} CONATEL /\allowbreak  Speedtest Ookla (2024)

\begin{longtable}[]{@{}ll@{}}
\toprule\noalign{}
Parâmetro & Valor \\
\midrule\noalign{}
\endhead
\bottomrule\noalign{}
\endlastfoot
Velocidade média download & 85 Mbps \\
Domicílios com internet (dept.) & 87\% \\
Tecnologia predominante & fibra/\allowbreak cabo \\
Opção rural & Starlink disponível (\textasciitilde USD 44/\allowbreak mês) \\
\end{longtable}

\subsubsection{Mercado Imobiliário e Terra
Rural}

\textbf{Fonte:} INDERT /\allowbreak  Clasificados.com.py (2024)

\begin{longtable}[]{@{}ll@{}}
\toprule\noalign{}
Tipo & Referência \\
\midrule\noalign{}
\endhead
\bottomrule\noalign{}
\endlastfoot
Terra agrícola alta prod. (USD/\allowbreak ha) & 14,800 \\
Imóvel urbano (USD/\allowbreak m²) & 1,400 \\
Aluguel 2 quartos (USD/\allowbreak mês) & 600 \\
\end{longtable}

\begin{quote}
Valores de referência departamental. Localidades menores podem ter
preços 20--40\% abaixo da capital departamental. \#\#\# Saúde
\end{quote}

\textbf{Fonte:} MSPBS /\allowbreak  IPS Paraguay (2024-2026), consolidação
departamental e proxy local

\begin{longtable}[]{@{}ll@{}}
\toprule\noalign{}
Serviço & Disponibilidade \\
\midrule\noalign{}
\endhead
\bottomrule\noalign{}
\endlastfoot
USF /\allowbreak  Posto de Saúde & sim \\
Hospital Regional & sim \\
IPS (seguro social) & sim \\
Farmácia & sim \\
Distância ao hospital de referência & \textasciitilde19 km (Ciudad del
Este) \\
\end{longtable}

\textbf{Principais estabelecimentos:} USF local; referencia hospitalar
em Ciudad del Este

\textbf{Observação para imigrantes:} Atendimento primario local ou em
raio curto; casos de maior complexidade seguem para Ciudad del Este.
Cobertura privada continua recomendada para especialidades e urgências
de maior complexidade.
\newpage
\secaoDiagnostico{Minga Pora}{\mapaConteudo{-24.8833}{-54.7500}}{Minga Porá é um dos enclaves de resiliência e produtividade mais
equilibrados do norte de Alto Paraná. Oferece solos de produtividade
mundial e uma abundância excepcional de recursos naturais (água e
terra), funcionando como um santuário de segurança alimentar e hídrica.
Sua força absoluta reside na combinação de isolamento demográfico com
uma base produtiva de alta performance. É o local ideal para quem busca
autonomia total em ambiente de paz social, transformando a riqueza de
recursos básicos em uma barreira de proteção contra crises externas.

\subsubsection{Por que sim}

\begin{itemize}
\tightlist
\item
  Solos de altíssima produtividade e abundância de águas puras
  superficiais.
\item
  Baixíssima densidade populacional (privacidade absoluta).
\item
  Localização estratégica fora do radar de fallout metropolitano.
\item
  Sem restrições da Lei de Fronteira.
\end{itemize}

\subsubsection{Por que não}

\begin{itemize}
\tightlist
\item
  Infraestrutura de saúde local limitada a atendimentos primários.
\item
  Dependência de grandes centros para serviços hospitalares complexos.
\item
  Pequena escala comercial local exige planejamento logístico de
  suprimentos técnicos.
\end{itemize}

\subsubsection{Recomendação
Técnica}

Altamente indicado para estabelecimentos de autossuficiência tecnificada
que valorizem o isolamento e produtividade. Requer autonomia energética
individual e investimento em segurança patrimonial. O GSS de 7.5 reflete
a solidez dos recursos e o isolamento em equilíbrio com a segurança
institucional local. Referencia atual de score: GSS 7.5 em 2026-03-05.

\subsubsection{}

\begin{itemize}
\tightlist
\item
  \begin{enumerate}
  \def\labelenumi{\arabic{enumi})}
  \tightlist
  \item
    Dados censitários validados (INE\footnote{Instituto Nacional de Estadística (Instituto Nacional de Estatística), órgão oficial de dados demográficos e censitários.} 2022).
  \end{enumerate}
\item
  \begin{enumerate}
  \def\labelenumi{\arabic{enumi})}
  \setcounter{enumi}{1}
  \tightlist
  \item
    Relatórios de produção agrícola mecanizada departamental (MAG).
  \end{enumerate}
\item
  \begin{enumerate}
  \def\labelenumi{\arabic{enumi})}
  \setcounter{enumi}{2}
  \tightlist
  \item
    Mapa de solos e bacias hídricas de Alto Paraná (Portal
    Geoestadístico).
  \end{enumerate}
\item
  \begin{enumerate}
  \def\labelenumi{\arabic{enumi})}
  \setcounter{enumi}{3}
  \tightlist
  \item
    Registro de segurança pública e governança municipal.
  \end{enumerate}
\end{itemize}}
\begin{table}[ht!]
\small \begin{tabular}{p{3.0cm}p{0.8cm}p{6.0cm}} \toprule
\textbf{Critério} & \textbf{Nota} & \textbf{Justificativa} \\ \midrule
A - Ameaças Estratégicas & 7.0 & Localização remota e discreta; excelente invisibilidade tática fora dos eixos de conflito primários \\
B - Risco Social & 8.0 & Densidade populacional baixa; isolamento social superior e privacidade absoluta \\
C - Riscos Naturais & 7.0 & Geologia muito estável e baixo risco de grandes desastres naturais \\
D - Autossuficiência & 8.5 & Recursos agrícolas e hídricos excelentes; solo de elite e proximidade com reservatório \\
E - Institucional & 7.0 & Estabilidade institucional presente e segurança regional funcional em ambiente rural \\
\midrule \textbf{GSS FINAL} & \textbf{7.5} & \textbf{Seguro (Altamente recomendado para quem busca isolamento tático rural em terras férteis com abundância de recursos básicos).} \\ \bottomrule \end{tabular}
\end{table}
\subsubsection{Vulnerabilidades e
Mitigação}\label{vuln-Minga_Pora-vulnerabilidades-e-mitigauxe7uxe3o}

\begin{itemize}
\tightlist
\item
  \textbf{Isolamento de Saúde:} Distância significativa de hospitais
  cirúrgicos (Mitigação: manutenção de estoque médico básico e kit
  trauma local).
\item
  \textbf{Dependência de Exportação:} Economia vinculada ao mercado
  global de commodities (Mitigação: diversificação produtiva para
  consumo interno e autonomia energética).
\item
  \textbf{Visibilidade Tática de Suprimentos:} Polo de silos pode ser
  alvo em crises de desabastecimento (Mitigação: residência rural
  afastada 5-10 km do eixo urbano industrial).
\end{itemize}
\subsubsection{Dossiê de Campo}\label{dossie-Minga_Pora-dossiuxea-de-campo}
\begin{description}[style=nextline,leftmargin=0.5cm,font=\bfseries]
\item[Coordenadas]  24°53'S 54°45'W.
\item[Topografia]  Relevo ondulado com extensas planícies mecanizáveis e solos de alta fertilidade (Terra Roxa). Margeado ao leste pelo reservatório de Itaipu. Altitude \textasciitilde 240m. Área vasta de \textasciitilde 881 km². Geologicamente estável, risco sísmico desprezível.
\item[Alvos Estratégicos]  Polo de produção agrícola intensiva (soja, milho); Silos de armazenamento de grande porte; Proximidade tática com o braço norte do reservatório de Itaipu. Ausência de alvos militares diretos, oferecendo excelente invisibilidade estratégica no norte do departamento.
\item[Fallout]  Ventos predominantes do Norte (NE) e Leste. Baixíssimo risco de fallout direto; posição isolada e protegida.
\item[População]  11.959 habitantes (Censo 2022 Final). População com forte herança de imigrantes brasileiros, altamente integrada ao sistema produtivo agroindustrial mecanizado.
\item[Densidade]  \textasciitilde 13,6 hab/\allowbreak km² (Densidade baixa, oferecendo isolamento tático superior e privacidade absoluta fora do pequeno núcleo urbano).
\item[Vias de Acesso]  Localização estratégica sobre a Ruta PY07. Conectividade asfáltica eficiente para Ciudad del Este e Salto del Guairá. Facilidade logística tática em tempos de paz, funcionando como um enclave produtivo seguro.
\item[Serviços]  Infraestrutura básica funcional. Centro de Saúde local e USF. Dependência de San Alberto ou Hernandarias para atendimentos médicos de alta complexidade.
\item[Custo de Vida]  Baixo; economia baseada na exportação de grãos e pecuária.
\item[Preço da Terra]  US\$ 10.000-15.000/\allowbreak ha (terras mecanizadas de elite - "Terra Roxa").
\item[Hidrologia]  Baixo risco de inundações sistêmicas devido à topografia favorável. Presença de arroios perenes e proximidade com o reservatório garantem abundância hídrica superficial.
\item[Clima]  Subtropical Úmido. Verões intensos e chuvosos; invernos amenos com geadas ocasionais benéficas.
\item[Energia]  Rede nacional (Itaipu); instável em áreas rurais remotas.
\item[Água]  Abundante via poços artesianos profundos e proximidade com o reservatório de Itaipu; excelente qualidade hídrica subterrânea.
\item[Qualidade do Solo]  Latossolo Vermelho de Elite; solo profundo, estruturado e extremamente fértil. Excelente aptidão para soja e milho.
\item[Resources Locais]  Grande produção de grãos e proteína animal. Elevadíssimo potencial para autossuficiência alimentar básica em nível de propriedade devido à vasta área produtiva e baixa densidade populacional rural.
\item[Segurança]  Zona rural excepcionalmente pacífica e segura. Baixíssima criminalidade urbana. Coesão social baseada na cultura produtora tradicional ("Minga" = trabalho comunitário). Segurança institucional presente, mas sob vigilância tática de fronteira fluvial.
\item[Leis Local]  Município consolidado e polo de expansão agrícola. \\ \textbf{Livre de Restrição} de Fronteira. \\ \textit{Aquisição de terra rural proibida para estrangeiros limítrofes.}
\end{description}
\paragraph{Solo (SoilGrids 2.0, média ponderada 0--30
cm)}

\begin{longtable}[]{@{}ll@{}}
\toprule\noalign{}
Parâmetro & Valor \\
\midrule\noalign{}
\endhead
\bottomrule\noalign{}
\endlastfoot
pH (H$_2$O) & 5.4 \\
Carbono orgânico (SOC) & 24.47 g/\allowbreak kg \\
Argila & 47.98 \% \\
Areia & 25.24 \% \\
Silte (calc.) & 26.8 \% \\
Densidade aparente & 1.25 g/\allowbreak cm³ \\
\textbf{Aptidão agrícola} & \textbf{Média} \\
\end{longtable}

Fonte: ISRIC SoilGrids 2.0 via WCS (acesso em 2026-03-21). Coords:
-24.8833°, -54.75°. Média ponderada camadas 0-5, 5-15, 15-30 cm.

\begin{itemize}
\tightlist
\item
  \textbf{Homicidios/\allowbreak 100k:} \textasciitilde10,0 (Abaixo da média
  departamental).
\end{itemize}

\subsection{10. Analise de Riscos}

\begin{itemize}
\tightlist
\item
  \textbf{Cenario curto prazo:} Manutenção da tranquilidade rural e
  rentabilidade da produção mecanizada.
\item
  \textbf{Cenario medio prazo:} Impacto de mudanças no controle de águas
  do reservatório sobre a logística local.
\end{itemize}

\subsection{11. Redacao Final}

\subsubsection{Diagnostico Integrado}

Minga Porá é um dos enclaves de resiliência e produtividade mais
equilibrados do norte de Alto Paraná. Oferece solos de produtividade
mundial e uma abundância excepcional de recursos naturais (água e
terra), funcionando como um santuário de segurança alimentar e hídrica.
Sua força absoluta reside na combinação de isolamento demográfico com
uma base produtiva de alta performance. É o local ideal para quem busca
autonomia total em ambiente de paz social, transformando a riqueza de
recursos básicos em uma barreira de proteção contra crises externas.

\subsubsection{Por que sim}

\begin{itemize}
\tightlist
\item
  Solos de altíssima produtividade e abundância de águas puras
  superficiais.
\item
  Baixíssima densidade populacional (privacidade absoluta).
\item
  Localização estratégica fora do radar de fallout metropolitano.
\item
  Sem restrições da Lei de Fronteira.
\end{itemize}

\subsubsection{Por que nao}

\begin{itemize}
\tightlist
\item
  Infraestrutura de saúde local limitada a atendimentos primários.
\item
  Dependência de grandes centros para serviços hospitalares complexos.
\item
  Pequena escala comercial local exige planejamento logístico de
  suprimentos técnicos.
\end{itemize}

\subsubsection{Recomendacao Tecnica}

Altamente indicado para estabelecimentos de autossuficiência tecnificada
que valorizem o isolamento e produtividade. Requer autonomia energética
individual e investimento em segurança patrimonial. O GSS de 7.5 reflete
a solidez dos recursos e o isolamento em equilíbrio com a segurança
institucional local. Referencia atual de score: GSS 7.5 em 2026-03-05.

\subsection{13.}

\begin{itemize}
\tightlist
\item
  \begin{enumerate}
  \def\labelenumi{\arabic{enumi})}
  \tightlist
  \item
    Dados censitários validados (INE\footnote{Instituto Nacional de Estadística (Instituto Nacional de Estatística), órgão oficial de dados demográficos e censitários.} 2022).
  \end{enumerate}
\item
  \begin{enumerate}
  \def\labelenumi{\arabic{enumi})}
  \setcounter{enumi}{1}
  \tightlist
  \item
    Relatórios de produção agrícola mecanizada departamental (MAG).
  \end{enumerate}
\item
  \begin{enumerate}
  \def\labelenumi{\arabic{enumi})}
  \setcounter{enumi}{2}
  \tightlist
  \item
    Mapa de solos e bacias hídricas de Alto Paraná (Portal
    Geoestadístico).
  \end{enumerate}
\item
  \begin{enumerate}
  \def\labelenumi{\arabic{enumi})}
  \setcounter{enumi}{3}
  \tightlist
  \item
    Registro de segurança pública e governança municipal.
  \end{enumerate}
\end{itemize}

\subsubsection{Combustível}

\textbf{Referência:} PETROPAR /\allowbreak  postos locais (2024) \textbf{Tipo de
localidade:} Interior

\begin{longtable}[]{@{}lll@{}}
\toprule\noalign{}
Combustível & USD/\allowbreak litro & Gs/\allowbreak litro (aprox.) \\
\midrule\noalign{}
\endhead
\bottomrule\noalign{}
\endlastfoot
Gasolina 93 oct & 0.97 & 7,178 \\
Gasolina 97 oct (premium) & 1.07 & 7,918 \\
Diesel & 0.90 & 6,660 \\
\end{longtable}

\begin{quote}
Preços podem variar ±5\% conforme posto e sazonalidade. Chaco e interior
remoto apresentam maior variação.
\end{quote}

\subsubsection{Cobertura Celular}

\textbf{Fonte:} CONATEL PY /\allowbreak  operadoras (2024)

\begin{longtable}[]{@{}ll@{}}
\toprule\noalign{}
Parâmetro & Valor \\
\midrule\noalign{}
\endhead
\bottomrule\noalign{}
\endlastfoot
Cobertura 4G população (dept.) & 95\% \\
Cobertura 4G área rural & 88\% \\
Melhor operadora & Tigo/\allowbreak Personal \\
Qualidade rural & boa \\
\end{longtable}

\begin{quote}
Para áreas rurais fora do núcleo urbano, recomenda-se chip Tigo como
principal e Personal como backup.
\end{quote}

\subsubsection{Internet}

\textbf{Fonte:} CONATEL /\allowbreak  Speedtest Ookla (2024)

\begin{longtable}[]{@{}ll@{}}
\toprule\noalign{}
Parâmetro & Valor \\
\midrule\noalign{}
\endhead
\bottomrule\noalign{}
\endlastfoot
Velocidade média download & 85 Mbps \\
Domicílios com internet (dept.) & 87\% \\
Tecnologia predominante & fibra/\allowbreak cabo \\
Opção rural & Starlink disponível (\textasciitilde USD 44/\allowbreak mês) \\
\end{longtable}

\subsubsection{Mercado Imobiliário e Terra
Rural}

\textbf{Fonte:} INDERT /\allowbreak  Clasificados.com.py (2024)

\begin{longtable}[]{@{}ll@{}}
\toprule\noalign{}
Tipo & Referência \\
\midrule\noalign{}
\endhead
\bottomrule\noalign{}
\endlastfoot
Terra agrícola alta prod. (USD/\allowbreak ha) & 14,800 \\
Imóvel urbano (USD/\allowbreak m²) & 1,400 \\
Aluguel 2 quartos (USD/\allowbreak mês) & 600 \\
\end{longtable}

\begin{quote}
Valores de referência departamental. Localidades menores podem ter
preços 20--40\% abaixo da capital departamental. \#\#\# Saúde
\end{quote}

\textbf{Fonte:} MSPBS /\allowbreak  IPS Paraguay (2024-2026), consolidação
departamental e proxy local

\begin{longtable}[]{@{}ll@{}}
\toprule\noalign{}
Serviço & Disponibilidade \\
\midrule\noalign{}
\endhead
\bottomrule\noalign{}
\endlastfoot
USF /\allowbreak  Posto de Saúde & sim \\
Hospital Regional & não \\
IPS (seguro social) & sim \\
Farmácia & sim \\
Distância ao hospital de referência & \textasciitilde98 km (Ciudad del
Este) \\
\end{longtable}

\textbf{Principais estabelecimentos:} USF local; referencia hospitalar
em Ciudad del Este

\textbf{Observação para imigrantes:} Atendimento primario local ou em
raio curto; casos de maior complexidade seguem para Ciudad del Este.
Cobertura privada continua recomendada para especialidades e urgências
de maior complexidade.
\subsubsection{Indicadores Sociais}

\textbf{Fontes:} PNUD 2020, INE\footnote{Instituto Nacional de Estadística (Instituto Nacional de Estatística), órgão oficial de dados demográficos e censitários.} Censo 2022, INE\footnote{Instituto Nacional de Estadística (Instituto Nacional de Estatística), órgão oficial de dados demográficos e censitários.} EPHC 2023, Ministerio
Público 2024\\
\textbf{Nota:} Valores marcados como \emph{(dept.)} referem-se ao
departamento de Alto Paraná; valores \emph{(dist.)} são específicos
deste distrito. Dados marcados \emph{(est.)} são estimativas.

\begin{longtable}[]{@{}
  >{\raggedright\arraybackslash}p{(\linewidth - 4\tabcolsep) * \real{0.4231}}
  >{\raggedright\arraybackslash}p{(\linewidth - 4\tabcolsep) * \real{0.2692}}
  >{\raggedright\arraybackslash}p{(\linewidth - 4\tabcolsep) * \real{0.3077}}@{}}
\toprule\noalign{}
\begin{minipage}[b]{\linewidth}\raggedright
Indicador
\end{minipage} & \begin{minipage}[b]{\linewidth}\raggedright
Valor
\end{minipage} & \begin{minipage}[b]{\linewidth}\raggedright
Âmbito
\end{minipage} \\
\midrule\noalign{}
\endhead
\bottomrule\noalign{}
\endlastfoot
IDH (2020) & 0,752 & dept. \\
Ranking IDH nacional & 3/\allowbreak 18 & dept. \\
Esperança de vida & 75,1 anos & dept. \\
Escolaridade média & 8.0 anos & dist. \\
RNB per capita (USD PPA) & 13.500 & dept. \\
Pobreza monetária (\%) & 16,3\% & dept. \\
Pobreza extrema (\%) & 3,1\% & dept. \\
Índice de Gini & 0,428 & dept. \\
Acesso a água potável (\%) & 88,7\% & dept. \\
Acesso a saneamento (\%) & 87,2\% & dept. \\
Taxa de homicídios (est., /\allowbreak 100k hab) & \textasciitilde8--12 (estimativa,
acima da média nacional) & dept. \\
Índice de segurança & baixo & dept. \\
População (Censo 2022) & 11.959 hab. & dist. \\
Idade mediana & 26.0 anos & dist. \\
Taxa de fecundidade & N/\allowbreak D & dist. \\
\end{longtable}
\subsubsection{Combustível}

\textbf{Referência:} PETROPAR /\allowbreak  postos locais (2024) \textbf{Tipo de
localidade:} Interior

\begin{longtable}[]{@{}lll@{}}
\toprule\noalign{}
Combustível & USD/\allowbreak litro & Gs/\allowbreak litro (aprox.) \\
\midrule\noalign{}
\endhead
\bottomrule\noalign{}
\endlastfoot
Gasolina 93 oct & 0.97 & 7,178 \\
Gasolina 97 oct (premium) & 1.07 & 7,918 \\
Diesel & 0.90 & 6,660 \\
\end{longtable}

\begin{quote}
Preços podem variar ±5\% conforme posto e sazonalidade. Chaco e interior
remoto apresentam maior variação.
\end{quote}

\subsubsection{Cobertura Celular}

\textbf{Fonte:} CONATEL PY /\allowbreak  operadoras (2024)

\begin{longtable}[]{@{}ll@{}}
\toprule\noalign{}
Parâmetro & Valor \\
\midrule\noalign{}
\endhead
\bottomrule\noalign{}
\endlastfoot
Cobertura 4G população (dept.) & 95\% \\
Cobertura 4G área rural & 88\% \\
Melhor operadora & Tigo/\allowbreak Personal \\
Qualidade rural & boa \\
\end{longtable}

\begin{quote}
Para áreas rurais fora do núcleo urbano, recomenda-se chip Tigo como
principal e Personal como backup.
\end{quote}

\subsubsection{Internet}

\textbf{Fonte:} CONATEL /\allowbreak  Speedtest Ookla (2024)

\begin{longtable}[]{@{}ll@{}}
\toprule\noalign{}
Parâmetro & Valor \\
\midrule\noalign{}
\endhead
\bottomrule\noalign{}
\endlastfoot
Velocidade média download & 85 Mbps \\
Domicílios com internet (dept.) & 87\% \\
Tecnologia predominante & fibra/\allowbreak cabo \\
Opção rural & Starlink disponível (\textasciitilde USD 44/\allowbreak mês) \\
\end{longtable}

\subsubsection{Mercado Imobiliário e Terra
Rural}

\textbf{Fonte:} INDERT /\allowbreak  Clasificados.com.py (2024)

\begin{longtable}[]{@{}ll@{}}
\toprule\noalign{}
Tipo & Referência \\
\midrule\noalign{}
\endhead
\bottomrule\noalign{}
\endlastfoot
Terra agrícola alta prod. (USD/\allowbreak ha) & 14,800 \\
Imóvel urbano (USD/\allowbreak m²) & 1,400 \\
Aluguel 2 quartos (USD/\allowbreak mês) & 600 \\
\end{longtable}

\begin{quote}
Valores de referência departamental. Localidades menores podem ter
preços 20--40\% abaixo da capital departamental. \#\#\# Saúde
\end{quote}

\textbf{Fonte:} MSPBS /\allowbreak  IPS Paraguay (2024-2026), consolidação
departamental e proxy local

\begin{longtable}[]{@{}ll@{}}
\toprule\noalign{}
Serviço & Disponibilidade \\
\midrule\noalign{}
\endhead
\bottomrule\noalign{}
\endlastfoot
USF /\allowbreak  Posto de Saúde & sim \\
Hospital Regional & não \\
IPS (seguro social) & sim \\
Farmácia & sim \\
Distância ao hospital de referência & \textasciitilde98 km (Ciudad del
Este) \\
\end{longtable}

\textbf{Principais estabelecimentos:} USF local; referencia hospitalar
em Ciudad del Este

\textbf{Observação para imigrantes:} Atendimento primario local ou em
raio curto; casos de maior complexidade seguem para Ciudad del Este.
Cobertura privada continua recomendada para especialidades e urgências
de maior complexidade.
\newpage
\secaoDiagnostico{Nacunday}{\mapaConteudo{-26.0500}{-54.7000}}{Ñacunday oferece um dos maiores santuários de isolamento e produtividade
no sudeste de Alto Paraná. Sua força absoluta reside na combinação de
baixíssima densidade populacional com solos de elite e abundância
hídrica protegida pelo Parque Nacional. É o local ideal para quem busca
privacidade total e segurança social em um ambiente cercado pela
natureza, exigindo do ocupante autonomia logística e atenção às normas
jurídicas da faixa de fronteira.

\subsubsection{Por que sim}

\begin{itemize}
\tightlist
\item
  Densidade populacional entre as menores da Região Oriental
  (privacidade absoluta).
\item
  Solo de alta produtividade e abundância de águas puros (Salto
  Ñacunday).
\item
  Localização ``fora do radar'' estratégico militar e industrial
  nacional.
\item
  Potencial excepcional para autossuficiência alimentar.
\end{itemize}

\subsubsection{Por que não}

\begin{itemize}
\tightlist
\item
  Localização em faixa de fronteira fluvial sensível (Lei 2532/\allowbreak 05).
\item
  Infraestrutura de saúde local limitada a atendimentos primários.
\item
  Dependência de estradas vicinais de terra para acesso interno.
\end{itemize}

\subsubsection{Recomendação
Técnica}

Altamente indicado para estabelecimentos de autossuficiência de médio
porte ou residências rurais de isolamento tático. Requer autonomia em
logística básica, saúde e segurança patrimonial. O GSS de 6.4 reflete a
segurança do isolamento em contraste com a limitação institucional
local. Referencia atual de score: GSS 6.4 em 2026-03-05.

\subsubsection{}

\begin{itemize}
\tightlist
\item
  \begin{enumerate}
  \def\labelenumi{\arabic{enumi})}
  \tightlist
  \item
    Dados censitários validados (INE\footnote{Instituto Nacional de Estadística (Instituto Nacional de Estatística), órgão oficial de dados demográficos e censitários.} 2022).
  \end{enumerate}
\item
  \begin{enumerate}
  \def\labelenumi{\arabic{enumi})}
  \setcounter{enumi}{1}
  \tightlist
  \item
    Relatórios de monitoramento ambiental do Parque Nacional Ñacunday
    (MADES).
  \end{enumerate}
\item
  \begin{enumerate}
  \def\labelenumi{\arabic{enumi})}
  \setcounter{enumi}{2}
  \tightlist
  \item
    Mapa de solos e aptidão agrícola de Itapúa/\allowbreak Alto Paraná (MAG).
  \end{enumerate}
\item
  \begin{enumerate}
  \def\labelenumi{\arabic{enumi})}
  \setcounter{enumi}{3}
  \tightlist
  \item
    Registro de segurança pública departamental.
  \end{enumerate}
\end{itemize}}
\begin{table}[ht!]
\small \begin{tabular}{p{3.0cm}p{0.8cm}p{6.0cm}} \toprule
\textbf{Critério} & \textbf{Nota} & \textbf{Justificativa} \\ \midrule
A - Ameaças Estratégicas & 6.5 & Localização remota e protegida por parques naturais; excelente invisibilidade tática \\
B - Risco Social & 7.5 & Baixíssima densidade demográfica; ideal para isolamento social absoluto \\
C - Riscos Naturais & 6.5 & Geologia estável; nota penalizada por riscos de erosão e isolamento pluvial \\
D - Autossuficiência & 7.0 & Bons recursos naturais - solo fértil, gado e água abundante - mas limitado pela logística local \\
E - Institucional & 4.5 & Município estável seguro, mas com infraestrutura institucional mínima e desafios de vigilância de fronteira \\
\midrule \textbf{GSS FINAL} & \textbf{6.4} & \textbf{Moderadamente Seguro (Ideal para quem busca o máximo de isolamento demográfico em território cercado pela natureza).} \\ \bottomrule \end{tabular}
\end{table}
\subsubsection{Vulnerabilidades e
Mitigação}\label{vuln-Nacunday-vulnerabilidades-e-mitigauxe7uxe3o}

\begin{itemize}
\tightlist
\item
  \textbf{Isolamento de Saúde:} Distância significativa de hospitais
  cirúrgicos (Mitigação: manutenção de estoque médico básico e kit
  trauma local).
\item
  \textbf{Logística Interna:} Acesso terrestre precário em chuvas
  intensas (Mitigação: estoque de suprimentos para 60 dias e posse de
  veículo 4x4 robusto).
\item
  \textbf{Restrições Legais de Fronteira:} Exige conformidade jurídica
  para investidores estrangeiros em zonas rurais
  limítrofes.
\end{itemize}
\subsubsection{Dossiê de Campo}\label{dossie-Nacunday-dossiuxea-de-campo}
\begin{description}[style=nextline,leftmargin=0.5cm,font=\bfseries]
\item[Coordenadas]  26°03'S 54°42'W.
\item[Topografia]  Relevo ondulado a acidentado, com presença de matas nativas preservadas e o imponente Salto Ñacunday. Margeado pelo Rio Paraná (fronteira fluvial com Argentina). Altitude \textasciitilde 180m. Área vasta de \textasciitilde 865 km². Geologicamente estável.
\item[Alvos Estratégicos]  Parque Nacional Ñacunday (Recurso ambiental e hídrico estratégico); Áreas de produção agrícola mecanizada; Proximidade tática com o Rio Paraná. Ausência de alvos militares de grande porte, oferecendo excelente invisibilidade estratégica.
\item[Fallout]  Ventos predominantes do Norte (NE) e Leste. Baixo risco de fallout direto.
\item[População]  6.468 habitantes (Censo 2022 Final). Perfil demográfico essencialmente rural (90\%) com forte cultura de preservação e agronegócio.
\item[Densidade]  \textasciitilde 7,5 hab/\allowbreak km² (Densidade extremamente baixa, oferecendo isolamento tático superior e privacidade estratégica absoluta fora dos pequenos núcleos povoados).
\item[Vias de Acesso]  Acesso via ramais de terra e conexões pavimentadas secundárias a partir da Ruta PY07. Localização isolada no sudeste do departamento, funcionando como um enclave seguro e discreto. Estradas internas de terra podem ser críticas em chuvas intensas (média 1.700 mm/\allowbreak ano).
\item[Serviços]  Unidade de Saúde da Família (USF) local. Dependência total de Santa Rita ou Ciudad del Este para saúde de alta complexidade e logística comercial avançada.
\item[Custo de Vida]  Muito baixo; economia baseada na produção primária e agricultura familiar.
\item[Preço da Terra]  US\$ 5.000-8.500/\allowbreak ha (áreas rurais agrícolas); US\$ 2.000-3.500/\allowbreak ha (áreas de campo natural/\allowbreak mata).
\item[Hidrologia]  Baixo risco de inundações sistêmicas devido à topografia elevada em relação ao Rio Paraná. Presença de arroios perenes de alta qualidade hídrica (ex: Rio Ñacunday). Risco de erosão em áreas de encosta mecanizadas sem cobertura.
\item[Clima]  Subtropical Úmido com influência da mata atlântica interiorana. Verões intensos e invernos com geadas ocasionais.
\item[Energia]  Rede nacional (Itaipu); instável em áreas rurais remotas devido ao isolamento geográfico.
\item[Água]  Abundante via Rio Paraná, nascentes protegidas e poços artesianos; excelente qualidade hídrica subterrânea.
\item[Qualidade do Solo]  Latossolo Vermelho de Alta Performance; solo profundo e muito fértil nas áreas de transição florestal. Excelente para soja, milho e reflorestamento.
\item[Resources Locais]  Grande produção de grãos, gado de corte e agricultura de subsistência. Elevadíssimo potencial para autossuficiência alimentar básica em nível de propriedade devido à vasta área produtiva.
\item[Segurança]  Zona rural pacífica e estável. Baixíssima criminalidade urbana. Vigilância necessária contra atividades ilícitas em fronteira fluvial. Coesão social baseada na cultura produtora tradicional. Baixa visibilidade para conflitos de massa.
\item[Leis Local: Município consolidado e polo de agronegócio e ecoturismo. Restrição de Fronteira (Lei 2532/\allowbreak 05)]  Aplicável a terras rurais na faixa de 50km da fronteira argentina.
\end{description}
\paragraph{Solo (SoilGrids 2.0, média ponderada 0--30
cm)}

\begin{longtable}[]{@{}ll@{}}
\toprule\noalign{}
Parâmetro & Valor \\
\midrule\noalign{}
\endhead
\bottomrule\noalign{}
\endlastfoot
pH (H$_2$O) & 5.3 \\
Carbono orgânico (SOC) & 19.57 g/\allowbreak kg \\
Argila & 44.45 \% \\
Areia & 18.37 \% \\
Silte (calc.) & 37.2 \% \\
Densidade aparente & 1.23 g/\allowbreak cm³ \\
\textbf{Aptidão agrícola} & \textbf{Média} \\
\end{longtable}

Fonte: ISRIC SoilGrids 2.0 via WCS (acesso em 2026-03-21). Coords:
-26.05°, -54.7°. Média ponderada camadas 0-5, 5-15, 15-30 cm.
\subsubsection{Indicadores Sociais}

\textbf{Fontes:} PNUD 2020, INE\footnote{Instituto Nacional de Estadística (Instituto Nacional de Estatística), órgão oficial de dados demográficos e censitários.} Censo 2022, INE\footnote{Instituto Nacional de Estadística (Instituto Nacional de Estatística), órgão oficial de dados demográficos e censitários.} EPHC 2023, Ministerio
Público 2024\\
\textbf{Nota:} Valores marcados como \emph{(dept.)} referem-se ao
departamento de Alto Paraná; valores \emph{(dist.)} são específicos
deste distrito. Dados marcados \emph{(est.)} são estimativas.

\begin{longtable}[]{@{}
  >{\raggedright\arraybackslash}p{(\linewidth - 4\tabcolsep) * \real{0.4231}}
  >{\raggedright\arraybackslash}p{(\linewidth - 4\tabcolsep) * \real{0.2692}}
  >{\raggedright\arraybackslash}p{(\linewidth - 4\tabcolsep) * \real{0.3077}}@{}}
\toprule\noalign{}
\begin{minipage}[b]{\linewidth}\raggedright
Indicador
\end{minipage} & \begin{minipage}[b]{\linewidth}\raggedright
Valor
\end{minipage} & \begin{minipage}[b]{\linewidth}\raggedright
Âmbito
\end{minipage} \\
\midrule\noalign{}
\endhead
\bottomrule\noalign{}
\endlastfoot
IDH (2020) & 0,752 & dept. \\
Ranking IDH nacional & 3/\allowbreak 18 & dept. \\
Esperança de vida & 75,1 anos & dept. \\
Escolaridade média & 7.5 anos & dist. \\
RNB per capita (USD PPA) & 13.500 & dept. \\
Pobreza monetária (\%) & 16,3\% & dept. \\
Pobreza extrema (\%) & 3,1\% & dept. \\
Índice de Gini & 0,428 & dept. \\
Acesso a água potável (\%) & 88,7\% & dept. \\
Acesso a saneamento (\%) & 87,2\% & dept. \\
Taxa de homicídios (est., /\allowbreak 100k hab) & \textasciitilde8--12 (estimativa,
acima da média nacional) & dept. \\
Índice de segurança & baixo & dept. \\
População (Censo 2022) & 6.468 hab. & dist. \\
Idade mediana & 24.0 anos & dist. \\
Taxa de fecundidade & N/\allowbreak D & dist. \\
\end{longtable}
\subsubsection{Combustível}

\textbf{Referência:} PETROPAR /\allowbreak  postos locais (2024) \textbf{Tipo de
localidade:} Interior

\begin{longtable}[]{@{}lll@{}}
\toprule\noalign{}
Combustível & USD/\allowbreak litro & Gs/\allowbreak litro (aprox.) \\
\midrule\noalign{}
\endhead
\bottomrule\noalign{}
\endlastfoot
Gasolina 93 oct & 0.97 & 7,178 \\
Gasolina 97 oct (premium) & 1.07 & 7,918 \\
Diesel & 0.90 & 6,660 \\
\end{longtable}

\begin{quote}
Preços podem variar ±5\% conforme posto e sazonalidade. Chaco e interior
remoto apresentam maior variação.
\end{quote}

\subsubsection{Cobertura Celular}

\textbf{Fonte:} CONATEL PY /\allowbreak  operadoras (2024)

\begin{longtable}[]{@{}ll@{}}
\toprule\noalign{}
Parâmetro & Valor \\
\midrule\noalign{}
\endhead
\bottomrule\noalign{}
\endlastfoot
Cobertura 4G população (dept.) & 95\% \\
Cobertura 4G área rural & 88\% \\
Melhor operadora & Tigo/\allowbreak Personal \\
Qualidade rural & boa \\
\end{longtable}

\begin{quote}
Para áreas rurais fora do núcleo urbano, recomenda-se chip Tigo como
principal e Personal como backup.
\end{quote}

\subsubsection{Internet}

\textbf{Fonte:} CONATEL /\allowbreak  Speedtest Ookla (2024)

\begin{longtable}[]{@{}ll@{}}
\toprule\noalign{}
Parâmetro & Valor \\
\midrule\noalign{}
\endhead
\bottomrule\noalign{}
\endlastfoot
Velocidade média download & 85 Mbps \\
Domicílios com internet (dept.) & 87\% \\
Tecnologia predominante & fibra/\allowbreak cabo \\
Opção rural & Starlink disponível (\textasciitilde USD 44/\allowbreak mês) \\
\end{longtable}

\subsubsection{Mercado Imobiliário e Terra
Rural}

\textbf{Fonte:} INDERT /\allowbreak  Clasificados.com.py (2024)

\begin{longtable}[]{@{}ll@{}}
\toprule\noalign{}
Tipo & Referência \\
\midrule\noalign{}
\endhead
\bottomrule\noalign{}
\endlastfoot
Terra agrícola alta prod. (USD/\allowbreak ha) & 14,800 \\
Imóvel urbano (USD/\allowbreak m²) & 1,400 \\
Aluguel 2 quartos (USD/\allowbreak mês) & 600 \\
\end{longtable}

\begin{quote}
Valores de referência departamental. Localidades menores podem ter
preços 20--40\% abaixo da capital departamental. \#\#\# Saúde
\end{quote}

\textbf{Fonte:} MSPBS /\allowbreak  IPS Paraguay (2024-2026), consolidação
departamental e proxy local

\begin{longtable}[]{@{}ll@{}}
\toprule\noalign{}
Serviço & Disponibilidade \\
\midrule\noalign{}
\endhead
\bottomrule\noalign{}
\endlastfoot
USF /\allowbreak  Posto de Saúde & sim \\
Hospital Regional & não \\
IPS (seguro social) & sim \\
Farmácia & sim \\
Distância ao hospital de referência & \textasciitilde65 km (Ciudad del
Este) \\
\end{longtable}

\textbf{Principais estabelecimentos:} USF local; referencia hospitalar
em Ciudad del Este

\textbf{Observação para imigrantes:} Atendimento primario local ou em
raio curto; casos de maior complexidade seguem para Ciudad del Este.
Cobertura privada continua recomendada para especialidades e urgências
de maior complexidade.
\newpage
\secaoDiagnostico{Naranjal}{\mapaConteudo{-25.9666}{-55.1833}}{Naranjal é um enclave de elite no coração produtivo de Alto Paraná.
Oferece uma combinação sólida de segurança institucional, abundância
excepcional de recursos alimentares e hídricos e uma densidade
demográfica baixíssima. Sua força absoluta reside na organização
comunitária via cooperativas e na qualidade mundial de seus solos. É o
local ideal para projetos de autossuficiência tecnificada que busquem
privacidade tática e segurança social, ciente da visibilidade
estratégica de um polo industrial desse porte.

\subsubsection{Por que sim}

\begin{itemize}
\tightlist
\item
  Solo de altíssima produtividade e ambiente social extremamente seguro
  e ordeiro.
\item
  Baixíssima densidade populacional (privacidade absoluta).
\item
  Conectividade rápida via Ruta PY06 para os principais centros de
  Itapúa e Alto Paraná.
\item
  Sem restrições da Lei de Fronteira.
\end{itemize}

\subsubsection{Por que não}

\begin{itemize}
\tightlist
\item
  Alvo estratégico logístico e alimentar em cenários de conflito
  nacional.
\item
  Valor da terra agrícola entre os mais altos do país.
\item
  Dependência tática de insumos externos para manutenção da escala
  industrial.
\end{itemize}

\subsubsection{Recomendação
Técnica}

Altamente indicado para estabelecimentos de autossuficiência familiar de
alto padrão ou investimentos agroindustriais. O GSS de 6.8 reflete a
excepcional solidez institucional e a riqueza de recursos da localidade.
Referencia atual de score: GSS 6.8 em 2026-03-05.

\subsubsection{}

\begin{itemize}
\tightlist
\item
  \begin{enumerate}
  \def\labelenumi{\arabic{enumi})}
  \tightlist
  \item
    Dados censitários validados (INE\footnote{Instituto Nacional de Estadística (Instituto Nacional de Estatística), órgão oficial de dados demográficos e censitários.} 2022).
  \end{enumerate}
\item
  \begin{enumerate}
  \def\labelenumi{\arabic{enumi})}
  \setcounter{enumi}{1}
  \tightlist
  \item
    Relatórios de produção e sustentabilidade da COPRONAR.
  \end{enumerate}
\item
  \begin{enumerate}
  \def\labelenumi{\arabic{enumi})}
  \setcounter{enumi}{2}
  \tightlist
  \item
    Mapa de solos e bacias hidrográficas de Alto Paraná (MAG).
  \end{enumerate}
\item
  \begin{enumerate}
  \def\labelenumi{\arabic{enumi})}
  \setcounter{enumi}{3}
  \tightlist
  \item
    Registro de segurança pública e governança cooperativa regional.
  \end{enumerate}
\end{itemize}}
\begin{table}[ht!]
\small \begin{tabular}{p{3.0cm}p{0.8cm}p{6.0cm}} \toprule
\textbf{Critério} & \textbf{Nota} & \textbf{Justificativa} \\ \midrule
A - Ameaças Estratégicas & 5.5 & Hub agroindustrial tático; alvo estratégico logístico de interesse nacional em crises \\
B - Risco Social & 7.0 & População pequena e densidade extremamente baixa; ideal para isolamento social \\
C - Riscos Naturais & 7.0 & Geologia muito estável e baixo risco de desastres naturais graves \\
D - Autossuficiência & 8.5 & Recursos alimentares e hídricos excepcionais: sede da COPRONAR e solos de elite \\
E - Institucional & 7.5 & Forte segurança institucional e social sustentada pelo setor privado cooperativista \\
\midrule \textbf{GSS FINAL} & \textbf{6.8} & \textbf{Seguro (Localidade altamente recomendada para quem busca isolamento rural produtivo com segurança superior).} \\ \bottomrule \end{tabular}
\end{table}
\subsubsection{Vulnerabilidades e
Mitigação}\label{vuln-Naranjal-vulnerabilidades-e-mitigauxe7uxe3o}

\begin{itemize}
\tightlist
\item
  \textbf{Visibilidade Tática de Suprimentos:} O polo agroindustrial é
  um alvo logístico óbvio em crises nacionais de fome (Mitigação:
  residência rural afastada 5-10 km do núcleo urbano industrial e
  integração com a rede de defesa cooperativa).
\item
  \textbf{Dependência de Serviços Externos:} Saúde especializada em
  cidades maiores (Mitigação: manutenção de estoque médico básico e kit
  trauma local).
\item
  \textbf{Especialização de Safra:} Dependência de insumos externos para
  alta produtividade (Mitigação: diversificação produtiva de pequena
  escala para consumo imediato).
\end{itemize}
\subsubsection{Dossiê de Campo}\label{dossie-Naranjal-dossiuxea-de-campo}
\begin{description}[style=nextline,leftmargin=0.5cm,font=\bfseries]
\item[Coordenadas]  25°58'S 55°11'W.
\item[Topografia]  Relevo predominantemente plano a suavemente ondulado, solo de altíssima produtividade mecanizado em 95\% da área. Altitude \textasciitilde 240m. Área de \textasciitilde 802 km². Geologicamente muito estável, risco sísmico desprezível.
\item[Alvos Estratégicos]  Cooperativa de Produção Agroindustrial Naranjal (COPRONAR) - Infraestrutura crítica de armazenamento e processamento de grãos; Grandes plantas de silos; Polo logístico regional do sul de Alto Paraná.
\item[Fallout]  Ventos predominantes do Norte (NE) e Leste. Baixo risco de fallout direto.
\item[População]  6.608 habitantes (Censo 2022 Final). População com forte herança de imigrantes brasileiros, alemães e italianos, com cultura de cooperação e eficiência produtiva.
\item[Densidade]  \textasciitilde 8,2 hab/\allowbreak km² (Densidade extremamente baixa, oferecendo isolamento tático superior e privacidade estratégica absoluta).
\item[Vias de Acesso]  Acesso asfáltico de alta qualidade ligando à Ruta PY06. Localização que permite rápido escoamento comercial, mas mantém a localidade fora dos eixos de trânsito intenso metropolitano.
\item[Serviços]  Infraestrutura de serviços de alta qualidade para o porte do distrito. Clínicas modernas, apoio técnico avançado da cooperativa e centros educacionais. Dependência de Santa Rita (35 km) para alta complexidade médica.
\item[Custo de Vida]  Médio; economia dolarizada pelo agronegócio, mas com acesso direto a produtos básicos a preços de origem.
\item[Preço da Terra]  US\$ 10.000-15.000/\allowbreak ha (terras mecanizadas de elite - "Terra Roxa").
\item[Hidrologia]  Baixo risco de inundações sistêmicas devido à gestão hídrica agrícola eficiente e topografia favorável. Presença de arroios perenes utilizados para irrigação e pecuária.
\item[Clima]  Subtropical Úmido. Verões quentes e invernos com geadas que favorecem a safra de trigo.
\item[Energia]  Rede nacional (Itaipu); boa estabilidade devido à importância industrial. Grande parte dos silos e indústrias possui sistemas de geração redundante (biomassa/\allowbreak térmica).
\item[Água]  Abundante via poços artesianos profundos e lençol freático produtivo (Sistema Aquífero Guarani).
\item[Qualidade do Solo]  Latossolo Vermelho de Elite; solo profundo e de altíssima fertilidade natural, mantido por técnicas de agricultura de precisão.
\item[Resources Locais]  Grande produção de soja, milho, trigo, suínos e aves. Elevadíssimo potencial para autossuficiência calórica e hídrica em nível de propriedade ou cooperativa.
\item[Segurança]  Uma das zonas mais seguras do Paraguai. Baixíssima criminalidade urbana violenta. Forte coesão social comunitária e segurança patrimonial privada altamente eficiente nas fazendas e indústrias. Baixa visibilidade para conflitos externos.
\item[Leis Local]  Município consolidado e polo de referência tecnológica ("Cidade Cooperativa"). \\ \textbf{Livre de Restrição} de Fronteira. \\ \textit{Aquisição de terra rural proibida para estrangeiros limítrofes.}
\end{description}
\paragraph{Solo (SoilGrids 2.0, média ponderada 0--30
cm)}

\begin{longtable}[]{@{}ll@{}}
\toprule\noalign{}
Parâmetro & Valor \\
\midrule\noalign{}
\endhead
\bottomrule\noalign{}
\endlastfoot
pH (H$_2$O) & 5.4 \\
Carbono orgânico (SOC) & 21.68 g/\allowbreak kg \\
Argila & 46.93 \% \\
Areia & 22.86 \% \\
Silte (calc.) & 30.2 \% \\
Densidade aparente & 1.23 g/\allowbreak cm³ \\
\textbf{Aptidão agrícola} & \textbf{Média} \\
\end{longtable}

Fonte: ISRIC SoilGrids 2.0 via WCS (acesso em 2026-03-21). Coords:
-25.9667°, -55.1833°. Média ponderada camadas 0-5, 5-15, 15-30 cm.
\subsubsection{Indicadores Sociais}

\textbf{Fontes:} PNUD 2020, INE\footnote{Instituto Nacional de Estadística (Instituto Nacional de Estatística), órgão oficial de dados demográficos e censitários.} Censo 2022, INE\footnote{Instituto Nacional de Estadística (Instituto Nacional de Estatística), órgão oficial de dados demográficos e censitários.} EPHC 2023, Ministerio
Público 2024\\
\textbf{Nota:} Valores marcados como \emph{(dept.)} referem-se ao
departamento de Alto Paraná; valores \emph{(dist.)} são específicos
deste distrito. Dados marcados \emph{(est.)} são estimativas.

\begin{longtable}[]{@{}
  >{\raggedright\arraybackslash}p{(\linewidth - 4\tabcolsep) * \real{0.4231}}
  >{\raggedright\arraybackslash}p{(\linewidth - 4\tabcolsep) * \real{0.2692}}
  >{\raggedright\arraybackslash}p{(\linewidth - 4\tabcolsep) * \real{0.3077}}@{}}
\toprule\noalign{}
\begin{minipage}[b]{\linewidth}\raggedright
Indicador
\end{minipage} & \begin{minipage}[b]{\linewidth}\raggedright
Valor
\end{minipage} & \begin{minipage}[b]{\linewidth}\raggedright
Âmbito
\end{minipage} \\
\midrule\noalign{}
\endhead
\bottomrule\noalign{}
\endlastfoot
IDH (2020) & 0,752 & dept. \\
Ranking IDH nacional & 3/\allowbreak 18 & dept. \\
Esperança de vida & 75,1 anos & dept. \\
Escolaridade média & 8.4 anos & dist. \\
RNB per capita (USD PPA) & 13.500 & dept. \\
Pobreza monetária (\%) & 16,3\% & dept. \\
Pobreza extrema (\%) & 3,1\% & dept. \\
Índice de Gini & 0,428 & dept. \\
Acesso a água potável (\%) & 88,7\% & dept. \\
Acesso a saneamento (\%) & 87,2\% & dept. \\
Taxa de homicídios (est., /\allowbreak 100k hab) & \textasciitilde8--12 (estimativa,
acima da média nacional) & dept. \\
Índice de segurança & baixo & dept. \\
População (Censo 2022) & 6.608 hab. & dist. \\
Idade mediana & 29.0 anos & dist. \\
Taxa de fecundidade & N/\allowbreak D & dist. \\
\end{longtable}
\subsubsection{Combustível}

\textbf{Referência:} PETROPAR /\allowbreak  postos locais (2024) \textbf{Tipo de
localidade:} Interior

\begin{longtable}[]{@{}lll@{}}
\toprule\noalign{}
Combustível & USD/\allowbreak litro & Gs/\allowbreak litro (aprox.) \\
\midrule\noalign{}
\endhead
\bottomrule\noalign{}
\endlastfoot
Gasolina 93 oct & 0.97 & 7,178 \\
Gasolina 97 oct (premium) & 1.07 & 7,918 \\
Diesel & 0.90 & 6,660 \\
\end{longtable}

\begin{quote}
Preços podem variar ±5\% conforme posto e sazonalidade. Chaco e interior
remoto apresentam maior variação.
\end{quote}

\subsubsection{Cobertura Celular}

\textbf{Fonte:} CONATEL PY /\allowbreak  operadoras (2024)

\begin{longtable}[]{@{}ll@{}}
\toprule\noalign{}
Parâmetro & Valor \\
\midrule\noalign{}
\endhead
\bottomrule\noalign{}
\endlastfoot
Cobertura 4G população (dept.) & 95\% \\
Cobertura 4G área rural & 88\% \\
Melhor operadora & Tigo/\allowbreak Personal \\
Qualidade rural & boa \\
\end{longtable}

\begin{quote}
Para áreas rurais fora do núcleo urbano, recomenda-se chip Tigo como
principal e Personal como backup.
\end{quote}

\subsubsection{Internet}

\textbf{Fonte:} CONATEL /\allowbreak  Speedtest Ookla (2024)

\begin{longtable}[]{@{}ll@{}}
\toprule\noalign{}
Parâmetro & Valor \\
\midrule\noalign{}
\endhead
\bottomrule\noalign{}
\endlastfoot
Velocidade média download & 85 Mbps \\
Domicílios com internet (dept.) & 87\% \\
Tecnologia predominante & fibra/\allowbreak cabo \\
Opção rural & Starlink disponível (\textasciitilde USD 44/\allowbreak mês) \\
\end{longtable}

\subsubsection{Mercado Imobiliário e Terra
Rural}

\textbf{Fonte:} INDERT /\allowbreak  Clasificados.com.py (2024)

\begin{longtable}[]{@{}ll@{}}
\toprule\noalign{}
Tipo & Referência \\
\midrule\noalign{}
\endhead
\bottomrule\noalign{}
\endlastfoot
Terra agrícola alta prod. (USD/\allowbreak ha) & 14,800 \\
Imóvel urbano (USD/\allowbreak m²) & 1,400 \\
Aluguel 2 quartos (USD/\allowbreak mês) & 600 \\
\end{longtable}

\begin{quote}
Valores de referência departamental. Localidades menores podem ter
preços 20--40\% abaixo da capital departamental. \#\#\# Saúde
\end{quote}

\textbf{Fonte:} MSPBS /\allowbreak  IPS Paraguay (2024-2026), consolidação
departamental e proxy local

\begin{longtable}[]{@{}ll@{}}
\toprule\noalign{}
Serviço & Disponibilidade \\
\midrule\noalign{}
\endhead
\bottomrule\noalign{}
\endlastfoot
USF /\allowbreak  Posto de Saúde & sim \\
Hospital Regional & não \\
IPS (seguro social) & sim \\
Farmácia & sim \\
Distância ao hospital de referência & \textasciitilde96 km (Ciudad del
Este) \\
\end{longtable}

\textbf{Principais estabelecimentos:} USF local; referencia hospitalar
em Ciudad del Este

\textbf{Observação para imigrantes:} Atendimento primario local ou em
raio curto; casos de maior complexidade seguem para Ciudad del Este.
Cobertura privada continua recomendada para especialidades e urgências
de maior complexidade.
\newpage
\secaoDiagnostico{Presidente Franco}{\mapaConteudo{-25.5666}{-54.6166}}{Presidente Franco (Alto Parana) apresenta perfil operacional consolidado
para comparacao territorial, com leitura conjunta de infraestrutura,
risco hidroclimatico e capacidade de continuidade de servicos
essenciais.

\subsubsection{Por que sim}

\begin{itemize}
\tightlist
\item
  Base institucional de referencia disponivel e rastreavel.
\item
  Estrutura comparativa (A-E/\allowbreak GSS) consistente para decisao relativa.
\item
  Plano de mitigacao acionavel ja definido no dossie.
\end{itemize}

\subsubsection{Por que não}

\begin{itemize}
\tightlist
\item
  Serie oficial distrital granular ainda pode apresentar lacunas
  temporais.
\item
  Sensibilidade do score exige monitoramento continuo de risco e
  conectividade.
\item
  Decisao final depende de validacao de campo e janela temporal recente.
\end{itemize}

\subsubsection{Pre-condicoes Minimas de
Decisao}

\begin{itemize}
\tightlist
\item
  Atualizar indicadores criticos em ciclo trimestral.
\item
  Confirmar trafegabilidade e contingencia hidrica/\allowbreak energetica local.
\item
  Validar checklist de resiliencia domiciliar/\allowbreak logistica antes de decisao
  final.
\end{itemize}

\subsubsection{Recomendação
Técnica}

Uso recomendado como candidato em analise multicriterio, condicionado ao
cumprimento das pre-condicoes acima. Referencia atual de score: GSS 7.7
em 2026-03-05; risco predominante: baixo.

\subsubsection{}

\begin{itemize}
\tightlist
\item
  \begin{enumerate}
  \def\labelenumi{\arabic{enumi})}
  \tightlist
  \item
    Delimitacao territorial validada por georreferencia institucional
    (Fonte: acesso: 2026-03-05).
  \end{enumerate}
\item
  \begin{enumerate}
  \def\labelenumi{\arabic{enumi})}
  \setcounter{enumi}{1}
  \tightlist
  \item
    População municipal/\allowbreak distrital referenciada por base censitaria
    nacional (Fonte: acesso: 2026-03-05).
  \end{enumerate}
\item
  \begin{enumerate}
  \def\labelenumi{\arabic{enumi})}
  \setcounter{enumi}{2}
  \tightlist
  \item
    Comparabilidade intra-departamental apoiada em indicadores
    distritais oficiais (Fonte: acesso: 2026-03-05).
  \end{enumerate}
\item
  \begin{enumerate}
  \def\labelenumi{\arabic{enumi})}
  \setcounter{enumi}{3}
  \tightlist
  \item
    Estado macro de conectividade viaria ancorado em informacoes de
    rodovias (Fonte: acesso: 2026-03-05).
  \end{enumerate}
\item
  \begin{enumerate}
  \def\labelenumi{\arabic{enumi})}
  \setcounter{enumi}{4}
  \tightlist
  \item
    Exposicao hidrometeorologica monitorada por avisos oficiais (Fonte:
    acesso: 2026-03-05).
  \end{enumerate}
\item
  \begin{enumerate}
  \def\labelenumi{\arabic{enumi})}
  \setcounter{enumi}{5}
  \tightlist
  \item
    Historico de contexto meteorologico apoiado em publicacoes tecnicas
    (Fonte: acesso: 2026-03-05).
  \end{enumerate}
\item
  \begin{enumerate}
  \def\labelenumi{\arabic{enumi})}
  \setcounter{enumi}{6}
  \tightlist
  \item
    Capacidade institucional de resposta a eventos apoiada em acoes da
    SEN\footnote{Secretaría de Emergencia Nacional (Secretaria de Emergência Nacional), órgão governamental de resposta a desastres.} (Fonte: acesso: 2026-03-05).
  \end{enumerate}
\item
  \begin{enumerate}
  \def\labelenumi{\arabic{enumi})}
  \setcounter{enumi}{7}
  \tightlist
  \item
    Infraestrutura energetica analisada via operador nacional (Fonte:
    acesso: 2026-03-05).
  \end{enumerate}
\item
  \begin{enumerate}
  \def\labelenumi{\arabic{enumi})}
  \setcounter{enumi}{8}
  \tightlist
  \item
    Referencias de obras e logistica nacional verificadas no MOPC
    (Fonte: acesso: 2026-03-05).
  \end{enumerate}
\item
  \begin{enumerate}
  \def\labelenumi{\arabic{enumi})}
  \setcounter{enumi}{9}
  \tightlist
  \item
    Contexto sociopolitico institucional referenciado por autoridade
    eleitoral (Fonte: acesso: 2026-03-05).
  \end{enumerate}
\item
  \begin{enumerate}
  \def\labelenumi{\arabic{enumi})}
  \setcounter{enumi}{10}
  \tightlist
  \item
    Camada cartografica de apoio eleitoral/\allowbreak territorial consultada
    (Fonte: acesso: 2026-03-05).
  \end{enumerate}
\item
  \begin{enumerate}
  \def\labelenumi{\arabic{enumi})}
  \setcounter{enumi}{11}
  \tightlist
  \item
    Dados publicos complementares para verificacao cruzada (Fonte:
    acesso: 2026-03-05).
  \end{enumerate}
\end{itemize}}
\begin{table}[ht!]
\small \begin{tabular}{p{3.0cm}p{0.8cm}p{6.0cm}} \toprule
\textbf{Critério} & \textbf{Nota} & \textbf{Justificativa} \\ \midrule
A - Ameaças Estratégicas & 8.7 & - \\
B - Risco Social & 6.7 & - \\
C - Riscos Naturais & 8.4 & - \\
D - Autossuficiência & 6.9 & - \\
E - Institucional & 7.6 & - \\
\midrule \textbf{GSS FINAL} & \textbf{7.7} & \textbf{Seguro (bom potencial com preparacao).} \\ \bottomrule \end{tabular}
\end{table}
\subsubsection{Vulnerabilidades e
Mitigação}\label{vuln-Presidente_Franco-vulnerabilidades-e-mitigauxe7uxe3o}

\begin{itemize}
\tightlist
\item
  Exposicao hidrometeorologica controlada (\textless45/\allowbreak 100).
\item
  Conectividade viaria favoravel (\textgreater=70/\allowbreak 100).
\item
  Estoque de contingencia para 14 dias (agua/\allowbreak alimentos/\allowbreak combustivel),
  priorizado quando risco hidro=20/\allowbreak 100.
\item
  Duplicar rotas/\allowbreak logistica quando conectividade viaria=78/\allowbreak 100 for
  \textless70; manter rota secundaria mapeada e testada trimestralmente.
\item
  Plano de energia resiliente com autonomia minima de 72h
  (geracao/\allowbreak backup), revisao semestral.
\end{itemize}
\subsubsection{Dossiê de Campo}\label{dossie-Presidente_Franco-dossiuxea-de-campo}
\begin{description}[style=nextline,leftmargin=0.5cm,font=\bfseries]
\item[Coordenadas]  25°34'S 54°37'W (geocodificado via Nominatim).
\end{description}
\paragraph{Solo (SoilGrids 2.0, média ponderada 0--30
cm)}

\begin{longtable}[]{@{}ll@{}}
\toprule\noalign{}
Parâmetro & Valor \\
\midrule\noalign{}
\endhead
\bottomrule\noalign{}
\endlastfoot
pH (H$_2$O) & 5.3 \\
Carbono orgânico (SOC) & 25.01 g/\allowbreak kg \\
Argila & 45.81 \% \\
Areia & 16.61 \% \\
Silte (calc.) & 37.6 \% \\
Densidade aparente & 1.18 g/\allowbreak cm³ \\
\textbf{Aptidão agrícola} & \textbf{Média} \\
\end{longtable}

Fonte: ISRIC SoilGrids 2.0 via WCS (acesso em 2026-03-21). Coords:
-25.5667°, -54.6167°. Média ponderada camadas 0-5, 5-15, 15-30 cm.
\subsubsection{Indicadores Sociais}

\textbf{Fontes:} PNUD 2020, INE\footnote{Instituto Nacional de Estadística (Instituto Nacional de Estatística), órgão oficial de dados demográficos e censitários.} Censo 2022, INE\footnote{Instituto Nacional de Estadística (Instituto Nacional de Estatística), órgão oficial de dados demográficos e censitários.} EPHC 2023, Ministerio
Público 2024\\
\textbf{Nota:} Valores marcados como \emph{(dept.)} referem-se ao
departamento de Alto Paraná; valores \emph{(dist.)} são específicos
deste distrito. Dados marcados \emph{(est.)} são estimativas.

\begin{longtable}[]{@{}
  >{\raggedright\arraybackslash}p{(\linewidth - 4\tabcolsep) * \real{0.4231}}
  >{\raggedright\arraybackslash}p{(\linewidth - 4\tabcolsep) * \real{0.2692}}
  >{\raggedright\arraybackslash}p{(\linewidth - 4\tabcolsep) * \real{0.3077}}@{}}
\toprule\noalign{}
\begin{minipage}[b]{\linewidth}\raggedright
Indicador
\end{minipage} & \begin{minipage}[b]{\linewidth}\raggedright
Valor
\end{minipage} & \begin{minipage}[b]{\linewidth}\raggedright
Âmbito
\end{minipage} \\
\midrule\noalign{}
\endhead
\bottomrule\noalign{}
\endlastfoot
IDH (2020) & 0,752 & dept. \\
Ranking IDH nacional & 3/\allowbreak 18 & dept. \\
Esperança de vida & 75,1 anos & dept. \\
Escolaridade média & 10.0 anos & dist. \\
RNB per capita (USD PPA) & 13.500 & dept. \\
Pobreza monetária (\%) & 16,3\% & dept. \\
Pobreza extrema (\%) & 3,1\% & dept. \\
Índice de Gini & 0,428 & dept. \\
Acesso a água potável (\%) & 88,7\% & dept. \\
Acesso a saneamento (\%) & 87,2\% & dept. \\
Taxa de homicídios (est., /\allowbreak 100k hab) & \textasciitilde8--12 (estimativa,
acima da média nacional) & dept. \\
Índice de segurança & baixo & dept. \\
População (Censo 2022) & 88.744 hab. & dist. \\
Idade mediana & 28.0 anos & dist. \\
Taxa de fecundidade & N/\allowbreak D & dist. \\
\end{longtable}
\subsubsection{Combustível}

\textbf{Referência:} PETROPAR /\allowbreak  postos locais (2024) \textbf{Tipo de
localidade:} Interior

\begin{longtable}[]{@{}lll@{}}
\toprule\noalign{}
Combustível & USD/\allowbreak litro & Gs/\allowbreak litro (aprox.) \\
\midrule\noalign{}
\endhead
\bottomrule\noalign{}
\endlastfoot
Gasolina 93 oct & 0.97 & 7,178 \\
Gasolina 97 oct (premium) & 1.07 & 7,918 \\
Diesel & 0.90 & 6,660 \\
\end{longtable}

\begin{quote}
Preços podem variar ±5\% conforme posto e sazonalidade. Chaco e interior
remoto apresentam maior variação.
\end{quote}

\subsubsection{Cobertura Celular}

\textbf{Fonte:} CONATEL PY /\allowbreak  operadoras (2024)

\begin{longtable}[]{@{}ll@{}}
\toprule\noalign{}
Parâmetro & Valor \\
\midrule\noalign{}
\endhead
\bottomrule\noalign{}
\endlastfoot
Cobertura 4G população (dept.) & 95\% \\
Cobertura 4G área rural & 88\% \\
Melhor operadora & Tigo/\allowbreak Personal \\
Qualidade rural & boa \\
\end{longtable}

\begin{quote}
Para áreas rurais fora do núcleo urbano, recomenda-se chip Tigo como
principal e Personal como backup.
\end{quote}

\subsubsection{Internet}

\textbf{Fonte:} CONATEL /\allowbreak  Speedtest Ookla (2024)

\begin{longtable}[]{@{}ll@{}}
\toprule\noalign{}
Parâmetro & Valor \\
\midrule\noalign{}
\endhead
\bottomrule\noalign{}
\endlastfoot
Velocidade média download & 85 Mbps \\
Domicílios com internet (dept.) & 87\% \\
Tecnologia predominante & fibra/\allowbreak cabo \\
Opção rural & Starlink disponível (\textasciitilde USD 44/\allowbreak mês) \\
\end{longtable}

\subsubsection{Mercado Imobiliário e Terra
Rural}

\textbf{Fonte:} INDERT /\allowbreak  Clasificados.com.py (2024)

\begin{longtable}[]{@{}ll@{}}
\toprule\noalign{}
Tipo & Referência \\
\midrule\noalign{}
\endhead
\bottomrule\noalign{}
\endlastfoot
Terra agrícola alta prod. (USD/\allowbreak ha) & 14,800 \\
Imóvel urbano (USD/\allowbreak m²) & 1,400 \\
Aluguel 2 quartos (USD/\allowbreak mês) & 600 \\
\end{longtable}

\begin{quote}
Valores de referência departamental. Localidades menores podem ter
preços 20--40\% abaixo da capital departamental. \#\#\# Saúde
\end{quote}

\textbf{Fonte:} MSPBS /\allowbreak  IPS Paraguay (2024-2026), consolidação
departamental e proxy local

\begin{longtable}[]{@{}ll@{}}
\toprule\noalign{}
Serviço & Disponibilidade \\
\midrule\noalign{}
\endhead
\bottomrule\noalign{}
\endlastfoot
USF /\allowbreak  Posto de Saúde & sim \\
Hospital Regional & não \\
IPS (seguro social) & sim \\
Farmácia & sim \\
Distância ao hospital de referência & \textasciitilde7 km (Ciudad del
Este) \\
\end{longtable}

\textbf{Principais estabelecimentos:} USF local; referencia hospitalar
em Ciudad del Este

\textbf{Observação para imigrantes:} Atendimento primario local ou em
raio curto; casos de maior complexidade seguem para Ciudad del Este.
Cobertura privada continua recomendada para especialidades e urgências
de maior complexidade.
\newpage
\secaoDiagnostico{San Alberto}{\mapaConteudo{-24.9666}{-54.9000}}{San Alberto é o símbolo da potência produtiva do norte de Alto Paraná.
Oferece solos de produtividade mundial, abundância massiva de recursos
hídricos e alimentares, e uma segurança social sustentada por uma
comunidade produtora altamente organizada. Sua força absoluta reside na
resiliência de sua base produtiva e na baixa densidade populacional,
garantindo privacidade e isolamento tático superior. É o local ideal
para quem busca autonomia total em ambiente de alta performance,
transformando a riqueza de recursos básicos em uma armadura definitiva
de proteção contra instabilidades externas.

\subsubsection{Por que sim}

\begin{itemize}
\tightlist
\item
  Solos de altíssima produtividade mundial e abundância de recursos
  básicos.
\item
  Ambiente social extremamente seguro, ordeiro e resiliente.
\item
  Isolamento demográfico rural de alto nível (privacidade absoluta).
\item
  Infraestrutura de serviços e saúde privada de primeiro nível.
\end{itemize}

\subsubsection{Por que não}

\begin{itemize}
\tightlist
\item
  Alvo estratégico prioritário para controle nacional de suprimentos em
  crises.
\item
  Valor da terra entre os mais elevados da região.
\item
  Dependência tática de insumos externos para manutenção da escala
  mecanizada.
\end{itemize}

\subsubsection{Recomendação
Técnica}

Altamente indicado para estabelecimentos de autossuficiência tecnificada
de alto padrão ou moradia tática. Requer autonomia energética individual
para mitigar riscos de rede. O GSS de 8.0 reflete a excepcional solidez
dos recursos e da governança local em equilíbrio com a visibilidade
estratégica. Referencia atual de score: GSS 8.0 em 2026-03-05.

\subsubsection{}

\begin{itemize}
\tightlist
\item
  \begin{enumerate}
  \def\labelenumi{\arabic{enumi})}
  \tightlist
  \item
    Dados censitários validados (INE\footnote{Instituto Nacional de Estadística (Instituto Nacional de Estatística), órgão oficial de dados demográficos e censitários.} 2022).
  \end{enumerate}
\item
  \begin{enumerate}
  \def\labelenumi{\arabic{enumi})}
  \setcounter{enumi}{1}
  \tightlist
  \item
    Relatórios de produção agrícola mecanizada de alta precisão (MAG).
  \end{enumerate}
\item
  \begin{enumerate}
  \def\labelenumi{\arabic{enumi})}
  \setcounter{enumi}{2}
  \tightlist
  \item
    Mapa de solos e capacidade de silos regional (Portal
    Geoestadístico).
  \end{enumerate}
\item
  \begin{enumerate}
  \def\labelenumi{\arabic{enumi})}
  \setcounter{enumi}{3}
  \tightlist
  \item
    Registro de segurança pública e governança municipal próspera.
  \end{enumerate}
\end{itemize}}
\begin{table}[ht!]
\small \begin{tabular}{p{3.0cm}p{0.8cm}p{6.0cm}} \toprule
\textbf{Critério} & \textbf{Nota} & \textbf{Justificativa} \\ \midrule
A - Ameaças Estratégicas & 7.5 & Hub agroindustrial de elite; isolamento tático superior fora dos eixos de conflito primários \\
B - Risco Social & 8.5 & Densidade populacional baixa; isolamento social de alto nível em ambiente rural organizado \\
C - Riscos Naturais & 7.5 & Geologia muito estável e topografia segura contra inundações \\
D - Autossuficiência & 8.5 & Recursos agrícolas e hídricos excepcionais; solo de elite mundial \\
E - Institucional & 8.0 & Forte segurança institucional privada e serviços de primeiro nível regional \\
\midrule \textbf{GSS FINAL} & \textbf{8.0} & \textbf{Seguro (Altamente recomendado para autossuficiência e residência tática em ambiente de alta produtividade).} \\ \bottomrule \end{tabular}
\end{table}
\subsubsection{Vulnerabilidades e
Mitigação}\label{vuln-San_Alberto-vulnerabilidades-e-mitigauxe7uxe3o}

\begin{itemize}
\tightlist
\item
  \textbf{Visibilidade Tática de Suprimentos:} Sendo um polo de grãos, é
  um alvo estratégico em crises de desabastecimento (Mitigação:
  residência rural afastada 5-10 km do núcleo urbano industrial).
\item
  \textbf{Dependência de Exportação:} Economia vinculada ao mercado
  global (Mitigação: diversificação produtiva para consumo interno e
  autonomia energética).
\item
  \textbf{Gargalo Logístico:} Dependência da Ruta PY07 (Mitigação:
  mapeamento de rotas rurais vicinais em direção a Mbaracayú).
\end{itemize}
\subsubsection{Dossiê de Campo}\label{dossie-San_Alberto-dossiuxea-de-campo}
\begin{description}[style=nextline,leftmargin=0.5cm,font=\bfseries]
\item[Coordenadas]  24°58'S 54°54'W.
\item[Topografia]  Relevo predominantemente plano a suavemente ondulado, com solos de altíssima fertilidade (Terra Roxa mecanizada). Altitude \textasciitilde 250m. Área de \textasciitilde 332 km². Geologicamente muito estável, risco sísmico nulo.
\item[Alvos Estratégicos]  Polo agroindustrial do norte de Alto Paraná; Grande concentração de silos e moinhos de grãos; Ponto logístico na Ruta PY07. Ausência de bases militares massivas, oferecendo excelente invisibilidade estratégica e proteção por anonimato rural produtivo.
\item[Fallout]  Ventos predominantes do Norte (NE) e Leste. Baixíssimo risco de fallout direto; posição favorável em relação aos centros metropolitanos.
\item[População]  11.162 habitantes (Censo 2022 Final). População cosmopolita com forte herança de imigrantes brasileiros, altamente tecnificada e integrada ao agronegócio de alta performance.
\item[Densidade]  \textasciitilde 33,6 hab/\allowbreak km² (Densidade baixa, garantindo privacidade absoluta e ausência de riscos sociais de massa urbana fora do núcleo comercial).
\item[Vias de Acesso]  Localização estratégica sobre a Ruta PY07. Conectividade asfáltica de elite para Ciudad del Este e Salto del Guairá. Facilidade logística tática em tempos de paz, funcionando como um enclave de riqueza protegido pela geografia.
\item[Serviços]  Infraestrutura de serviços privados de alta qualidade. Sanatórios modernos, rede bancária completa e suporte técnico agroindustrial avançado.
\item[Custo de Vida]  Médio; economia dolarizada pelo agronegócio e alta valorização das áreas mecanizadas.
\item[Preço da Terra]  US\$ 12.000-18.000/\allowbreak ha (terras mecanizadas de elite mundial).
\item[Hidrologia]  Baixo risco de inundações sistêmicas devido à topografia favorável e gestão hídrica eficiente. Presença de arroios perenes de alta qualidade.
\item[Clima]  Subtropical Úmido. Verões quentes; invernos frescos com geadas ocasionais benéficas.
\item[Energia]  Rede nacional (Itaipu); boa estabilidade devido à importância industrial. Nodo de distribuição regional robusto.
\item[Água]  Abundante via poços artesianos profundos e lençol freático produtivo; excelente qualidade hídrica subterrânea.
\item[Qualidade do Solo]  Latossolo Vermelho de Elite; solo profundo, estruturado e de fertilidade máxima mundial. Excelente aptidão para soja, milho e trigo.
\item[Resources Locais]  Grande polo produtor de grãos e proteína animal. Elevadíssimo potencial para autossuficiência alimentar básica em nível de propriedade devido à vasta área produtiva e baixa população rural.
\item[Segurança]  Uma das cidades rurais mais seguras e prósperas do país. Baixíssima criminalidade urbana. Coesão social fortíssima baseada na cultura produtora. Presença de segurança patrimonial privada de alta performance em todas as instalações agroindustriais.
\item[Leis Local]  Município consolidado e polo de influência no norte do departamento ("Capital do Trabalho"). \\ \textbf{Livre de Restrição} de Fronteira. \\ \textit{Aquisição de terra rural proibida para estrangeiros limítrofes.}
\end{description}
\paragraph{Solo (SoilGrids 2.0, média ponderada 0--30
cm)}

\begin{longtable}[]{@{}ll@{}}
\toprule\noalign{}
Parâmetro & Valor \\
\midrule\noalign{}
\endhead
\bottomrule\noalign{}
\endlastfoot
pH (H$_2$O) & 5.4 \\
Carbono orgânico (SOC) & 23.37 g/\allowbreak kg \\
Argila & 44.02 \% \\
Areia & 27.35 \% \\
Silte (calc.) & 28.6 \% \\
Densidade aparente & 1.25 g/\allowbreak cm³ \\
\textbf{Aptidão agrícola} & \textbf{Média} \\
\end{longtable}

Fonte: ISRIC SoilGrids 2.0 via WCS (acesso em 2026-03-21). Coords:
-24.9667°, -54.9°. Média ponderada camadas 0-5, 5-15, 15-30 cm.

\begin{itemize}
\tightlist
\item
  \textbf{Homicidios/\allowbreak 100k:} \textasciitilde5,0 (Muito baixo/\allowbreak Produtivo).
\end{itemize}

\subsection{10. Analise de Riscos}

\begin{itemize}
\tightlist
\item
  \textbf{Cenario curto prazo:} Manutenção da prosperidade econômica e
  expansão das agroindústrias.
\item
  \textbf{Cenario medio prazo:} Impacto de mudanças táticas na capital
  sobre o controle das rotas de escoamento do norte.
\end{itemize}

\subsection{11. Redacao Final}

\subsubsection{Diagnostico Integrado}

San Alberto é o símbolo da potência produtiva do norte de Alto Paraná.
Oferece solos de produtividade mundial, abundância massiva de recursos
hídricos e alimentares, e uma segurança social sustentada por uma
comunidade produtora altamente organizada. Sua força absoluta reside na
resiliência de sua base produtiva e na baixa densidade populacional,
garantindo privacidade e isolamento tático superior. É o local ideal
para quem busca autonomia total em ambiente de alta performance,
transformando a riqueza de recursos básicos em uma armadura definitiva
de proteção contra instabilidades externas.

\subsubsection{Por que sim}

\begin{itemize}
\tightlist
\item
  Solos de altíssima produtividade mundial e abundância de recursos
  básicos.
\item
  Ambiente social extremamente seguro, ordeiro e resiliente.
\item
  Isolamento demográfico rural de alto nível (privacidade absoluta).
\item
  Infraestrutura de serviços e saúde privada de primeiro nível.
\end{itemize}

\subsubsection{Por que nao}

\begin{itemize}
\tightlist
\item
  Alvo estratégico prioritário para controle nacional de suprimentos em
  crises.
\item
  Valor da terra entre os mais elevados da região.
\item
  Dependência tática de insumos externos para manutenção da escala
  mecanizada.
\end{itemize}

\subsubsection{Recomendacao Tecnica}

Altamente indicado para estabelecimentos de autossuficiência tecnificada
de alto padrão ou moradia tática. Requer autonomia energética individual
para mitigar riscos de rede. O GSS de 8.0 reflete a excepcional solidez
dos recursos e da governança local em equilíbrio com a visibilidade
estratégica. Referencia atual de score: GSS 8.0 em 2026-03-05.

\subsection{13.}

\begin{itemize}
\tightlist
\item
  \begin{enumerate}
  \def\labelenumi{\arabic{enumi})}
  \tightlist
  \item
    Dados censitários validados (INE\footnote{Instituto Nacional de Estadística (Instituto Nacional de Estatística), órgão oficial de dados demográficos e censitários.} 2022).
  \end{enumerate}
\item
  \begin{enumerate}
  \def\labelenumi{\arabic{enumi})}
  \setcounter{enumi}{1}
  \tightlist
  \item
    Relatórios de produção agrícola mecanizada de alta precisão (MAG).
  \end{enumerate}
\item
  \begin{enumerate}
  \def\labelenumi{\arabic{enumi})}
  \setcounter{enumi}{2}
  \tightlist
  \item
    Mapa de solos e capacidade de silos regional (Portal
    Geoestadístico).
  \end{enumerate}
\item
  \begin{enumerate}
  \def\labelenumi{\arabic{enumi})}
  \setcounter{enumi}{3}
  \tightlist
  \item
    Registro de segurança pública e governança municipal próspera.
  \end{enumerate}
\end{itemize}

\subsubsection{Combustível}

\textbf{Referência:} PETROPAR /\allowbreak  postos locais (2024) \textbf{Tipo de
localidade:} Interior

\begin{longtable}[]{@{}lll@{}}
\toprule\noalign{}
Combustível & USD/\allowbreak litro & Gs/\allowbreak litro (aprox.) \\
\midrule\noalign{}
\endhead
\bottomrule\noalign{}
\endlastfoot
Gasolina 93 oct & 0.97 & 7,178 \\
Gasolina 97 oct (premium) & 1.07 & 7,918 \\
Diesel & 0.90 & 6,660 \\
\end{longtable}

\begin{quote}
Preços podem variar ±5\% conforme posto e sazonalidade. Chaco e interior
remoto apresentam maior variação.
\end{quote}

\subsubsection{Cobertura Celular}

\textbf{Fonte:} CONATEL PY /\allowbreak  operadoras (2024)

\begin{longtable}[]{@{}ll@{}}
\toprule\noalign{}
Parâmetro & Valor \\
\midrule\noalign{}
\endhead
\bottomrule\noalign{}
\endlastfoot
Cobertura 4G população (dept.) & 95\% \\
Cobertura 4G área rural & 88\% \\
Melhor operadora & Tigo/\allowbreak Personal \\
Qualidade rural & boa \\
\end{longtable}

\begin{quote}
Para áreas rurais fora do núcleo urbano, recomenda-se chip Tigo como
principal e Personal como backup.
\end{quote}

\subsubsection{Internet}

\textbf{Fonte:} CONATEL /\allowbreak  Speedtest Ookla (2024)

\begin{longtable}[]{@{}ll@{}}
\toprule\noalign{}
Parâmetro & Valor \\
\midrule\noalign{}
\endhead
\bottomrule\noalign{}
\endlastfoot
Velocidade média download & 85 Mbps \\
Domicílios com internet (dept.) & 87\% \\
Tecnologia predominante & fibra/\allowbreak cabo \\
Opção rural & Starlink disponível (\textasciitilde USD 44/\allowbreak mês) \\
\end{longtable}

\subsubsection{Mercado Imobiliário e Terra
Rural}

\textbf{Fonte:} INDERT /\allowbreak  Clasificados.com.py (2024)

\begin{longtable}[]{@{}ll@{}}
\toprule\noalign{}
Tipo & Referência \\
\midrule\noalign{}
\endhead
\bottomrule\noalign{}
\endlastfoot
Terra agrícola alta prod. (USD/\allowbreak ha) & 14,800 \\
Imóvel urbano (USD/\allowbreak m²) & 1,400 \\
Aluguel 2 quartos (USD/\allowbreak mês) & 600 \\
\end{longtable}

\begin{quote}
Valores de referência departamental. Localidades menores podem ter
preços 20--40\% abaixo da capital departamental. \#\#\# Saúde
\end{quote}

\textbf{Fonte:} MSPBS /\allowbreak  IPS Paraguay (2024-2026), consolidação
departamental e proxy local

\begin{longtable}[]{@{}ll@{}}
\toprule\noalign{}
Serviço & Disponibilidade \\
\midrule\noalign{}
\endhead
\bottomrule\noalign{}
\endlastfoot
USF /\allowbreak  Posto de Saúde & sim \\
Hospital Regional & não \\
IPS (seguro social) & sim \\
Farmácia & sim \\
Distância ao hospital de referência & \textasciitilde84 km (Ciudad del
Este) \\
\end{longtable}

\textbf{Principais estabelecimentos:} USF local; referencia hospitalar
em Ciudad del Este

\textbf{Observação para imigrantes:} Atendimento primario local ou em
raio curto; casos de maior complexidade seguem para Ciudad del Este.
Cobertura privada continua recomendada para especialidades e urgências
de maior complexidade.
\subsubsection{Indicadores Sociais}

\textbf{Fontes:} PNUD 2020, INE\footnote{Instituto Nacional de Estadística (Instituto Nacional de Estatística), órgão oficial de dados demográficos e censitários.} Censo 2022, INE\footnote{Instituto Nacional de Estadística (Instituto Nacional de Estatística), órgão oficial de dados demográficos e censitários.} EPHC 2023, Ministerio
Público 2024\\
\textbf{Nota:} Valores marcados como \emph{(dept.)} referem-se ao
departamento de Alto Paraná; valores \emph{(dist.)} são específicos
deste distrito. Dados marcados \emph{(est.)} são estimativas.

\begin{longtable}[]{@{}
  >{\raggedright\arraybackslash}p{(\linewidth - 4\tabcolsep) * \real{0.4231}}
  >{\raggedright\arraybackslash}p{(\linewidth - 4\tabcolsep) * \real{0.2692}}
  >{\raggedright\arraybackslash}p{(\linewidth - 4\tabcolsep) * \real{0.3077}}@{}}
\toprule\noalign{}
\begin{minipage}[b]{\linewidth}\raggedright
Indicador
\end{minipage} & \begin{minipage}[b]{\linewidth}\raggedright
Valor
\end{minipage} & \begin{minipage}[b]{\linewidth}\raggedright
Âmbito
\end{minipage} \\
\midrule\noalign{}
\endhead
\bottomrule\noalign{}
\endlastfoot
IDH (2020) & 0,752 & dept. \\
Ranking IDH nacional & 3/\allowbreak 18 & dept. \\
Esperança de vida & 75,1 anos & dept. \\
Escolaridade média & 8.5 anos & dist. \\
RNB per capita (USD PPA) & 13.500 & dept. \\
Pobreza monetária (\%) & 16,3\% & dept. \\
Pobreza extrema (\%) & 3,1\% & dept. \\
Índice de Gini & 0,428 & dept. \\
Acesso a água potável (\%) & 88,7\% & dept. \\
Acesso a saneamento (\%) & 87,2\% & dept. \\
Taxa de homicídios (est., /\allowbreak 100k hab) & \textasciitilde8--12 (estimativa,
acima da média nacional) & dept. \\
Índice de segurança & baixo & dept. \\
População (Censo 2022) & 11.162 hab. & dist. \\
Idade mediana & 29.0 anos & dist. \\
Taxa de fecundidade & N/\allowbreak D & dist. \\
\end{longtable}
\subsubsection{Combustível}

\textbf{Referência:} PETROPAR /\allowbreak  postos locais (2024) \textbf{Tipo de
localidade:} Interior

\begin{longtable}[]{@{}lll@{}}
\toprule\noalign{}
Combustível & USD/\allowbreak litro & Gs/\allowbreak litro (aprox.) \\
\midrule\noalign{}
\endhead
\bottomrule\noalign{}
\endlastfoot
Gasolina 93 oct & 0.97 & 7,178 \\
Gasolina 97 oct (premium) & 1.07 & 7,918 \\
Diesel & 0.90 & 6,660 \\
\end{longtable}

\begin{quote}
Preços podem variar ±5\% conforme posto e sazonalidade. Chaco e interior
remoto apresentam maior variação.
\end{quote}

\subsubsection{Cobertura Celular}

\textbf{Fonte:} CONATEL PY /\allowbreak  operadoras (2024)

\begin{longtable}[]{@{}ll@{}}
\toprule\noalign{}
Parâmetro & Valor \\
\midrule\noalign{}
\endhead
\bottomrule\noalign{}
\endlastfoot
Cobertura 4G população (dept.) & 95\% \\
Cobertura 4G área rural & 88\% \\
Melhor operadora & Tigo/\allowbreak Personal \\
Qualidade rural & boa \\
\end{longtable}

\begin{quote}
Para áreas rurais fora do núcleo urbano, recomenda-se chip Tigo como
principal e Personal como backup.
\end{quote}

\subsubsection{Internet}

\textbf{Fonte:} CONATEL /\allowbreak  Speedtest Ookla (2024)

\begin{longtable}[]{@{}ll@{}}
\toprule\noalign{}
Parâmetro & Valor \\
\midrule\noalign{}
\endhead
\bottomrule\noalign{}
\endlastfoot
Velocidade média download & 85 Mbps \\
Domicílios com internet (dept.) & 87\% \\
Tecnologia predominante & fibra/\allowbreak cabo \\
Opção rural & Starlink disponível (\textasciitilde USD 44/\allowbreak mês) \\
\end{longtable}

\subsubsection{Mercado Imobiliário e Terra
Rural}

\textbf{Fonte:} INDERT /\allowbreak  Clasificados.com.py (2024)

\begin{longtable}[]{@{}ll@{}}
\toprule\noalign{}
Tipo & Referência \\
\midrule\noalign{}
\endhead
\bottomrule\noalign{}
\endlastfoot
Terra agrícola alta prod. (USD/\allowbreak ha) & 14,800 \\
Imóvel urbano (USD/\allowbreak m²) & 1,400 \\
Aluguel 2 quartos (USD/\allowbreak mês) & 600 \\
\end{longtable}

\begin{quote}
Valores de referência departamental. Localidades menores podem ter
preços 20--40\% abaixo da capital departamental. \#\#\# Saúde
\end{quote}

\textbf{Fonte:} MSPBS /\allowbreak  IPS Paraguay (2024-2026), consolidação
departamental e proxy local

\begin{longtable}[]{@{}ll@{}}
\toprule\noalign{}
Serviço & Disponibilidade \\
\midrule\noalign{}
\endhead
\bottomrule\noalign{}
\endlastfoot
USF /\allowbreak  Posto de Saúde & sim \\
Hospital Regional & não \\
IPS (seguro social) & sim \\
Farmácia & sim \\
Distância ao hospital de referência & \textasciitilde84 km (Ciudad del
Este) \\
\end{longtable}

\textbf{Principais estabelecimentos:} USF local; referencia hospitalar
em Ciudad del Este

\textbf{Observação para imigrantes:} Atendimento primario local ou em
raio curto; casos de maior complexidade seguem para Ciudad del Este.
Cobertura privada continua recomendada para especialidades e urgências
de maior complexidade.
\newpage
\secaoDiagnostico{San Cristobal}{\mapaConteudo{-25.9166}{-55.4500}}{San Cristóbal é o enclave de resiliência e produtividade do sudoeste de
Alto Paraná. Oferece solos de produtividade mundial e um ambiente social
tranquilo, ``fora do radar'' das grandes crises demográficas
metropolitanas. Sua força absoluta reside na abundância de recursos
naturais e na eficiência da organização comunitária via cooperativas. É
o local ideal para quem busca autonomia total em ambiente de alta
performance, transformando a riqueza de recursos básicos em uma barreira
de proteção tática contra instabilidades externas.

\subsubsection{Por que sim}

\begin{itemize}
\tightlist
\item
  Solos de altíssima produtividade e abundância de recursos hídricos.
\item
  Ambiente social extremamente seguro, ordeiro e resiliente.
\item
  Isolamento demográfico rural de alto nível (privacidade absoluta).
\item
  Sem restrições da Lei de Fronteira.
\end{itemize}

\subsubsection{Por que não}

\begin{itemize}
\tightlist
\item
  Alvo estratégico logístico e alimentar em cenários de conflito.
\item
  Valor da terra entre os mais altos da região.
\item
  Dependência tática de centros maiores para serviços de alta
  complexidade.
\end{itemize}

\subsubsection{Recomendação
Técnica}

Altamente indicado para estabelecimentos de autossuficiência tecnificada
que busquem solo de alta performance. Requer autonomia energética
individual e investimento em segurança patrimonial. O GSS de 7.5 reflete
a solidez dos recursos e da governança local em equilíbrio com o
isolamento tático. Referencia atual de score: GSS 7.5 em 2026-03-05.

\subsubsection{}

\begin{itemize}
\tightlist
\item
  \begin{enumerate}
  \def\labelenumi{\arabic{enumi})}
  \tightlist
  \item
    Dados censitários validados (INE\footnote{Instituto Nacional de Estadística (Instituto Nacional de Estatística), órgão oficial de dados demográficos e censitários.} 2022).
  \end{enumerate}
\item
  \begin{enumerate}
  \def\labelenumi{\arabic{enumi})}
  \setcounter{enumi}{1}
  \tightlist
  \item
    Relatórios de produção agrícola mecanizada departamental (MAG).
  \end{enumerate}
\item
  \begin{enumerate}
  \def\labelenumi{\arabic{enumi})}
  \setcounter{enumi}{2}
  \tightlist
  \item
    Mapa de solos e hidrografia de Alto Paraná (Portal Geoestadístico).
  \end{enumerate}
\item
  \begin{enumerate}
  \def\labelenumi{\arabic{enumi})}
  \setcounter{enumi}{3}
  \tightlist
  \item
    Registro de segurança pública e governança municipal.
  \end{enumerate}
\end{itemize}}
\begin{table}[ht!]
\small \begin{tabular}{p{3.0cm}p{0.8cm}p{6.0cm}} \toprule
\textbf{Critério} & \textbf{Nota} & \textbf{Justificativa} \\ \midrule
A - Ameaças Estratégicas & 7.0 & Hub agroindustrial estratégico; boa invisibilidade tática fora dos eixos de conflito primários \\
B - Risco Social & 8.0 & Densidade populacional extremamente baixa; isolamento social de alto nível em ambiente rural organizado \\
C - Riscos Naturais & 7.0 & Geologia muito estável e topografia segura contra inundações \\
D - Autossuficiência & 8.5 & Recursos agrícolas e hídricos excepcionais; solo de elite e silos modernos \\
E - Institucional & 7.0 & Forte segurança institucional privada e serviços funcionais próximos \\
\midrule \textbf{GSS FINAL} & \textbf{7.5} & \textbf{Seguro (Altamente recomendado para quem busca isolamento tático rural em terras férteis com excelente suporte organizacional).} \\ \bottomrule \end{tabular}
\end{table}
\subsubsection{Vulnerabilidades e
Mitigação}\label{vuln-San_Cristobal-vulnerabilidades-e-mitigauxe7uxe3o}

\begin{itemize}
\tightlist
\item
  \textbf{Visibilidade Tática de Suprimentos:} Sendo um polo
  agroindustrial, é um alvo estratégico para controle de suprimentos em
  crises (Mitigação: residência rural afastada 5-10 km do núcleo urbano
  industrial e integração com as redes de defesa cooperativa).
\item
  \textbf{Dependência de Exportação:} Economia vinculada ao mercado
  global (Mitigação: diversificação produtiva para consumo interno e
  autonomia energética).
\item
  \textbf{Fragilidade Elétrica Rural:} Quedas ocasionais em tempestades
  (Mitigação: instalação de sistemas solares fotovoltaicos potentes).
\end{itemize}
\subsubsection{Dossiê de Campo}\label{dossie-San_Cristobal-dossiuxea-de-campo}
\begin{description}[style=nextline,leftmargin=0.5cm,font=\bfseries]
\item[Coordenadas]  25°55'S 55°27'W.
\item[Topografia]  Relevo predominantemente plano a suavemente ondulado, com solos de altíssima produtividade (Terra Roxa mecanizada). Altitude \textasciitilde 240m. Área vasta de \textasciitilde 1.027 km². Geologicamente muito estável, risco sísmico nulo.
\item[Alvos Estratégicos]  Polo agroindustrial do sudoeste de Alto Paraná; Grande concentração de silos e moinhos da Cooperativa Pindo; Proximidade tática com o Rio Monday. Ausência de alvos militares diretos, oferecendo excelente invisibilidade estratégica, sendo um dos motores alimentares regionais.
\item[Fallout]  Ventos predominantes do Norte (NE) e Leste. Baixíssimo risco de fallout direto; posição isolada e protegida pela distância dos grandes centros urbanos.
\item[População]  8.563 habitantes (Censo 2022 Final). População com forte herança de imigrantes brasileiros e europeus, altamente integrada ao sistema cooperativista mecanizado.
\item[Densidade]  \textasciitilde 8,3 hab/\allowbreak km² (Densidade extremamente baixa, garantindo privacidade absoluta e ausência de riscos sociais de massa urbana).
\item[Vias de Acesso]  Acesso via ramais pavimentados de alta qualidade que conectam à Ruta PY06. Localização protegida por estar fora do radar dos grandes fluxos de trânsito internacional pesado, funcionando como um enclave de riqueza e tranquilidade rural.
\item[Serviços]  Infraestrutura de serviços de primeiro nível regional sustentada pela Cooperativa Pindo. Sanatórios modernos, suporte técnico de precisão e rede bancária consolidada.
\item[Custo de Vida]  Médio; economia dolarizada pelo agronegócio e alta valorização imobiliária das terras mecanizadas.
\item[Preço da Terra]  US\$ 12.000-18.000/\allowbreak ha (terras mecanizadas de elite - "Terra Roxa").
\item[Hidrologia]  Baixo risco de inundações sistêmicas devido à topografia favorável e gestão hídrica de precisão. Presença de arroios perenes de alta qualidade hídrica utilizados para irrigação e pecuária.
\item[Clima]  Subtropical Úmido. Verões quentes; invernos frescos com geadas benéficas ao trigo.
\item[Energia]  Rede nacional (Itaipu); boa estabilidade devido à importância industrial das cooperativas locais. Uso comum de sistemas solares e geradores potentes em propriedades de elite.
\item[Água]  Abundante via poços artesianos profundos e lençol freático produtivo; excelente qualidade hídrica subterrânea.
\item[Qualidade do Solo]  Latossolo Vermelho de Elite; solo profundo, estruturado e de fertilidade máxima mundial.
\item[Resources Locais]  Grande produção de grãos (soja, milho, trigo) e proteína animal. Elevadíssimo potencial para autossuficiência alimentar básica em nível de propriedade devido à vasta área produtiva e baixa população.
\item[Segurança]  Uma das zonas rurais mais seguras e estáveis do Paraguai. Baixíssima criminalidade urbana. Coesão social fortíssima baseada na cultura cooperativista e produtora. Presença de segurança patrimonial privada de alta performance.
\item[Leis Local]  Município consolidado e polo de influência regional. \\ \textbf{Livre de Restrição} de Fronteira. \\ \textit{Aquisição de terra rural proibida para estrangeiros limítrofes.}
\end{description}
\paragraph{Solo (SoilGrids 2.0, média ponderada 0--30
cm)}

\begin{longtable}[]{@{}ll@{}}
\toprule\noalign{}
Parâmetro & Valor \\
\midrule\noalign{}
\endhead
\bottomrule\noalign{}
\endlastfoot
pH (H$_2$O) & 5.35 \\
Carbono orgânico (SOC) & 21.54 g/\allowbreak kg \\
Argila & 43.54 \% \\
Areia & 26.86 \% \\
Silte (calc.) & 29.6 \% \\
Densidade aparente & 1.21 g/\allowbreak cm³ \\
\textbf{Aptidão agrícola} & \textbf{Média} \\
\end{longtable}

Fonte: ISRIC SoilGrids 2.0 via WCS (acesso em 2026-03-21). Coords:
-25.9167°, -55.45°. Média ponderada camadas 0-5, 5-15, 15-30 cm.

\begin{itemize}
\tightlist
\item
  \textbf{Homicidios/\allowbreak 100k:} \textasciitilde5,0 (Abaixo da média
  nacional).
\end{itemize}

\subsection{10. Analise de Riscos}

\begin{itemize}
\tightlist
\item
  \textbf{Cenario curto prazo:} Manutenção da prosperidade econômica e
  estabilidade social cooperativista.
\item
  \textbf{Cenario medio prazo:} Impacto de mudanças climáticas na
  rentabilidade das safras mecanizadas.
\end{itemize}

\subsection{11. Redacao Final}

\subsubsection{Diagnostico Integrado}

San Cristóbal é o enclave de resiliência e produtividade do sudoeste de
Alto Paraná. Oferece solos de produtividade mundial e um ambiente social
tranquilo, ``fora do radar'' das grandes crises demográficas
metropolitanas. Sua força absoluta reside na abundância de recursos
naturais e na eficiência da organização comunitária via cooperativas. É
o local ideal para quem busca autonomia total em ambiente de alta
performance, transformando a riqueza de recursos básicos em uma barreira
de proteção tática contra instabilidades externas.

\subsubsection{Por que sim}

\begin{itemize}
\tightlist
\item
  Solos de altíssima produtividade e abundância de recursos hídricos.
\item
  Ambiente social extremamente seguro, ordeiro e resiliente.
\item
  Isolamento demográfico rural de alto nível (privacidade absoluta).
\item
  Sem restrições da Lei de Fronteira.
\end{itemize}

\subsubsection{Por que nao}

\begin{itemize}
\tightlist
\item
  Alvo estratégico logístico e alimentar em cenários de conflito.
\item
  Valor da terra entre os mais altos da região.
\item
  Dependência tática de centros maiores para serviços de alta
  complexidade.
\end{itemize}

\subsubsection{Recomendacao Tecnica}

Altamente indicado para estabelecimentos de autossuficiência tecnificada
que busquem solo de alta performance. Requer autonomia energética
individual e investimento em segurança patrimonial. O GSS de 7.5 reflete
a solidez dos recursos e da governança local em equilíbrio com o
isolamento tático. Referencia atual de score: GSS 7.5 em 2026-03-05.

\subsection{13.}

\begin{itemize}
\tightlist
\item
  \begin{enumerate}
  \def\labelenumi{\arabic{enumi})}
  \tightlist
  \item
    Dados censitários validados (INE\footnote{Instituto Nacional de Estadística (Instituto Nacional de Estatística), órgão oficial de dados demográficos e censitários.} 2022).
  \end{enumerate}
\item
  \begin{enumerate}
  \def\labelenumi{\arabic{enumi})}
  \setcounter{enumi}{1}
  \tightlist
  \item
    Relatórios de produção agrícola mecanizada departamental (MAG).
  \end{enumerate}
\item
  \begin{enumerate}
  \def\labelenumi{\arabic{enumi})}
  \setcounter{enumi}{2}
  \tightlist
  \item
    Mapa de solos e hidrografia de Alto Paraná (Portal Geoestadístico).
  \end{enumerate}
\item
  \begin{enumerate}
  \def\labelenumi{\arabic{enumi})}
  \setcounter{enumi}{3}
  \tightlist
  \item
    Registro de segurança pública e governança municipal.
  \end{enumerate}
\end{itemize}

\subsubsection{Combustível}

\textbf{Referência:} PETROPAR /\allowbreak  postos locais (2024) \textbf{Tipo de
localidade:} Interior

\begin{longtable}[]{@{}lll@{}}
\toprule\noalign{}
Combustível & USD/\allowbreak litro & Gs/\allowbreak litro (aprox.) \\
\midrule\noalign{}
\endhead
\bottomrule\noalign{}
\endlastfoot
Gasolina 93 oct & 0.97 & 7,178 \\
Gasolina 97 oct (premium) & 1.07 & 7,918 \\
Diesel & 0.90 & 6,660 \\
\end{longtable}

\begin{quote}
Preços podem variar ±5\% conforme posto e sazonalidade. Chaco e interior
remoto apresentam maior variação.
\end{quote}

\subsubsection{Cobertura Celular}

\textbf{Fonte:} CONATEL PY /\allowbreak  operadoras (2024)

\begin{longtable}[]{@{}ll@{}}
\toprule\noalign{}
Parâmetro & Valor \\
\midrule\noalign{}
\endhead
\bottomrule\noalign{}
\endlastfoot
Cobertura 4G população (dept.) & 95\% \\
Cobertura 4G área rural & 88\% \\
Melhor operadora & Tigo/\allowbreak Personal \\
Qualidade rural & boa \\
\end{longtable}

\begin{quote}
Para áreas rurais fora do núcleo urbano, recomenda-se chip Tigo como
principal e Personal como backup.
\end{quote}

\subsubsection{Internet}

\textbf{Fonte:} CONATEL /\allowbreak  Speedtest Ookla (2024)

\begin{longtable}[]{@{}ll@{}}
\toprule\noalign{}
Parâmetro & Valor \\
\midrule\noalign{}
\endhead
\bottomrule\noalign{}
\endlastfoot
Velocidade média download & 85 Mbps \\
Domicílios com internet (dept.) & 87\% \\
Tecnologia predominante & fibra/\allowbreak cabo \\
Opção rural & Starlink disponível (\textasciitilde USD 44/\allowbreak mês) \\
\end{longtable}

\subsubsection{Mercado Imobiliário e Terra
Rural}

\textbf{Fonte:} INDERT /\allowbreak  Clasificados.com.py (2024)

\begin{longtable}[]{@{}ll@{}}
\toprule\noalign{}
Tipo & Referência \\
\midrule\noalign{}
\endhead
\bottomrule\noalign{}
\endlastfoot
Terra agrícola alta prod. (USD/\allowbreak ha) & 14,800 \\
Imóvel urbano (USD/\allowbreak m²) & 1,400 \\
Aluguel 2 quartos (USD/\allowbreak mês) & 600 \\
\end{longtable}

\begin{quote}
Valores de referência departamental. Localidades menores podem ter
preços 20--40\% abaixo da capital departamental. \#\#\# Saúde
\end{quote}

\textbf{Fonte:} MSPBS /\allowbreak  IPS Paraguay (2024-2026), consolidação
departamental e proxy local

\begin{longtable}[]{@{}ll@{}}
\toprule\noalign{}
Serviço & Disponibilidade \\
\midrule\noalign{}
\endhead
\bottomrule\noalign{}
\endlastfoot
USF /\allowbreak  Posto de Saúde & sim \\
Hospital Regional & não \\
IPS (seguro social) & sim \\
Farmácia & sim \\
Distância ao hospital de referência & \textasciitilde104 km (Ciudad del
Este) \\
\end{longtable}

\textbf{Principais estabelecimentos:} USF local; referencia hospitalar
em Ciudad del Este

\textbf{Observação para imigrantes:} Atendimento primario local ou em
raio curto; casos de maior complexidade seguem para Ciudad del Este.
Cobertura privada continua recomendada para especialidades e urgências
de maior complexidade.
\subsubsection{Indicadores Sociais}

\textbf{Fontes:} PNUD 2020, INE\footnote{Instituto Nacional de Estadística (Instituto Nacional de Estatística), órgão oficial de dados demográficos e censitários.} Censo 2022, INE\footnote{Instituto Nacional de Estadística (Instituto Nacional de Estatística), órgão oficial de dados demográficos e censitários.} EPHC 2023, Ministerio
Público 2024\\
\textbf{Nota:} Valores marcados como \emph{(dept.)} referem-se ao
departamento de Alto Paraná; valores \emph{(dist.)} são específicos
deste distrito. Dados marcados \emph{(est.)} são estimativas.

\begin{longtable}[]{@{}
  >{\raggedright\arraybackslash}p{(\linewidth - 4\tabcolsep) * \real{0.4231}}
  >{\raggedright\arraybackslash}p{(\linewidth - 4\tabcolsep) * \real{0.2692}}
  >{\raggedright\arraybackslash}p{(\linewidth - 4\tabcolsep) * \real{0.3077}}@{}}
\toprule\noalign{}
\begin{minipage}[b]{\linewidth}\raggedright
Indicador
\end{minipage} & \begin{minipage}[b]{\linewidth}\raggedright
Valor
\end{minipage} & \begin{minipage}[b]{\linewidth}\raggedright
Âmbito
\end{minipage} \\
\midrule\noalign{}
\endhead
\bottomrule\noalign{}
\endlastfoot
IDH (2020) & 0,752 & dept. \\
Ranking IDH nacional & 3/\allowbreak 18 & dept. \\
Esperança de vida & 75,1 anos & dept. \\
Escolaridade média & 8.4 anos & dist. \\
RNB per capita (USD PPA) & 13.500 & dept. \\
Pobreza monetária (\%) & 16,3\% & dept. \\
Pobreza extrema (\%) & 3,1\% & dept. \\
Índice de Gini & 0,428 & dept. \\
Acesso a água potável (\%) & 88,7\% & dept. \\
Acesso a saneamento (\%) & 87,2\% & dept. \\
Taxa de homicídios (est., /\allowbreak 100k hab) & \textasciitilde8--12 (estimativa,
acima da média nacional) & dept. \\
Índice de segurança & baixo & dept. \\
População (Censo 2022) & 8.563 hab. & dist. \\
Idade mediana & 27.0 anos & dist. \\
Taxa de fecundidade & N/\allowbreak D & dist. \\
\end{longtable}
\subsubsection{Combustível}

\textbf{Referência:} PETROPAR /\allowbreak  postos locais (2024) \textbf{Tipo de
localidade:} Interior

\begin{longtable}[]{@{}lll@{}}
\toprule\noalign{}
Combustível & USD/\allowbreak litro & Gs/\allowbreak litro (aprox.) \\
\midrule\noalign{}
\endhead
\bottomrule\noalign{}
\endlastfoot
Gasolina 93 oct & 0.97 & 7,178 \\
Gasolina 97 oct (premium) & 1.07 & 7,918 \\
Diesel & 0.90 & 6,660 \\
\end{longtable}

\begin{quote}
Preços podem variar ±5\% conforme posto e sazonalidade. Chaco e interior
remoto apresentam maior variação.
\end{quote}

\subsubsection{Cobertura Celular}

\textbf{Fonte:} CONATEL PY /\allowbreak  operadoras (2024)

\begin{longtable}[]{@{}ll@{}}
\toprule\noalign{}
Parâmetro & Valor \\
\midrule\noalign{}
\endhead
\bottomrule\noalign{}
\endlastfoot
Cobertura 4G população (dept.) & 95\% \\
Cobertura 4G área rural & 88\% \\
Melhor operadora & Tigo/\allowbreak Personal \\
Qualidade rural & boa \\
\end{longtable}

\begin{quote}
Para áreas rurais fora do núcleo urbano, recomenda-se chip Tigo como
principal e Personal como backup.
\end{quote}

\subsubsection{Internet}

\textbf{Fonte:} CONATEL /\allowbreak  Speedtest Ookla (2024)

\begin{longtable}[]{@{}ll@{}}
\toprule\noalign{}
Parâmetro & Valor \\
\midrule\noalign{}
\endhead
\bottomrule\noalign{}
\endlastfoot
Velocidade média download & 85 Mbps \\
Domicílios com internet (dept.) & 87\% \\
Tecnologia predominante & fibra/\allowbreak cabo \\
Opção rural & Starlink disponível (\textasciitilde USD 44/\allowbreak mês) \\
\end{longtable}

\subsubsection{Mercado Imobiliário e Terra
Rural}

\textbf{Fonte:} INDERT /\allowbreak  Clasificados.com.py (2024)

\begin{longtable}[]{@{}ll@{}}
\toprule\noalign{}
Tipo & Referência \\
\midrule\noalign{}
\endhead
\bottomrule\noalign{}
\endlastfoot
Terra agrícola alta prod. (USD/\allowbreak ha) & 14,800 \\
Imóvel urbano (USD/\allowbreak m²) & 1,400 \\
Aluguel 2 quartos (USD/\allowbreak mês) & 600 \\
\end{longtable}

\begin{quote}
Valores de referência departamental. Localidades menores podem ter
preços 20--40\% abaixo da capital departamental. \#\#\# Saúde
\end{quote}

\textbf{Fonte:} MSPBS /\allowbreak  IPS Paraguay (2024-2026), consolidação
departamental e proxy local

\begin{longtable}[]{@{}ll@{}}
\toprule\noalign{}
Serviço & Disponibilidade \\
\midrule\noalign{}
\endhead
\bottomrule\noalign{}
\endlastfoot
USF /\allowbreak  Posto de Saúde & sim \\
Hospital Regional & não \\
IPS (seguro social) & sim \\
Farmácia & sim \\
Distância ao hospital de referência & \textasciitilde104 km (Ciudad del
Este) \\
\end{longtable}

\textbf{Principais estabelecimentos:} USF local; referencia hospitalar
em Ciudad del Este

\textbf{Observação para imigrantes:} Atendimento primario local ou em
raio curto; casos de maior complexidade seguem para Ciudad del Este.
Cobertura privada continua recomendada para especialidades e urgências
de maior complexidade.
\newpage
\secaoDiagnostico{Santa Fe del Parana}{\mapaConteudo{-25.2333}{-54.7000}}{Santa Fe del Parana (Alto Parana) apresenta perfil operacional
consolidado para comparacao territorial, com leitura conjunta de
infraestrutura, risco hidroclimatico e capacidade de continuidade de
servicos essenciais.

\subsubsection{Por que sim}

\begin{itemize}
\tightlist
\item
  Base institucional de referencia disponivel e rastreavel.
\item
  Estrutura comparativa (A-E/\allowbreak GSS) consistente para decisao relativa.
\item
  Plano de mitigacao acionavel ja definido no dossie.
\end{itemize}

\subsubsection{Por que não}

\begin{itemize}
\tightlist
\item
  Serie oficial distrital granular ainda pode apresentar lacunas
  temporais.
\item
  Sensibilidade do score exige monitoramento continuo de risco e
  conectividade.
\item
  Decisao final depende de validacao de campo e janela temporal recente.
\end{itemize}

\subsubsection{Pre-condicoes Minimas de
Decisao}

\begin{itemize}
\tightlist
\item
  Atualizar indicadores criticos em ciclo trimestral.
\item
  Confirmar trafegabilidade e contingencia hidrica/\allowbreak energetica local.
\item
  Validar checklist de resiliencia domiciliar/\allowbreak logistica antes de decisao
  final.
\end{itemize}

\subsubsection{Recomendação
Técnica}

Uso recomendado como candidato em analise multicriterio, condicionado ao
cumprimento das pre-condicoes acima. Referencia atual de score: GSS 7.0
em 2026-03-05; risco predominante: moderado.

\subsubsection{}

\begin{itemize}
\tightlist
\item
  \begin{enumerate}
  \def\labelenumi{\arabic{enumi})}
  \tightlist
  \item
    Delimitacao territorial validada por georreferencia institucional
    (Fonte: acesso: 2026-03-05).
  \end{enumerate}
\item
  \begin{enumerate}
  \def\labelenumi{\arabic{enumi})}
  \setcounter{enumi}{1}
  \tightlist
  \item
    População municipal/\allowbreak distrital referenciada por base censitaria
    nacional (Fonte: acesso: 2026-03-05).
  \end{enumerate}
\item
  \begin{enumerate}
  \def\labelenumi{\arabic{enumi})}
  \setcounter{enumi}{2}
  \tightlist
  \item
    Comparabilidade intra-departamental apoiada em indicadores
    distritais oficiais (Fonte: acesso: 2026-03-05).
  \end{enumerate}
\item
  \begin{enumerate}
  \def\labelenumi{\arabic{enumi})}
  \setcounter{enumi}{3}
  \tightlist
  \item
    Estado macro de conectividade viaria ancorado em informacoes de
    rodovias (Fonte: acesso: 2026-03-05).
  \end{enumerate}
\item
  \begin{enumerate}
  \def\labelenumi{\arabic{enumi})}
  \setcounter{enumi}{4}
  \tightlist
  \item
    Exposicao hidrometeorologica monitorada por avisos oficiais (Fonte:
    acesso: 2026-03-05).
  \end{enumerate}
\item
  \begin{enumerate}
  \def\labelenumi{\arabic{enumi})}
  \setcounter{enumi}{5}
  \tightlist
  \item
    Historico de contexto meteorologico apoiado em publicacoes tecnicas
    (Fonte: acesso: 2026-03-05).
  \end{enumerate}
\item
  \begin{enumerate}
  \def\labelenumi{\arabic{enumi})}
  \setcounter{enumi}{6}
  \tightlist
  \item
    Capacidade institucional de resposta a eventos apoiada em acoes da
    SEN\footnote{Secretaría de Emergencia Nacional (Secretaria de Emergência Nacional), órgão governamental de resposta a desastres.} (Fonte: acesso: 2026-03-05).
  \end{enumerate}
\item
  \begin{enumerate}
  \def\labelenumi{\arabic{enumi})}
  \setcounter{enumi}{7}
  \tightlist
  \item
    Infraestrutura energetica analisada via operador nacional (Fonte:
    acesso: 2026-03-05).
  \end{enumerate}
\item
  \begin{enumerate}
  \def\labelenumi{\arabic{enumi})}
  \setcounter{enumi}{8}
  \tightlist
  \item
    Referencias de obras e logistica nacional verificadas no MOPC
    (Fonte: acesso: 2026-03-05).
  \end{enumerate}
\item
  \begin{enumerate}
  \def\labelenumi{\arabic{enumi})}
  \setcounter{enumi}{9}
  \tightlist
  \item
    Contexto sociopolitico institucional referenciado por autoridade
    eleitoral (Fonte: acesso: 2026-03-05).
  \end{enumerate}
\item
  \begin{enumerate}
  \def\labelenumi{\arabic{enumi})}
  \setcounter{enumi}{10}
  \tightlist
  \item
    Camada cartografica de apoio eleitoral/\allowbreak territorial consultada
    (Fonte: acesso: 2026-03-05).
  \end{enumerate}
\item
  \begin{enumerate}
  \def\labelenumi{\arabic{enumi})}
  \setcounter{enumi}{11}
  \tightlist
  \item
    Dados publicos complementares para verificacao cruzada (Fonte:
    acesso: 2026-03-05).
  \end{enumerate}
\end{itemize}}
\begin{table}[ht!]
\small \begin{tabular}{p{3.0cm}p{0.8cm}p{6.0cm}} \toprule
\textbf{Critério} & \textbf{Nota} & \textbf{Justificativa} \\ \midrule
A - Ameaças Estratégicas & 6.9 & - \\
B - Risco Social & 7.8 & - \\
C - Riscos Naturais & 5.1 & - \\
D - Autossuficiência & 7.3 & - \\
E - Institucional & 7.6 & - \\
\midrule \textbf{GSS FINAL} & \textbf{7.0} & \textbf{Seguro (bom potencial com preparacao).} \\ \bottomrule \end{tabular}
\end{table}
\subsubsection{Vulnerabilidades e
Mitigação}\label{vuln-Santa_Fe_del_Parana-vulnerabilidades-e-mitigauxe7uxe3o}

\begin{itemize}
\tightlist
\item
  Exposicao hidrometeorologica moderada (45-64/\allowbreak 100).
\item
  Conectividade viaria intermediaria (50-69/\allowbreak 100).
\item
  Estoque de contingencia para 14 dias (agua/\allowbreak alimentos/\allowbreak combustivel),
  priorizado quando risco hidro=61/\allowbreak 100.
\item
  Duplicar rotas/\allowbreak logistica quando conectividade viaria=57/\allowbreak 100 for
  \textless70; manter rota secundaria mapeada e testada trimestralmente.
\item
  Plano de energia resiliente com autonomia minima de 72h
  (geracao/\allowbreak backup), revisao semestral.
\end{itemize}
\subsubsection{Dossiê de Campo}\label{dossie-Santa_Fe_del_Parana-dossiuxea-de-campo}
\begin{description}[style=nextline,leftmargin=0.5cm,font=\bfseries]
\item[Coordenadas]  25°14'S 54°42'W (geocodificado via Nominatim).
\end{description}
\paragraph{Solo (SoilGrids 2.0, média ponderada 0--30
cm)}

\begin{longtable}[]{@{}ll@{}}
\toprule\noalign{}
Parâmetro & Valor \\
\midrule\noalign{}
\endhead
\bottomrule\noalign{}
\endlastfoot
pH (H$_2$O) & 5.42 \\
Carbono orgânico (SOC) & 22.47 g/\allowbreak kg \\
Argila & 51.92 \% \\
Areia & 21.64 \% \\
Silte (calc.) & 26.4 \% \\
Densidade aparente & 1.18 g/\allowbreak cm³ \\
\textbf{Aptidão agrícola} & \textbf{Média} \\
\end{longtable}

Fonte: ISRIC SoilGrids 2.0 via WCS (acesso em 2026-03-21). Coords:
-25.2333°, -54.7°. Média ponderada camadas 0-5, 5-15, 15-30 cm.
\subsubsection{Indicadores Sociais}

\textbf{Fontes:} PNUD 2020, INE\footnote{Instituto Nacional de Estadística (Instituto Nacional de Estatística), órgão oficial de dados demográficos e censitários.} Censo 2022, INE\footnote{Instituto Nacional de Estadística (Instituto Nacional de Estatística), órgão oficial de dados demográficos e censitários.} EPHC 2023, Ministerio
Público 2024\\
\textbf{Nota:} Valores marcados como \emph{(dept.)} referem-se ao
departamento de Alto Paraná; valores \emph{(dist.)} são específicos
deste distrito. Dados marcados \emph{(est.)} são estimativas.

\begin{longtable}[]{@{}
  >{\raggedright\arraybackslash}p{(\linewidth - 4\tabcolsep) * \real{0.4231}}
  >{\raggedright\arraybackslash}p{(\linewidth - 4\tabcolsep) * \real{0.2692}}
  >{\raggedright\arraybackslash}p{(\linewidth - 4\tabcolsep) * \real{0.3077}}@{}}
\toprule\noalign{}
\begin{minipage}[b]{\linewidth}\raggedright
Indicador
\end{minipage} & \begin{minipage}[b]{\linewidth}\raggedright
Valor
\end{minipage} & \begin{minipage}[b]{\linewidth}\raggedright
Âmbito
\end{minipage} \\
\midrule\noalign{}
\endhead
\bottomrule\noalign{}
\endlastfoot
IDH (2020) & 0,752 & dept. \\
Ranking IDH nacional & 3/\allowbreak 18 & dept. \\
Esperança de vida & 75,1 anos & dept. \\
Escolaridade média & 7.5 anos & dist. \\
RNB per capita (USD PPA) & 13.500 & dept. \\
Pobreza monetária (\%) & 16,3\% & dept. \\
Pobreza extrema (\%) & 3,1\% & dept. \\
Índice de Gini & 0,428 & dept. \\
Acesso a água potável (\%) & 88,7\% & dept. \\
Acesso a saneamento (\%) & 87,2\% & dept. \\
Taxa de homicídios (est., /\allowbreak 100k hab) & \textasciitilde8--12 (estimativa,
acima da média nacional) & dept. \\
Índice de segurança & baixo & dept. \\
População (Censo 2022) & 3.981 hab. & dist. \\
Idade mediana & 28.0 anos & dist. \\
Taxa de fecundidade & N/\allowbreak D & dist. \\
\end{longtable}
\subsubsection{Combustível}

\textbf{Referência:} PETROPAR /\allowbreak  postos locais (2024) \textbf{Tipo de
localidade:} Interior

\begin{longtable}[]{@{}lll@{}}
\toprule\noalign{}
Combustível & USD/\allowbreak litro & Gs/\allowbreak litro (aprox.) \\
\midrule\noalign{}
\endhead
\bottomrule\noalign{}
\endlastfoot
Gasolina 93 oct & 0.97 & 7,178 \\
Gasolina 97 oct (premium) & 1.07 & 7,918 \\
Diesel & 0.90 & 6,660 \\
\end{longtable}

\begin{quote}
Preços podem variar ±5\% conforme posto e sazonalidade. Chaco e interior
remoto apresentam maior variação.
\end{quote}

\subsubsection{Cobertura Celular}

\textbf{Fonte:} CONATEL PY /\allowbreak  operadoras (2024)

\begin{longtable}[]{@{}ll@{}}
\toprule\noalign{}
Parâmetro & Valor \\
\midrule\noalign{}
\endhead
\bottomrule\noalign{}
\endlastfoot
Cobertura 4G população (dept.) & 95\% \\
Cobertura 4G área rural & 88\% \\
Melhor operadora & Tigo/\allowbreak Personal \\
Qualidade rural & boa \\
\end{longtable}

\begin{quote}
Para áreas rurais fora do núcleo urbano, recomenda-se chip Tigo como
principal e Personal como backup.
\end{quote}

\subsubsection{Internet}

\textbf{Fonte:} CONATEL /\allowbreak  Speedtest Ookla (2024)

\begin{longtable}[]{@{}ll@{}}
\toprule\noalign{}
Parâmetro & Valor \\
\midrule\noalign{}
\endhead
\bottomrule\noalign{}
\endlastfoot
Velocidade média download & 85 Mbps \\
Domicílios com internet (dept.) & 87\% \\
Tecnologia predominante & fibra/\allowbreak cabo \\
Opção rural & Starlink disponível (\textasciitilde USD 44/\allowbreak mês) \\
\end{longtable}

\subsubsection{Mercado Imobiliário e Terra
Rural}

\textbf{Fonte:} INDERT /\allowbreak  Clasificados.com.py (2024)

\begin{longtable}[]{@{}ll@{}}
\toprule\noalign{}
Tipo & Referência \\
\midrule\noalign{}
\endhead
\bottomrule\noalign{}
\endlastfoot
Terra agrícola alta prod. (USD/\allowbreak ha) & 14,800 \\
Imóvel urbano (USD/\allowbreak m²) & 1,400 \\
Aluguel 2 quartos (USD/\allowbreak mês) & 600 \\
\end{longtable}

\begin{quote}
Valores de referência departamental. Localidades menores podem ter
preços 20--40\% abaixo da capital departamental. \#\#\# Saúde
\end{quote}

\textbf{Fonte:} MSPBS /\allowbreak  IPS Paraguay (2024-2026), consolidação
departamental e proxy local

\begin{longtable}[]{@{}ll@{}}
\toprule\noalign{}
Serviço & Disponibilidade \\
\midrule\noalign{}
\endhead
\bottomrule\noalign{}
\endlastfoot
USF /\allowbreak  Posto de Saúde & sim \\
Hospital Regional & não \\
IPS (seguro social) & sim \\
Farmácia & sim \\
Distância ao hospital de referência & \textasciitilde42 km (Ciudad del
Este) \\
\end{longtable}

\textbf{Principais estabelecimentos:} USF local; referencia hospitalar
em Ciudad del Este

\textbf{Observação para imigrantes:} Atendimento primario local ou em
raio curto; casos de maior complexidade seguem para Ciudad del Este.
Cobertura privada continua recomendada para especialidades e urgências
de maior complexidade.
\newpage
\secaoDiagnostico{Santa Rita}{\mapaConteudo{-25.7833}{-55.0833}}{Santa Rita é o símbolo da potência agroindustrial paraguaia. Oferece
infraestrutura de serviços, segurança e produtividade que estão entre as
melhores do país. Sua força absoluta reside na qualidade mundial de seus
solos e na solidez de sua base econômica. É o local ideal para quem
busca integrar-se a um ambiente produtivo vibrante e seguro, mantendo-se
ciente da visibilidade estratégica que um polo de suprimentos deste
porte atrai em cenários de instabilidade nacional.

\subsubsection{Por que sim}

\begin{itemize}
\tightlist
\item
  Solo de altíssima produtividade (``terra roxa'') e abundância de
  recursos alimentares.
\item
  Ambiente social seguro e infraestrutura de serviços privados de alta
  qualidade.
\item
  Conectividade direta via Ruta PY06 para polos metropolitanos.
\item
  Sem restrições da Lei de Fronteira.
\end{itemize}

\subsubsection{Por que não}

\begin{itemize}
\tightlist
\item
  Alvo estratégico logístico e alimentar de alta visibilidade nacional.
\item
  Valor da terra entre os mais altos do Paraguai.
\item
  Vulnerabilidade social em cenários de interrupção drástica de
  suprimentos rodoviários.
\end{itemize}

\subsubsection{Recomendação
Técnica}

Altamente indicado para investimentos agroindustriais ou residências
rurais permanentes. Requer plano de autonomia energética e hídrica
individual para mitigar riscos de dependência de redes em crises. O GSS
de 6.7 reflete a solidez dos recursos equilibrada pela importância
estratégica da localidade. Referencia atual de score: GSS 6.7 em
2026-03-05.

\subsubsection{}

\begin{itemize}
\tightlist
\item
  \begin{enumerate}
  \def\labelenumi{\arabic{enumi})}
  \tightlist
  \item
    Dados censitários validados (INE\footnote{Instituto Nacional de Estadística (Instituto Nacional de Estatística), órgão oficial de dados demográficos e censitários.} 2022).
  \end{enumerate}
\item
  \begin{enumerate}
  \def\labelenumi{\arabic{enumi})}
  \setcounter{enumi}{1}
  \tightlist
  \item
    Relatórios de produtividade agrícola e eventos setoriais (Expo Santa
    Rita).
  \end{enumerate}
\item
  \begin{enumerate}
  \def\labelenumi{\arabic{enumi})}
  \setcounter{enumi}{2}
  \tightlist
  \item
    Mapa de solos e aptidão do Sul de Alto Paraná (MAG).
  \end{enumerate}
\item
  \begin{enumerate}
  \def\labelenumi{\arabic{enumi})}
  \setcounter{enumi}{3}
  \tightlist
  \item
    Registro de segurança pública e governança municipal próspera.
  \end{enumerate}
\end{itemize}}
\begin{table}[ht!]
\small \begin{tabular}{p{3.0cm}p{0.8cm}p{6.0cm}} \toprule
\textbf{Critério} & \textbf{Nota} & \textbf{Justificativa} \\ \midrule
A - Ameaças Estratégicas & 5.0 & Hub agroindustrial vital; alvo tático logístico de alta visibilidade e importância econômica \\
B - Risco Social & 6.0 & População regional e densidade urbana manejada; forte estabilidade social produtora \\
C - Riscos Naturais & 7.0 & Geologia muito estável e baixo risco de desastres naturais graves \\
D - Autossuficiência & 8.5 & Recursos agroindustriais excepcionais: capital da soja, proteína animal e solos de elite \\
E - Institucional & 7.5 & Forte segurança institucional e serviços comerciais/\allowbreak privados de primeiro nível \\
\midrule \textbf{GSS FINAL} & \textbf{6.7} & \textbf{Moderadamente Seguro (Localidade de elite para suporte logístico e produção de recursos, ciente da visibilidade estratégica).} \\ \bottomrule \end{tabular}
\end{table}
\subsubsection{Vulnerabilidades e
Mitigação}\label{vuln-Santa_Rita-vulnerabilidades-e-mitigauxe7uxe3o}

\begin{itemize}
\tightlist
\item
  \textbf{Visibilidade Tática:} Polo agroindustrial é um alvo
  estratégico para controle de suprimentos em crises (Mitigação:
  residência rural afastada 10-15 km do eixo urbano, mantendo autonomia
  total).
\item
  \textbf{Dependência de Exportação:} Economia vinculada ao mercado
  global de commodities (Mitigação: investimento em diversificação
  produtiva para consumo interno e energia solar off-grid).
\item
  \textbf{Gargalo Rodoviário:} Dependência da Ruta PY06 para escoamento
  e acesso a serviços complexos (Mitigação: mapeamento de rotas rurais
  vicinais em direção a Caazapá).
\end{itemize}
\subsubsection{Dossiê de Campo}\label{dossie-Santa_Rita-dossiuxea-de-campo}
\begin{description}[style=nextline,leftmargin=0.5cm,font=\bfseries]
\item[Coordenadas]  25°47'S 55°05'W.
\item[Topografia]  Relevo ondulado com extensas planícies de solos de elite. Altitude \textasciitilde 220m. Área de \textasciitilde 698 km². Geologicamente muito estável, risco sísmico praticamente nulo.
\item[Alvos Estratégicos]  Capital do Agronegócio de Alto Paraná (Maior polo de soja e milho do departamento); Grandes plantas de silos e moinhos; Nó logístico vital da Ruta PY06 (Eixo Encarnación-CDE); Subestação regional da ANDE\footnote{Administración Nacional de Electricidade (Administração Nacional de Eletricidade), companhia estatal de energia.}. Alta visibilidade econômica nacional.
\item[Fallout]  Ventos predominantes do Norte (NE) e Leste. Baixo risco de fallout direto.
\item[População]  27.249 habitantes (Censo 2022 Final). População cosmopolita com forte componente de imigrantes brasileiros e perfil produtor altamente dinâmico.
\item[Densidade]  \textasciitilde 39,0 hab/\allowbreak km² (Densidade moderada, com núcleo urbano comercial denso e colônias rurais altamente organizadas).
\item[Vias de Acesso]  Localização centralizada sobre a Ruta PY06 (Asfaltada). Conectividade de elite para Ciudad del Este (70 km) e Encarnación. Risco elevado de tornar-se ponto de controle tático e sofrer com o trânsito intenso em cenários de instabilidade nacional.
\item[Serviços]  Excelente infraestrutura de serviços privados. Sanatórios de médio porte, rede bancária completa e suporte técnico agroindustrial de primeiro nível. Dependência de CDE para procedimentos de altíssima complexidade.
\item[Custo de Vida]  Médio; economia dolarizada pelo agronegócio, com alta valorização imobiliária urbana e comercial.
\item[Preço da Terra]  US\$ 10.000-15.000/\allowbreak ha (terras mecanizadas de elite - "Terra Roxa").
\item[Hidrologia]  Baixo risco de inundações sistêmicas devido à drenagem natural eficiente do relevo ondulado. Presença de arroios perenes integrados à paisagem rural e urbana.
\item[Clima]  Subtropical Úmido. Verões quentes; invernos com geadas ocasionais. Precipitação média 1.700 mm/\allowbreak ano.
\item[Energia]  Rede nacional (Itaipu); boa estabilidade devido à importância industrial. Nodo de distribuição regional.
\item[Água]  Abundante via poços artesianos profundos e sistemas de junta de saneamento modernos.
\item[Qualidade do Solo]  Latossolo Vermelho de Elite; solo profundo e de altíssima produtividade mundial. Excelente aptidão para soja, milho, trigo e suinocultura intensiva.
\item[Resources Locais]  Grande polo produtor de grãos e proteína animal (suínos/\allowbreak aves). Elevadíssimo potencial para autossuficiência alimentar e hídrica regional, dependente de logística interna.
\item[Segurança]  Cidade próspera com segurança social superior à média nacional. Criminalidade urbana comum monitorada. Forte coesão comunitária baseada na cultura do agronegócio e cooperativismo. Presença de segurança patrimonial privada robusta.
\item[Leis Local]  Município consolidado e polo de influência ("Capital da Soja"). \\ \textbf{Livre de Restrição} de Fronteira. \\ \textit{Aquisição de terra rural proibida para estrangeiros limítrofes.}
\end{description}
\paragraph{Solo (SoilGrids 2.0, média ponderada 0--30
cm)}

\begin{longtable}[]{@{}ll@{}}
\toprule\noalign{}
Parâmetro & Valor \\
\midrule\noalign{}
\endhead
\bottomrule\noalign{}
\endlastfoot
pH (H$_2$O) & 5.4 \\
Carbono orgânico (SOC) & 21.78 g/\allowbreak kg \\
Argila & 45.75 \% \\
Areia & 20.8 \% \\
Silte (calc.) & 33.5 \% \\
Densidade aparente & 1.2 g/\allowbreak cm³ \\
\textbf{Aptidão agrícola} & \textbf{Média} \\
\end{longtable}

Fonte: ISRIC SoilGrids 2.0 via WCS (acesso em 2026-03-21). Coords:
-25.7833°, -55.0833°. Média ponderada camadas 0-5, 5-15, 15-30 cm.
\subsubsection{Indicadores Sociais}

\textbf{Fontes:} PNUD 2020, INE\footnote{Instituto Nacional de Estadística (Instituto Nacional de Estatística), órgão oficial de dados demográficos e censitários.} Censo 2022, INE\footnote{Instituto Nacional de Estadística (Instituto Nacional de Estatística), órgão oficial de dados demográficos e censitários.} EPHC 2023, Ministerio
Público 2024\\
\textbf{Nota:} Valores marcados como \emph{(dept.)} referem-se ao
departamento de Alto Paraná; valores \emph{(dist.)} são específicos
deste distrito. Dados marcados \emph{(est.)} são estimativas.

\begin{longtable}[]{@{}
  >{\raggedright\arraybackslash}p{(\linewidth - 4\tabcolsep) * \real{0.4231}}
  >{\raggedright\arraybackslash}p{(\linewidth - 4\tabcolsep) * \real{0.2692}}
  >{\raggedright\arraybackslash}p{(\linewidth - 4\tabcolsep) * \real{0.3077}}@{}}
\toprule\noalign{}
\begin{minipage}[b]{\linewidth}\raggedright
Indicador
\end{minipage} & \begin{minipage}[b]{\linewidth}\raggedright
Valor
\end{minipage} & \begin{minipage}[b]{\linewidth}\raggedright
Âmbito
\end{minipage} \\
\midrule\noalign{}
\endhead
\bottomrule\noalign{}
\endlastfoot
IDH (2020) & 0,752 & dept. \\
Ranking IDH nacional & 3/\allowbreak 18 & dept. \\
Esperança de vida & 75,1 anos & dept. \\
Escolaridade média & 10.0 anos & dist. \\
RNB per capita (USD PPA) & 13.500 & dept. \\
Pobreza monetária (\%) & 16,3\% & dept. \\
Pobreza extrema (\%) & 3,1\% & dept. \\
Índice de Gini & 0,428 & dept. \\
Acesso a água potável (\%) & 88,7\% & dept. \\
Acesso a saneamento (\%) & 87,2\% & dept. \\
Taxa de homicídios (est., /\allowbreak 100k hab) & \textasciitilde8--12 (estimativa,
acima da média nacional) & dept. \\
Índice de segurança & baixo & dept. \\
População (Censo 2022) & 27.249 hab. & dist. \\
Idade mediana & 28.0 anos & dist. \\
Taxa de fecundidade & N/\allowbreak D & dist. \\
\end{longtable}
\subsubsection{Combustível}

\textbf{Referência:} PETROPAR /\allowbreak  postos locais (2024) \textbf{Tipo de
localidade:} Interior

\begin{longtable}[]{@{}lll@{}}
\toprule\noalign{}
Combustível & USD/\allowbreak litro & Gs/\allowbreak litro (aprox.) \\
\midrule\noalign{}
\endhead
\bottomrule\noalign{}
\endlastfoot
Gasolina 93 oct & 0.97 & 7,178 \\
Gasolina 97 oct (premium) & 1.07 & 7,918 \\
Diesel & 0.90 & 6,660 \\
\end{longtable}

\begin{quote}
Preços podem variar ±5\% conforme posto e sazonalidade. Chaco e interior
remoto apresentam maior variação.
\end{quote}

\subsubsection{Cobertura Celular}

\textbf{Fonte:} CONATEL PY /\allowbreak  operadoras (2024)

\begin{longtable}[]{@{}ll@{}}
\toprule\noalign{}
Parâmetro & Valor \\
\midrule\noalign{}
\endhead
\bottomrule\noalign{}
\endlastfoot
Cobertura 4G população (dept.) & 95\% \\
Cobertura 4G área rural & 88\% \\
Melhor operadora & Tigo/\allowbreak Personal \\
Qualidade rural & boa \\
\end{longtable}

\begin{quote}
Para áreas rurais fora do núcleo urbano, recomenda-se chip Tigo como
principal e Personal como backup.
\end{quote}

\subsubsection{Internet}

\textbf{Fonte:} CONATEL /\allowbreak  Speedtest Ookla (2024)

\begin{longtable}[]{@{}ll@{}}
\toprule\noalign{}
Parâmetro & Valor \\
\midrule\noalign{}
\endhead
\bottomrule\noalign{}
\endlastfoot
Velocidade média download & 85 Mbps \\
Domicílios com internet (dept.) & 87\% \\
Tecnologia predominante & fibra/\allowbreak cabo \\
Opção rural & Starlink disponível (\textasciitilde USD 44/\allowbreak mês) \\
\end{longtable}

\subsubsection{Mercado Imobiliário e Terra
Rural}

\textbf{Fonte:} INDERT /\allowbreak  Clasificados.com.py (2024)

\begin{longtable}[]{@{}ll@{}}
\toprule\noalign{}
Tipo & Referência \\
\midrule\noalign{}
\endhead
\bottomrule\noalign{}
\endlastfoot
Terra agrícola alta prod. (USD/\allowbreak ha) & 14,800 \\
Imóvel urbano (USD/\allowbreak m²) & 1,400 \\
Aluguel 2 quartos (USD/\allowbreak mês) & 600 \\
\end{longtable}

\begin{quote}
Valores de referência departamental. Localidades menores podem ter
preços 20--40\% abaixo da capital departamental. \#\#\# Saúde
\end{quote}

\textbf{Fonte:} MSPBS /\allowbreak  IPS Paraguay (2024-2026), consolidação
departamental e proxy local

\begin{longtable}[]{@{}ll@{}}
\toprule\noalign{}
Serviço & Disponibilidade \\
\midrule\noalign{}
\endhead
\bottomrule\noalign{}
\endlastfoot
USF /\allowbreak  Posto de Saúde & sim \\
Hospital Regional & não \\
IPS (seguro social) & sim \\
Farmácia & sim \\
Distância ao hospital de referência & \textasciitilde72 km (Ciudad del
Este) \\
\end{longtable}

\textbf{Principais estabelecimentos:} USF local; referencia hospitalar
em Ciudad del Este

\textbf{Observação para imigrantes:} Atendimento primario local ou em
raio curto; casos de maior complexidade seguem para Ciudad del Este.
Cobertura privada continua recomendada para especialidades e urgências
de maior complexidade.
\newpage
\secaoDiagnostico{Santa Rosa del Monday}{\mapaConteudo{-25.8833}{-54.9666}}{Santa Rosa del Monday é a fortaleza agroindustrial de elite do Paraguai.
Oferece solos de produtividade mundial, abundância massiva de recursos
hídricos e alimentares, e uma segurança social sustentada por uma
comunidade de produtores altamente organizada e próspera. Sua força
absoluta reside na resiliência de sua base produtiva e na baixíssima
densidade populacional, garantindo privacidade e isolamento tático
superior. É o local ideal para quem busca autonomia total em ambiente de
alta performance, transformando a riqueza de recursos básicos em uma
armadura definitiva de proteção econômica e alimentar contra colapsos
sistêmicos externos.

\subsubsection{Por que sim}

\begin{itemize}
\tightlist
\item
  Localidade de elite no Ranking Nacional de Segurança Estratégica.
\item
  Solos de altíssima produtividade mundial (``terra roxa'') e clima
  favorável.
\item
  Ambiente social extremamente seguro, ordeiro e resiliente.
\item
  Isolamento demográfico rural de alto nível (privacidade absoluta).
\end{itemize}

\subsubsection{Por que não}

\begin{itemize}
\tightlist
\item
  Alvo estratégico prioritário para controle nacional de alimentos em
  crises extremas.
\item
  Valor da terra entre os mais elevados do país.
\item
  Dependência tática de insumos externos para manutenção da escala
  mecanizada.
\end{itemize}

\subsubsection{Recomendação
Técnica}

Altamente indicado para estabelecimentos de autossuficiência tecnificada
de alto padrão ou investimentos produtivos. Requer plano de autonomia
hídrica e energética individual para mitigar riscos de rede. O GSS de
8.2 reflete a excepcional força dos recursos e da segurança
institucional da localidade. Referencia atual de score: GSS 8.2 em
2026-03-05.

\subsubsection{}

\begin{itemize}
\tightlist
\item
  \begin{enumerate}
  \def\labelenumi{\arabic{enumi})}
  \tightlist
  \item
    Dados censitários validados (INE\footnote{Instituto Nacional de Estadística (Instituto Nacional de Estatística), órgão oficial de dados demográficos e censitários.} 2022).
  \end{enumerate}
\item
  \begin{enumerate}
  \def\labelenumi{\arabic{enumi})}
  \setcounter{enumi}{1}
  \tightlist
  \item
    Relatórios de produção agrícola mecanizada de alta precisão (MAG).
  \end{enumerate}
\item
  \begin{enumerate}
  \def\labelenumi{\arabic{enumi})}
  \setcounter{enumi}{2}
  \tightlist
  \item
    Mapa de solos e bacias hídricas do sul de Alto Paraná (MAG).
  \end{enumerate}
\item
  \begin{enumerate}
  \def\labelenumi{\arabic{enumi})}
  \setcounter{enumi}{3}
  \tightlist
  \item
    Registro de segurança pública e governança cooperativa regional.
  \end{enumerate}
\end{itemize}}
\begin{table}[ht!]
\small \begin{tabular}{p{3.0cm}p{0.8cm}p{6.0cm}} \toprule
\textbf{Critério} & \textbf{Nota} & \textbf{Justificativa} \\ \midrule
A - Ameaças Estratégicas & 7.5 & Hub agroindustrial de elite; isolamento tático superior fora dos eixos de conflito, visibilidade como polo de recursos \\
B - Risco Social & 8.5 & Densidade populacional extremamente baixa; isolamento social de alto nível em ambiente rural organizado \\
C - Riscos Naturais & 7.5 & Geologia muito estável e topografia segura contra desastres naturais \\
D - Autossuficiência & 9.5 & Recursos agrícolas e hídricos excepcionais de nível mundial; polo de segurança alimentar \\
E - Institucional & 8.0 & Forte segurança institucional privada e serviços de primeiro nível sustentados pelo agronegócio \\
\midrule \textbf{GSS FINAL} & \textbf{8.2} & \textbf{Seguro (Localidade de elite absoluta para autossuficiência e residência tática, oferecendo o máximo em recursos e segurança social rural).} \\ \bottomrule \end{tabular}
\end{table}
\subsubsection{Vulnerabilidades e
Mitigação}\label{vuln-Santa_Rosa_del_Monday-vulnerabilidades-e-mitigauxe7uxe3o}

\begin{itemize}
\tightlist
\item
  \textbf{Visibilidade Tática de Suprimentos:} Sendo um polo de elite de
  grãos, é um alvo estratégico para controle de suprimentos em crises
  nacionais (Mitigação: residência rural afastada 5-10 km do núcleo
  urbano industrial e integração com as redes de defesa das
  cooperativas).
\item
  \textbf{Dependência de Exportação:} Economia vinculada ao mercado
  global (Mitigação: diversificação produtiva para consumo interno e
  autonomia energética solar total).
\item
  \textbf{Escala Comercial Especializada:} Dependência de grandes
  centros para insumos médicos de altíssima complexidade (Mitigação:
  manutenção de estoque médico avançado e kit trauma local).
\end{itemize}
\subsubsection{Dossiê de Campo}\label{dossie-Santa_Rosa_del_Monday-dossiuxea-de-campo}
\begin{description}[style=nextline,leftmargin=0.5cm,font=\bfseries]
\item[Coordenadas]  25°53'S 54°58'W.
\item[Topografia]  Relevo predominantemente plano a suavemente ondulado, com solos de altíssima produtividade (Terra Roxa mecanizada em 98\% da área). Altitude \textasciitilde 240m. Área de \textasciitilde 500 km². Geologicamente muito estável, risco sísmico nulo.
\item[Alvos Estratégicos]  Capital Nacional da Soja (Um dos maiores polos produtores de grãos por metro quadrado do país); Grandes silos e plantas de beneficiamento da Cooperativa Lar e Cooperativa Raul Peña; Proximidade tática com o Rio Monday. Ausência de alvos militares diretos, oferecendo excelente invisibilidade estratégica, sendo um centro de gravidade alimentar nacional.
\item[Fallout]  Ventos predominantes do Norte (NE) e Leste. Baixíssimo risco de fallout direto; posição isolada no sul do departamento, protegida da massa urbana de CDE.
\item[População]  5.923 habitantes (Censo 2022 Final). População composta por produtores rurais altamente tecnificados, com forte componente demográfico de imigrantes brasileiros e europeus.
\item[Densidade]  \textasciitilde 11,8 hab/\allowbreak km² (Densidade extremamente baixa, garantindo privacidade absoluta e ausência de riscos sociais de massa urbana).
\item[Vias de Acesso]  Acesso via ramais pavimentados de alta qualidade que conectam à Ruta PY06. Localização que favorece a discrição, funcionando como um enclave de riqueza e tranquilidade rural fora dos eixos de trânsito internacional pesado.
\item[Serviços]  Infraestrutura de serviços privados de primeiro nível. Sanatórios modernos, suporte técnico de precisão e rede bancária consolidada sustentada pelo agronegócio.
\item[Custo de Vida]  Médio-alto; economia dolarizada pelo agronegócio e alta valorização imobiliária das terras mecanizadas.
\item[Preço da Terra]  US\$ 15.000-25.000/\allowbreak ha (terras mecanizadas de elite mundial - "Terra Roxa").
\item[Hidrologia]  Baixo risco de inundações sistêmicas devido à topografia favorável e gestão hídrica de precisão nas fazendas. Presença de arroios perenes tributários do Rio Monday.
\item[Clima]  Subtropical Úmido. Verões quentes; invernos frescos com geadas benéficas ao trigo.
\item[Energia]  Rede nacional (Itaipu); boa estabilidade devido à importância das agroindústrias locais e cooperativas. Uso extensivo de sistemas solares e geradores potentes em propriedades de elite.
\item[Água]  Abundante via poços artesianos profundos e lençol freático produtivo; excelente qualidade hídrica subterrânea (Aquífero Guarani).
\item[Qualidade do Solo]  Latossolo Vermelho de Elite; solo profundo, estruturado e de fertilidade máxima mundial. Base da potência agrícola regional.
\item[Resources Locais]  Polo de Segurança Alimentar de Elite. Imensa produção de soja, milho (incluindo variedades GSS de alto rendimento), trigo e proteína animal (suínos/\allowbreak aves). Elevadíssimo potencial para autossuficiência absoluta em nível de propriedade ou regional.
\item[Segurança]  Uma das zonas rurais mais seguras e estáveis do Paraguai. Baixíssima criminalidade urbana. Coesão social fortíssima baseada na cultura cooperativista e produtora. Presença de segurança patrimonial privada de alta performance.
\item[Leis Local]  Município consolidado e polo de referência tecnológica ("Cidade da Soja"). \\ \textbf{Livre de Restrição} de Fronteira. \\ \textit{Aquisição de terra rural proibida para estrangeiros limítrofes.}
\end{description}
\paragraph{Solo (SoilGrids 2.0, média ponderada 0--30
cm)}

\begin{longtable}[]{@{}ll@{}}
\toprule\noalign{}
Parâmetro & Valor \\
\midrule\noalign{}
\endhead
\bottomrule\noalign{}
\endlastfoot
pH (H$_2$O) & 5.4 \\
Carbono orgânico (SOC) & 26.74 g/\allowbreak kg \\
Argila & 45.46 \% \\
Areia & 20.18 \% \\
Silte (calc.) & 34.4 \% \\
Densidade aparente & 1.17 g/\allowbreak cm³ \\
\textbf{Aptidão agrícola} & \textbf{Média} \\
\end{longtable}

Fonte: ISRIC SoilGrids 2.0 via WCS (acesso em 2026-03-21). Coords:
-25.8833°, -54.9667°. Média ponderada camadas 0-5, 5-15, 15-30 cm.

\begin{itemize}
\tightlist
\item
  \textbf{Homicidios/\allowbreak 100k:} \textasciitilde4,0 (Recorde de paz rural).
\end{itemize}

\subsection{10. Analise de Riscos}

\begin{itemize}
\tightlist
\item
  \textbf{Cenario curto prazo:} Manutenção da prosperidade econômica e
  expansão das agroindústrias mecanizadas.
\item
  \textbf{Cenario medio prazo:} Impacto de mudanças climáticas na
  frequência de secas e resiliência das variedades GSS de milho.
\end{itemize}

\subsection{11. Redacao Final}

\subsubsection{Diagnostico Integrado}

Santa Rosa del Monday é a fortaleza agroindustrial de elite do Paraguai.
Oferece solos de produtividade mundial, abundância massiva de recursos
hídricos e alimentares, e uma segurança social sustentada por uma
comunidade de produtores altamente organizada e próspera. Sua força
absoluta reside na resiliência de sua base produtiva e na baixíssima
densidade populacional, garantindo privacidade e isolamento tático
superior. É o local ideal para quem busca autonomia total em ambiente de
alta performance, transformando a riqueza de recursos básicos em uma
armadura definitiva de proteção econômica e alimentar contra colapsos
sistêmicos externos.

\subsubsection{Por que sim}

\begin{itemize}
\tightlist
\item
  Localidade de elite no Ranking Nacional de Segurança Estratégica.
\item
  Solos de altíssima produtividade mundial (``terra roxa'') e clima
  favorável.
\item
  Ambiente social extremamente seguro, ordeiro e resiliente.
\item
  Isolamento demográfico rural de alto nível (privacidade absoluta).
\end{itemize}

\subsubsection{Por que nao}

\begin{itemize}
\tightlist
\item
  Alvo estratégico prioritário para controle nacional de alimentos em
  crises extremas.
\item
  Valor da terra entre os mais elevados do país.
\item
  Dependência tática de insumos externos para manutenção da escala
  mecanizada.
\end{itemize}

\subsubsection{Recomendacao Tecnica}

Altamente indicado para estabelecimentos de autossuficiência tecnificada
de alto padrão ou investimentos produtivos. Requer plano de autonomia
hídrica e energética individual para mitigar riscos de rede. O GSS de
8.2 reflete a excepcional força dos recursos e da segurança
institucional da localidade. Referencia atual de score: GSS 8.2 em
2026-03-05.

\subsection{13.}

\begin{itemize}
\tightlist
\item
  \begin{enumerate}
  \def\labelenumi{\arabic{enumi})}
  \tightlist
  \item
    Dados censitários validados (INE\footnote{Instituto Nacional de Estadística (Instituto Nacional de Estatística), órgão oficial de dados demográficos e censitários.} 2022).
  \end{enumerate}
\item
  \begin{enumerate}
  \def\labelenumi{\arabic{enumi})}
  \setcounter{enumi}{1}
  \tightlist
  \item
    Relatórios de produção agrícola mecanizada de alta precisão (MAG).
  \end{enumerate}
\item
  \begin{enumerate}
  \def\labelenumi{\arabic{enumi})}
  \setcounter{enumi}{2}
  \tightlist
  \item
    Mapa de solos e bacias hídricas do sul de Alto Paraná (MAG).
  \end{enumerate}
\item
  \begin{enumerate}
  \def\labelenumi{\arabic{enumi})}
  \setcounter{enumi}{3}
  \tightlist
  \item
    Registro de segurança pública e governança cooperativa regional.
  \end{enumerate}
\end{itemize}

\subsubsection{Combustível}

\textbf{Referência:} PETROPAR /\allowbreak  postos locais (2024) \textbf{Tipo de
localidade:} Interior

\begin{longtable}[]{@{}lll@{}}
\toprule\noalign{}
Combustível & USD/\allowbreak litro & Gs/\allowbreak litro (aprox.) \\
\midrule\noalign{}
\endhead
\bottomrule\noalign{}
\endlastfoot
Gasolina 93 oct & 0.97 & 7,178 \\
Gasolina 97 oct (premium) & 1.07 & 7,918 \\
Diesel & 0.90 & 6,660 \\
\end{longtable}

\begin{quote}
Preços podem variar ±5\% conforme posto e sazonalidade. Chaco e interior
remoto apresentam maior variação.
\end{quote}

\subsubsection{Cobertura Celular}

\textbf{Fonte:} CONATEL PY /\allowbreak  operadoras (2024)

\begin{longtable}[]{@{}ll@{}}
\toprule\noalign{}
Parâmetro & Valor \\
\midrule\noalign{}
\endhead
\bottomrule\noalign{}
\endlastfoot
Cobertura 4G população (dept.) & 95\% \\
Cobertura 4G área rural & 88\% \\
Melhor operadora & Tigo/\allowbreak Personal \\
Qualidade rural & boa \\
\end{longtable}

\begin{quote}
Para áreas rurais fora do núcleo urbano, recomenda-se chip Tigo como
principal e Personal como backup.
\end{quote}

\subsubsection{Internet}

\textbf{Fonte:} CONATEL /\allowbreak  Speedtest Ookla (2024)

\begin{longtable}[]{@{}ll@{}}
\toprule\noalign{}
Parâmetro & Valor \\
\midrule\noalign{}
\endhead
\bottomrule\noalign{}
\endlastfoot
Velocidade média download & 85 Mbps \\
Domicílios com internet (dept.) & 87\% \\
Tecnologia predominante & fibra/\allowbreak cabo \\
Opção rural & Starlink disponível (\textasciitilde USD 44/\allowbreak mês) \\
\end{longtable}

\subsubsection{Mercado Imobiliário e Terra
Rural}

\textbf{Fonte:} INDERT /\allowbreak  Clasificados.com.py (2024)

\begin{longtable}[]{@{}ll@{}}
\toprule\noalign{}
Tipo & Referência \\
\midrule\noalign{}
\endhead
\bottomrule\noalign{}
\endlastfoot
Terra agrícola alta prod. (USD/\allowbreak ha) & 14,800 \\
Imóvel urbano (USD/\allowbreak m²) & 1,400 \\
Aluguel 2 quartos (USD/\allowbreak mês) & 600 \\
\end{longtable}

\begin{quote}
Valores de referência departamental. Localidades menores podem ter
preços 20--40\% abaixo da capital departamental. \#\#\# Saúde
\end{quote}

\textbf{Fonte:} MSPBS /\allowbreak  IPS Paraguay (2024-2026), consolidação
departamental e proxy local

\begin{longtable}[]{@{}ll@{}}
\toprule\noalign{}
Serviço & Disponibilidade \\
\midrule\noalign{}
\endhead
\bottomrule\noalign{}
\endlastfoot
USF /\allowbreak  Posto de Saúde & sim \\
Hospital Regional & não \\
IPS (seguro social) & sim \\
Farmácia & sim \\
Distância ao hospital de referência & \textasciitilde55 km (Ciudad del
Este) \\
\end{longtable}

\textbf{Principais estabelecimentos:} USF local; referencia hospitalar
em Ciudad del Este

\textbf{Observação para imigrantes:} Atendimento primario local ou em
raio curto; casos de maior complexidade seguem para Ciudad del Este.
Cobertura privada continua recomendada para especialidades e urgências
de maior complexidade.
\subsubsection{Indicadores Sociais}

\textbf{Fontes:} PNUD 2020, INE\footnote{Instituto Nacional de Estadística (Instituto Nacional de Estatística), órgão oficial de dados demográficos e censitários.} Censo 2022, INE\footnote{Instituto Nacional de Estadística (Instituto Nacional de Estatística), órgão oficial de dados demográficos e censitários.} EPHC 2023, Ministerio
Público 2024\\
\textbf{Nota:} Valores marcados como \emph{(dept.)} referem-se ao
departamento de Alto Paraná; valores \emph{(dist.)} são específicos
deste distrito. Dados marcados \emph{(est.)} são estimativas.

\begin{longtable}[]{@{}
  >{\raggedright\arraybackslash}p{(\linewidth - 4\tabcolsep) * \real{0.4231}}
  >{\raggedright\arraybackslash}p{(\linewidth - 4\tabcolsep) * \real{0.2692}}
  >{\raggedright\arraybackslash}p{(\linewidth - 4\tabcolsep) * \real{0.3077}}@{}}
\toprule\noalign{}
\begin{minipage}[b]{\linewidth}\raggedright
Indicador
\end{minipage} & \begin{minipage}[b]{\linewidth}\raggedright
Valor
\end{minipage} & \begin{minipage}[b]{\linewidth}\raggedright
Âmbito
\end{minipage} \\
\midrule\noalign{}
\endhead
\bottomrule\noalign{}
\endlastfoot
IDH (2020) & 0,752 & dept. \\
Ranking IDH nacional & 3/\allowbreak 18 & dept. \\
Esperança de vida & 75,1 anos & dept. \\
Escolaridade média & 8.4 anos & dist. \\
RNB per capita (USD PPA) & 13.500 & dept. \\
Pobreza monetária (\%) & 16,3\% & dept. \\
Pobreza extrema (\%) & 3,1\% & dept. \\
Índice de Gini & 0,428 & dept. \\
Acesso a água potável (\%) & 88,7\% & dept. \\
Acesso a saneamento (\%) & 87,2\% & dept. \\
Taxa de homicídios (est., /\allowbreak 100k hab) & \textasciitilde8--12 (estimativa,
acima da média nacional) & dept. \\
Índice de segurança & baixo & dept. \\
População (Censo 2022) & 5.923 hab. & dist. \\
Idade mediana & 30.0 anos & dist. \\
Taxa de fecundidade & N/\allowbreak D & dist. \\
\end{longtable}
\subsubsection{Combustível}

\textbf{Referência:} PETROPAR /\allowbreak  postos locais (2024) \textbf{Tipo de
localidade:} Interior

\begin{longtable}[]{@{}lll@{}}
\toprule\noalign{}
Combustível & USD/\allowbreak litro & Gs/\allowbreak litro (aprox.) \\
\midrule\noalign{}
\endhead
\bottomrule\noalign{}
\endlastfoot
Gasolina 93 oct & 0.97 & 7,178 \\
Gasolina 97 oct (premium) & 1.07 & 7,918 \\
Diesel & 0.90 & 6,660 \\
\end{longtable}

\begin{quote}
Preços podem variar ±5\% conforme posto e sazonalidade. Chaco e interior
remoto apresentam maior variação.
\end{quote}

\subsubsection{Cobertura Celular}

\textbf{Fonte:} CONATEL PY /\allowbreak  operadoras (2024)

\begin{longtable}[]{@{}ll@{}}
\toprule\noalign{}
Parâmetro & Valor \\
\midrule\noalign{}
\endhead
\bottomrule\noalign{}
\endlastfoot
Cobertura 4G população (dept.) & 95\% \\
Cobertura 4G área rural & 88\% \\
Melhor operadora & Tigo/\allowbreak Personal \\
Qualidade rural & boa \\
\end{longtable}

\begin{quote}
Para áreas rurais fora do núcleo urbano, recomenda-se chip Tigo como
principal e Personal como backup.
\end{quote}

\subsubsection{Internet}

\textbf{Fonte:} CONATEL /\allowbreak  Speedtest Ookla (2024)

\begin{longtable}[]{@{}ll@{}}
\toprule\noalign{}
Parâmetro & Valor \\
\midrule\noalign{}
\endhead
\bottomrule\noalign{}
\endlastfoot
Velocidade média download & 85 Mbps \\
Domicílios com internet (dept.) & 87\% \\
Tecnologia predominante & fibra/\allowbreak cabo \\
Opção rural & Starlink disponível (\textasciitilde USD 44/\allowbreak mês) \\
\end{longtable}

\subsubsection{Mercado Imobiliário e Terra
Rural}

\textbf{Fonte:} INDERT /\allowbreak  Clasificados.com.py (2024)

\begin{longtable}[]{@{}ll@{}}
\toprule\noalign{}
Tipo & Referência \\
\midrule\noalign{}
\endhead
\bottomrule\noalign{}
\endlastfoot
Terra agrícola alta prod. (USD/\allowbreak ha) & 14,800 \\
Imóvel urbano (USD/\allowbreak m²) & 1,400 \\
Aluguel 2 quartos (USD/\allowbreak mês) & 600 \\
\end{longtable}

\begin{quote}
Valores de referência departamental. Localidades menores podem ter
preços 20--40\% abaixo da capital departamental. \#\#\# Saúde
\end{quote}

\textbf{Fonte:} MSPBS /\allowbreak  IPS Paraguay (2024-2026), consolidação
departamental e proxy local

\begin{longtable}[]{@{}ll@{}}
\toprule\noalign{}
Serviço & Disponibilidade \\
\midrule\noalign{}
\endhead
\bottomrule\noalign{}
\endlastfoot
USF /\allowbreak  Posto de Saúde & sim \\
Hospital Regional & não \\
IPS (seguro social) & sim \\
Farmácia & sim \\
Distância ao hospital de referência & \textasciitilde55 km (Ciudad del
Este) \\
\end{longtable}

\textbf{Principais estabelecimentos:} USF local; referencia hospitalar
em Ciudad del Este

\textbf{Observação para imigrantes:} Atendimento primario local ou em
raio curto; casos de maior complexidade seguem para Ciudad del Este.
Cobertura privada continua recomendada para especialidades e urgências
de maior complexidade.
\newpage
\secaoDiagnostico{Tavapy}{\mapaConteudo{-25.6666}{-55.0000}}{Tavapy (Alto Parana) apresenta perfil operacional consolidado para
comparacao territorial, com leitura conjunta de infraestrutura, risco
hidroclimatico e capacidade de continuidade de servicos essenciais.

\subsubsection{Por que sim}

\begin{itemize}
\tightlist
\item
  Base institucional de referencia disponivel e rastreavel.
\item
  Estrutura comparativa (A-E/\allowbreak GSS) consistente para decisao relativa.
\item
  Plano de mitigacao acionavel ja definido no dossie.
\end{itemize}

\subsubsection{Por que não}

\begin{itemize}
\tightlist
\item
  Serie oficial distrital granular ainda pode apresentar lacunas
  temporais.
\item
  Sensibilidade do score exige monitoramento continuo de risco e
  conectividade.
\item
  Decisao final depende de validacao de campo e janela temporal recente.
\end{itemize}

\subsubsection{Pre-condicoes Minimas de
Decisao}

\begin{itemize}
\tightlist
\item
  Atualizar indicadores criticos em ciclo trimestral.
\item
  Confirmar trafegabilidade e contingencia hidrica/\allowbreak energetica local.
\item
  Validar checklist de resiliencia domiciliar/\allowbreak logistica antes de decisao
  final.
\end{itemize}

\subsubsection{Recomendação
Técnica}

Uso recomendado como candidato em analise multicriterio, condicionado ao
cumprimento das pre-condicoes acima. Referencia atual de score: GSS 7.3
em 2026-03-05; risco predominante: baixo.

\subsubsection{}

\begin{itemize}
\tightlist
\item
  \begin{enumerate}
  \def\labelenumi{\arabic{enumi})}
  \tightlist
  \item
    Delimitacao territorial validada por georreferencia institucional
    (Fonte: acesso: 2026-03-05).
  \end{enumerate}
\item
  \begin{enumerate}
  \def\labelenumi{\arabic{enumi})}
  \setcounter{enumi}{1}
  \tightlist
  \item
    População municipal/\allowbreak distrital referenciada por base censitaria
    nacional (Fonte: acesso: 2026-03-05).
  \end{enumerate}
\item
  \begin{enumerate}
  \def\labelenumi{\arabic{enumi})}
  \setcounter{enumi}{2}
  \tightlist
  \item
    Comparabilidade intra-departamental apoiada em indicadores
    distritais oficiais (Fonte: acesso: 2026-03-05).
  \end{enumerate}
\item
  \begin{enumerate}
  \def\labelenumi{\arabic{enumi})}
  \setcounter{enumi}{3}
  \tightlist
  \item
    Estado macro de conectividade viaria ancorado em informacoes de
    rodovias (Fonte: acesso: 2026-03-05).
  \end{enumerate}
\item
  \begin{enumerate}
  \def\labelenumi{\arabic{enumi})}
  \setcounter{enumi}{4}
  \tightlist
  \item
    Exposicao hidrometeorologica monitorada por avisos oficiais (Fonte:
    acesso: 2026-03-05).
  \end{enumerate}
\item
  \begin{enumerate}
  \def\labelenumi{\arabic{enumi})}
  \setcounter{enumi}{5}
  \tightlist
  \item
    Historico de contexto meteorologico apoiado em publicacoes tecnicas
    (Fonte: acesso: 2026-03-05).
  \end{enumerate}
\item
  \begin{enumerate}
  \def\labelenumi{\arabic{enumi})}
  \setcounter{enumi}{6}
  \tightlist
  \item
    Capacidade institucional de resposta a eventos apoiada em acoes da
    SEN\footnote{Secretaría de Emergencia Nacional (Secretaria de Emergência Nacional), órgão governamental de resposta a desastres.} (Fonte: acesso: 2026-03-05).
  \end{enumerate}
\item
  \begin{enumerate}
  \def\labelenumi{\arabic{enumi})}
  \setcounter{enumi}{7}
  \tightlist
  \item
    Infraestrutura energetica analisada via operador nacional (Fonte:
    acesso: 2026-03-05).
  \end{enumerate}
\item
  \begin{enumerate}
  \def\labelenumi{\arabic{enumi})}
  \setcounter{enumi}{8}
  \tightlist
  \item
    Referencias de obras e logistica nacional verificadas no MOPC
    (Fonte: acesso: 2026-03-05).
  \end{enumerate}
\item
  \begin{enumerate}
  \def\labelenumi{\arabic{enumi})}
  \setcounter{enumi}{9}
  \tightlist
  \item
    Contexto sociopolitico institucional referenciado por autoridade
    eleitoral (Fonte: acesso: 2026-03-05).
  \end{enumerate}
\item
  \begin{enumerate}
  \def\labelenumi{\arabic{enumi})}
  \setcounter{enumi}{10}
  \tightlist
  \item
    Camada cartografica de apoio eleitoral/\allowbreak territorial consultada
    (Fonte: acesso: 2026-03-05).
  \end{enumerate}
\item
  \begin{enumerate}
  \def\labelenumi{\arabic{enumi})}
  \setcounter{enumi}{11}
  \tightlist
  \item
    Dados publicos complementares para verificacao cruzada (Fonte:
    acesso: 2026-03-05).
  \end{enumerate}
\end{itemize}}
\begin{table}[ht!]
\small \begin{tabular}{p{3.0cm}p{0.8cm}p{6.0cm}} \toprule
\textbf{Critério} & \textbf{Nota} & \textbf{Justificativa} \\ \midrule
A - Ameaças Estratégicas & 7.8 & - \\
B - Risco Social & 7.2 & - \\
C - Riscos Naturais & 6.8 & - \\
D - Autossuficiência & 6.9 & - \\
E - Institucional & 7.6 & - \\
\midrule \textbf{GSS FINAL} & \textbf{7.3} & \textbf{Seguro (bom potencial com preparacao).} \\ \bottomrule \end{tabular}
\end{table}
\subsubsection{Vulnerabilidades e
Mitigação}\label{vuln-Tavapy-vulnerabilidades-e-mitigauxe7uxe3o}

\begin{itemize}
\tightlist
\item
  Exposicao hidrometeorologica controlada (\textless45/\allowbreak 100).
\item
  Conectividade viaria intermediaria (50-69/\allowbreak 100).
\item
  Estoque de contingencia para 14 dias (agua/\allowbreak alimentos/\allowbreak combustivel),
  priorizado quando risco hidro=40/\allowbreak 100.
\item
  Duplicar rotas/\allowbreak logistica quando conectividade viaria=67/\allowbreak 100 for
  \textless70; manter rota secundaria mapeada e testada trimestralmente.
\item
  Plano de energia resiliente com autonomia minima de 72h
  (geracao/\allowbreak backup), revisao semestral.
\end{itemize}
\subsubsection{Dossiê de Campo}\label{dossie-Tavapy-dossiuxea-de-campo}
\begin{description}[style=nextline,leftmargin=0.5cm,font=\bfseries]
\item[Coordenadas]  25°40'S 55°00'W (geocodificado via Nominatim).
\end{description}
\paragraph{Solo (SoilGrids 2.0, média ponderada 0--30
cm)}

\begin{longtable}[]{@{}ll@{}}
\toprule\noalign{}
Parâmetro & Valor \\
\midrule\noalign{}
\endhead
\bottomrule\noalign{}
\endlastfoot
pH (H$_2$O) & 5.4 \\
Carbono orgânico (SOC) & 21.24 g/\allowbreak kg \\
Argila & 40.53 \% \\
Areia & 23.51 \% \\
Silte (calc.) & 36.0 \% \\
Densidade aparente & 1.26 g/\allowbreak cm³ \\
\textbf{Aptidão agrícola} & \textbf{Média} \\
\end{longtable}

Fonte: ISRIC SoilGrids 2.0 via WCS (acesso em 2026-03-21). Coords:
-25.6667°, -55.0°. Média ponderada camadas 0-5, 5-15, 15-30 cm.
\subsubsection{Indicadores Sociais}

\textbf{Fontes:} PNUD 2020, INE\footnote{Instituto Nacional de Estadística (Instituto Nacional de Estatística), órgão oficial de dados demográficos e censitários.} Censo 2022, INE\footnote{Instituto Nacional de Estadística (Instituto Nacional de Estatística), órgão oficial de dados demográficos e censitários.} EPHC 2023, Ministerio
Público 2024\\
\textbf{Nota:} Valores marcados como \emph{(dept.)} referem-se ao
departamento de Alto Paraná; valores \emph{(dist.)} são específicos
deste distrito. Dados marcados \emph{(est.)} são estimativas.

\begin{longtable}[]{@{}
  >{\raggedright\arraybackslash}p{(\linewidth - 4\tabcolsep) * \real{0.4231}}
  >{\raggedright\arraybackslash}p{(\linewidth - 4\tabcolsep) * \real{0.2692}}
  >{\raggedright\arraybackslash}p{(\linewidth - 4\tabcolsep) * \real{0.3077}}@{}}
\toprule\noalign{}
\begin{minipage}[b]{\linewidth}\raggedright
Indicador
\end{minipage} & \begin{minipage}[b]{\linewidth}\raggedright
Valor
\end{minipage} & \begin{minipage}[b]{\linewidth}\raggedright
Âmbito
\end{minipage} \\
\midrule\noalign{}
\endhead
\bottomrule\noalign{}
\endlastfoot
IDH (2020) & 0,752 & dept. \\
Ranking IDH nacional & 3/\allowbreak 18 & dept. \\
Esperança de vida & 75,1 anos & dept. \\
Escolaridade média & 8.4 anos & dist. \\
RNB per capita (USD PPA) & 13.500 & dept. \\
Pobreza monetária (\%) & 16,3\% & dept. \\
Pobreza extrema (\%) & 3,1\% & dept. \\
Índice de Gini & 0,428 & dept. \\
Acesso a água potável (\%) & 88,7\% & dept. \\
Acesso a saneamento (\%) & 87,2\% & dept. \\
Taxa de homicídios (est., /\allowbreak 100k hab) & \textasciitilde8--12 (estimativa,
acima da média nacional) & dept. \\
Índice de segurança & baixo & dept. \\
População (Censo 2022) & 6.903 hab. & dist. \\
Idade mediana & 27.0 anos & dist. \\
Taxa de fecundidade & N/\allowbreak D & dist. \\
\end{longtable}
\subsubsection{Combustível}

\textbf{Referência:} PETROPAR /\allowbreak  postos locais (2024) \textbf{Tipo de
localidade:} Interior

\begin{longtable}[]{@{}lll@{}}
\toprule\noalign{}
Combustível & USD/\allowbreak litro & Gs/\allowbreak litro (aprox.) \\
\midrule\noalign{}
\endhead
\bottomrule\noalign{}
\endlastfoot
Gasolina 93 oct & 0.97 & 7,178 \\
Gasolina 97 oct (premium) & 1.07 & 7,918 \\
Diesel & 0.90 & 6,660 \\
\end{longtable}

\begin{quote}
Preços podem variar ±5\% conforme posto e sazonalidade. Chaco e interior
remoto apresentam maior variação.
\end{quote}

\subsubsection{Cobertura Celular}

\textbf{Fonte:} CONATEL PY /\allowbreak  operadoras (2024)

\begin{longtable}[]{@{}ll@{}}
\toprule\noalign{}
Parâmetro & Valor \\
\midrule\noalign{}
\endhead
\bottomrule\noalign{}
\endlastfoot
Cobertura 4G população (dept.) & 95\% \\
Cobertura 4G área rural & 88\% \\
Melhor operadora & Tigo/\allowbreak Personal \\
Qualidade rural & boa \\
\end{longtable}

\begin{quote}
Para áreas rurais fora do núcleo urbano, recomenda-se chip Tigo como
principal e Personal como backup.
\end{quote}

\subsubsection{Internet}

\textbf{Fonte:} CONATEL /\allowbreak  Speedtest Ookla (2024)

\begin{longtable}[]{@{}ll@{}}
\toprule\noalign{}
Parâmetro & Valor \\
\midrule\noalign{}
\endhead
\bottomrule\noalign{}
\endlastfoot
Velocidade média download & 85 Mbps \\
Domicílios com internet (dept.) & 87\% \\
Tecnologia predominante & fibra/\allowbreak cabo \\
Opção rural & Starlink disponível (\textasciitilde USD 44/\allowbreak mês) \\
\end{longtable}

\subsubsection{Mercado Imobiliário e Terra
Rural}

\textbf{Fonte:} INDERT /\allowbreak  Clasificados.com.py (2024)

\begin{longtable}[]{@{}ll@{}}
\toprule\noalign{}
Tipo & Referência \\
\midrule\noalign{}
\endhead
\bottomrule\noalign{}
\endlastfoot
Terra agrícola alta prod. (USD/\allowbreak ha) & 14,800 \\
Imóvel urbano (USD/\allowbreak m²) & 1,400 \\
Aluguel 2 quartos (USD/\allowbreak mês) & 600 \\
\end{longtable}

\begin{quote}
Valores de referência departamental. Localidades menores podem ter
preços 20--40\% abaixo da capital departamental. \#\#\# Saúde
\end{quote}

\textbf{Fonte:} MSPBS /\allowbreak  IPS Paraguay (2024-2026), consolidação
departamental e proxy local

\begin{longtable}[]{@{}ll@{}}
\toprule\noalign{}
Serviço & Disponibilidade \\
\midrule\noalign{}
\endhead
\bottomrule\noalign{}
\endlastfoot
USF /\allowbreak  Posto de Saúde & sim \\
Hospital Regional & não \\
IPS (seguro social) & sim \\
Farmácia & sim \\
Distância ao hospital de referência & \textasciitilde52 km (Ciudad del
Este) \\
\end{longtable}

\textbf{Principais estabelecimentos:} USF local; referencia hospitalar
em Ciudad del Este

\textbf{Observação para imigrantes:} Atendimento primario local ou em
raio curto; casos de maior complexidade seguem para Ciudad del Este.
Cobertura privada continua recomendada para especialidades e urgências
de maior complexidade.
\newpage
\secaoDiagnostico{Yguazu}{\mapaConteudo{-25.4666}{-55.0000}}{Yguazu (Alto Parana) apresenta perfil operacional consolidado para
comparacao territorial, com leitura conjunta de infraestrutura, risco
hidroclimatico e capacidade de continuidade de servicos essenciais.

\subsubsection{Por que sim}

\begin{itemize}
\tightlist
\item
  Base institucional de referencia disponivel e rastreavel.
\item
  Estrutura comparativa (A-E/\allowbreak GSS) consistente para decisao relativa.
\item
  Plano de mitigacao acionavel ja definido no dossie.
\end{itemize}

\subsubsection{Por que não}

\begin{itemize}
\tightlist
\item
  Serie oficial distrital granular ainda pode apresentar lacunas
  temporais.
\item
  Sensibilidade do score exige monitoramento continuo de risco e
  conectividade.
\item
  Decisao final depende de validacao de campo e janela temporal recente.
\end{itemize}

\subsubsection{Pre-condicoes Minimas de
Decisao}

\begin{itemize}
\tightlist
\item
  Atualizar indicadores criticos em ciclo trimestral.
\item
  Confirmar trafegabilidade e contingencia hidrica/\allowbreak energetica local.
\item
  Validar checklist de resiliencia domiciliar/\allowbreak logistica antes de decisao
  final.
\end{itemize}

\subsubsection{Recomendação
Técnica}

Uso recomendado como candidato em analise multicriterio, condicionado ao
cumprimento das pre-condicoes acima. Referencia atual de score: GSS 5.4
em 2026-03-05; risco predominante: alto.

\subsubsection{}

\begin{itemize}
\tightlist
\item
  \begin{enumerate}
  \def\labelenumi{\arabic{enumi})}
  \tightlist
  \item
    Delimitacao territorial validada por georreferencia institucional
    (Fonte: acesso: 2026-03-05).
  \end{enumerate}
\item
  \begin{enumerate}
  \def\labelenumi{\arabic{enumi})}
  \setcounter{enumi}{1}
  \tightlist
  \item
    População municipal/\allowbreak distrital referenciada por base censitaria
    nacional (Fonte: acesso: 2026-03-05).
  \end{enumerate}
\item
  \begin{enumerate}
  \def\labelenumi{\arabic{enumi})}
  \setcounter{enumi}{2}
  \tightlist
  \item
    Comparabilidade intra-departamental apoiada em indicadores
    distritais oficiais (Fonte: acesso: 2026-03-05).
  \end{enumerate}
\item
  \begin{enumerate}
  \def\labelenumi{\arabic{enumi})}
  \setcounter{enumi}{3}
  \tightlist
  \item
    Estado macro de conectividade viaria ancorado em informacoes de
    rodovias (Fonte: acesso: 2026-03-05).
  \end{enumerate}
\item
  \begin{enumerate}
  \def\labelenumi{\arabic{enumi})}
  \setcounter{enumi}{4}
  \tightlist
  \item
    Exposicao hidrometeorologica monitorada por avisos oficiais (Fonte:
    acesso: 2026-03-05).
  \end{enumerate}
\item
  \begin{enumerate}
  \def\labelenumi{\arabic{enumi})}
  \setcounter{enumi}{5}
  \tightlist
  \item
    Historico de contexto meteorologico apoiado em publicacoes tecnicas
    (Fonte: acesso: 2026-03-05).
  \end{enumerate}
\item
  \begin{enumerate}
  \def\labelenumi{\arabic{enumi})}
  \setcounter{enumi}{6}
  \tightlist
  \item
    Capacidade institucional de resposta a eventos apoiada em acoes da
    SEN\footnote{Secretaría de Emergencia Nacional (Secretaria de Emergência Nacional), órgão governamental de resposta a desastres.} (Fonte: acesso: 2026-03-05).
  \end{enumerate}
\item
  \begin{enumerate}
  \def\labelenumi{\arabic{enumi})}
  \setcounter{enumi}{7}
  \tightlist
  \item
    Infraestrutura energetica analisada via operador nacional (Fonte:
    acesso: 2026-03-05).
  \end{enumerate}
\item
  \begin{enumerate}
  \def\labelenumi{\arabic{enumi})}
  \setcounter{enumi}{8}
  \tightlist
  \item
    Referencias de obras e logistica nacional verificadas no MOPC
    (Fonte: acesso: 2026-03-05).
  \end{enumerate}
\item
  \begin{enumerate}
  \def\labelenumi{\arabic{enumi})}
  \setcounter{enumi}{9}
  \tightlist
  \item
    Contexto sociopolitico institucional referenciado por autoridade
    eleitoral (Fonte: acesso: 2026-03-05).
  \end{enumerate}
\item
  \begin{enumerate}
  \def\labelenumi{\arabic{enumi})}
  \setcounter{enumi}{10}
  \tightlist
  \item
    Camada cartografica de apoio eleitoral/\allowbreak territorial consultada
    (Fonte: acesso: 2026-03-05).
  \end{enumerate}
\item
  \begin{enumerate}
  \def\labelenumi{\arabic{enumi})}
  \setcounter{enumi}{11}
  \tightlist
  \item
    Dados publicos complementares para verificacao cruzada (Fonte:
    acesso: 2026-03-05).
  \end{enumerate}
\end{itemize}}
\begin{table}[ht!]
\small \begin{tabular}{p{3.0cm}p{0.8cm}p{6.0cm}} \toprule
\textbf{Critério} & \textbf{Nota} & \textbf{Justificativa} \\ \midrule
A - Ameaças Estratégicas & 6.2 & - \\
B - Risco Social & 5.6 & - \\
C - Riscos Naturais & 3.8 & - \\
D - Autossuficiência & 5.0 & - \\
E - Institucional & 5.8 & - \\
\midrule \textbf{GSS FINAL} & \textbf{5.4} & \textbf{Sensivel (requer mitigacao robusta).} \\ \bottomrule \end{tabular}
\end{table}
\subsubsection{Vulnerabilidades e
Mitigação}\label{vuln-Yguazu-vulnerabilidades-e-mitigauxe7uxe3o}

\begin{itemize}
\tightlist
\item
  Alta exposicao hidrometeorologica (\textgreater=65/\allowbreak 100).
\item
  Conectividade viaria intermediaria (50-69/\allowbreak 100).
\item
  Estoque de contingencia para 14 dias (agua/\allowbreak alimentos/\allowbreak combustivel),
  priorizado quando risco hidro=78/\allowbreak 100.
\item
  Duplicar rotas/\allowbreak logistica quando conectividade viaria=52/\allowbreak 100 for
  \textless70; manter rota secundaria mapeada e testada trimestralmente.
\item
  Plano de energia resiliente com autonomia minima de 72h
  (geracao/\allowbreak backup), revisao semestral.
\end{itemize}
\subsubsection{Dossiê de Campo}\label{dossie-Yguazu-dossiuxea-de-campo}
\begin{description}[style=nextline,leftmargin=0.5cm,font=\bfseries]
\item[Coordenadas]  25°28'S 55°00'W (geocodificado via Nominatim).
\end{description}
\paragraph{Solo (SoilGrids 2.0, média ponderada 0--30
cm)}

\begin{longtable}[]{@{}ll@{}}
\toprule\noalign{}
Parâmetro & Valor \\
\midrule\noalign{}
\endhead
\bottomrule\noalign{}
\endlastfoot
pH (H$_2$O) & 5.2 \\
Carbono orgânico (SOC) & 19.41 g/\allowbreak kg \\
Argila & 38.66 \% \\
Areia & 27.98 \% \\
Silte (calc.) & 33.4 \% \\
Densidade aparente & 1.27 g/\allowbreak cm³ \\
\textbf{Aptidão agrícola} & \textbf{Média} \\
\end{longtable}

Fonte: ISRIC SoilGrids 2.0 via WCS (acesso em 2026-03-21). Coords:
-25.4667°, -55.0°. Média ponderada camadas 0-5, 5-15, 15-30 cm.
\subsubsection{Indicadores Sociais}

\textbf{Fontes:} PNUD 2020, INE\footnote{Instituto Nacional de Estadística (Instituto Nacional de Estatística), órgão oficial de dados demográficos e censitários.} Censo 2022, INE\footnote{Instituto Nacional de Estadística (Instituto Nacional de Estatística), órgão oficial de dados demográficos e censitários.} EPHC 2023, Ministerio
Público 2024\\
\textbf{Nota:} Valores marcados como \emph{(dept.)} referem-se ao
departamento de Alto Paraná; valores \emph{(dist.)} são específicos
deste distrito. Dados marcados \emph{(est.)} são estimativas.

\begin{longtable}[]{@{}
  >{\raggedright\arraybackslash}p{(\linewidth - 4\tabcolsep) * \real{0.4231}}
  >{\raggedright\arraybackslash}p{(\linewidth - 4\tabcolsep) * \real{0.2692}}
  >{\raggedright\arraybackslash}p{(\linewidth - 4\tabcolsep) * \real{0.3077}}@{}}
\toprule\noalign{}
\begin{minipage}[b]{\linewidth}\raggedright
Indicador
\end{minipage} & \begin{minipage}[b]{\linewidth}\raggedright
Valor
\end{minipage} & \begin{minipage}[b]{\linewidth}\raggedright
Âmbito
\end{minipage} \\
\midrule\noalign{}
\endhead
\bottomrule\noalign{}
\endlastfoot
IDH (2020) & 0,752 & dept. \\
Ranking IDH nacional & 3/\allowbreak 18 & dept. \\
Esperança de vida & 75,1 anos & dept. \\
Escolaridade média & 8.7 anos & dist. \\
RNB per capita (USD PPA) & 13.500 & dept. \\
Pobreza monetária (\%) & 16,3\% & dept. \\
Pobreza extrema (\%) & 3,1\% & dept. \\
Índice de Gini & 0,428 & dept. \\
Acesso a água potável (\%) & 88,7\% & dept. \\
Acesso a saneamento (\%) & 87,2\% & dept. \\
Taxa de homicídios (est., /\allowbreak 100k hab) & \textasciitilde8--12 (estimativa,
acima da média nacional) & dept. \\
Índice de segurança & baixo & dept. \\
População (Censo 2022) & 8.131 hab. & dist. \\
Idade mediana & 29.0 anos & dist. \\
Taxa de fecundidade & N/\allowbreak D & dist. \\
\end{longtable}
\subsubsection{Combustível}

\textbf{Referência:} PETROPAR /\allowbreak  postos locais (2024) \textbf{Tipo de
localidade:} Interior

\begin{longtable}[]{@{}lll@{}}
\toprule\noalign{}
Combustível & USD/\allowbreak litro & Gs/\allowbreak litro (aprox.) \\
\midrule\noalign{}
\endhead
\bottomrule\noalign{}
\endlastfoot
Gasolina 93 oct & 0.97 & 7,178 \\
Gasolina 97 oct (premium) & 1.07 & 7,918 \\
Diesel & 0.90 & 6,660 \\
\end{longtable}

\begin{quote}
Preços podem variar ±5\% conforme posto e sazonalidade. Chaco e interior
remoto apresentam maior variação.
\end{quote}

\subsubsection{Cobertura Celular}

\textbf{Fonte:} CONATEL PY /\allowbreak  operadoras (2024)

\begin{longtable}[]{@{}ll@{}}
\toprule\noalign{}
Parâmetro & Valor \\
\midrule\noalign{}
\endhead
\bottomrule\noalign{}
\endlastfoot
Cobertura 4G população (dept.) & 95\% \\
Cobertura 4G área rural & 88\% \\
Melhor operadora & Tigo/\allowbreak Personal \\
Qualidade rural & boa \\
\end{longtable}

\begin{quote}
Para áreas rurais fora do núcleo urbano, recomenda-se chip Tigo como
principal e Personal como backup.
\end{quote}

\subsubsection{Internet}

\textbf{Fonte:} CONATEL /\allowbreak  Speedtest Ookla (2024)

\begin{longtable}[]{@{}ll@{}}
\toprule\noalign{}
Parâmetro & Valor \\
\midrule\noalign{}
\endhead
\bottomrule\noalign{}
\endlastfoot
Velocidade média download & 85 Mbps \\
Domicílios com internet (dept.) & 87\% \\
Tecnologia predominante & fibra/\allowbreak cabo \\
Opção rural & Starlink disponível (\textasciitilde USD 44/\allowbreak mês) \\
\end{longtable}

\subsubsection{Mercado Imobiliário e Terra
Rural}

\textbf{Fonte:} INDERT /\allowbreak  Clasificados.com.py (2024)

\begin{longtable}[]{@{}ll@{}}
\toprule\noalign{}
Tipo & Referência \\
\midrule\noalign{}
\endhead
\bottomrule\noalign{}
\endlastfoot
Terra agrícola alta prod. (USD/\allowbreak ha) & 14,800 \\
Imóvel urbano (USD/\allowbreak m²) & 1,400 \\
Aluguel 2 quartos (USD/\allowbreak mês) & 600 \\
\end{longtable}

\begin{quote}
Valores de referência departamental. Localidades menores podem ter
preços 20--40\% abaixo da capital departamental. \#\#\# Saúde
\end{quote}

\textbf{Fonte:} MSPBS /\allowbreak  IPS Paraguay (2024-2026), consolidação
departamental e proxy local

\begin{longtable}[]{@{}ll@{}}
\toprule\noalign{}
Serviço & Disponibilidade \\
\midrule\noalign{}
\endhead
\bottomrule\noalign{}
\endlastfoot
USF /\allowbreak  Posto de Saúde & sim \\
Hospital Regional & não \\
IPS (seguro social) & sim \\
Farmácia & sim \\
Distância ao hospital de referência & \textasciitilde49 km (Ciudad del
Este) \\
\end{longtable}

\textbf{Principais estabelecimentos:} USF local; referencia hospitalar
em Ciudad del Este

\textbf{Observação para imigrantes:} Atendimento primario local ou em
raio curto; casos de maior complexidade seguem para Ciudad del Este.
Cobertura privada continua recomendada para especialidades e urgências
de maior complexidade.
